% Expression de l’information génétique — SVTEEHB 3ème — Cours signé OnBuch+
% Programme officiel MINESEC. Compiler avec : tectonic N03-expression-information-genetique.tex
\newcommand{\DOCMATIERE}{SVTEEHB}
\newcommand{\DOCNIVEAU}{3ème}
\newcommand{\DOCMODULE}{SVTEEHB – classe de 3ème}
\newcommand{\DOCLECON}{N03}
\newcommand{\DOCTITRE}{Expression de l’information génétique}
% OnBuch+ — préambule « cours signés » ENRICHI (compilé avec tectonic / XeLaTeX)
% Système visuel « Pop » d'OnBuch+ : fond crème (jamais blanc), bordures encre
% épaisses, ombres « dures » décalées, titres Archivo Black en capitales.
% Version MATHÉMATIQUES : graphes (pgfplots/tikz), boîtes pédagogiques
% (définition, propriété, méthode, attention, le savais-tu, exercice résolu,
% exercice, TP, bilan), page de garde et signature OnBuch+ ancrée en bas.
%
% Macros attendues dans main.tex AVANT \input{preamble} :
%   \DOCMATIERE  \DOCNIVEAU  \DOCMODULE  \DOCLECON (numéro)  \DOCTITRE
\documentclass[11pt]{article}

\usepackage{fontspec}
\usepackage[french]{babel}
\usepackage[a4paper,margin=2cm,top=3.4cm,bottom=2.8cm,headheight=1.4cm,footskip=1.3cm]{geometry}
\usepackage{amsmath,amssymb,amsfonts}
\usepackage{mathtools} % \xmapsto, \coloneqq... souvent utilisés par les modèles
\usepackage{cancel} % simplifications barrées (bilans de forces)
% \abs{z} (module, valeur absolue) et \norm{u} (norme), utilisés par les modèles
\providecommand{\abs}{}\let\abs\relax\DeclarePairedDelimiter{\abs}{\lvert}{\rvert}
\providecommand{\norm}{}\let\norm\relax\DeclarePairedDelimiter{\norm}{\lVert}{\rVert}
\usepackage[table]{xcolor} % [table] : fournit aussi \rowcolors (alternance de lignes)
\usepackage{pagecolor}
\usepackage{tikz}
\usetikzlibrary{arrows.meta,positioning,calc,shapes.geometric,shapes.misc,shapes.arrows,decorations.pathmorphing,decorations.pathreplacing,decorations.markings,patterns,shadows,mindmap,trees,fit,backgrounds,matrix,intersections,angles,quotes}
\usepackage{pgfplots}
\usepackage[version=4]{mhchem} % équations nucléaires \ce{^{235}_{92}U}
\usepackage[european,siunitx]{circuitikz} % schémas électriques
\pgfplotsset{compat=1.18}
\usepgfplotslibrary{fillbetween,groupplots} % aires entre courbes (calcul intégral)
\usepackage{enumitem}
\usepackage{fancyhdr}
% [most] charge toutes les bibliothèques tcolorbox (listings, minted, raster,
% documentation...) et gonfle inutilement la mémoire TeX sur un document long
% (~300 boîtes) : on ne charge que ce qui est réellement utilisé.
\usepackage[skins,breakable]{tcolorbox}
\usepackage{graphicx}
\usepackage{colortbl}
\usepackage{booktabs}
\usepackage{tabularx}
\usepackage{multirow}
\usepackage{diagbox} % case d'en-tête diagonale des tableaux à double entrée
% Colonnes extensibles centrée / gauche / droite (C, L, R), souvent utilisées par les modèles
\newcolumntype{C}{>{\centering\arraybackslash}X}
\newcolumntype{L}{>{\raggedright\arraybackslash}X}
\newcolumntype{R}{>{\raggedleft\arraybackslash}X}
\usepackage{array}
\usepackage{multicol}
\usepackage{siunitx}
% parse-numbers=false : accepte aussi \SI{2,5 \times 10^{-3}}{...}
\sisetup{locale=FR,output-decimal-marker={,},group-minimum-digits=4,parse-numbers=false}
\DeclareSIUnit\ohm{\ensuremath{\symup{\Omega}}} % U+2126 absent de la police
\DeclareSIUnit\division{div} % réglages d'oscilloscope (ms/div, V/div)
\DeclareSIUnit\tr{tr} % tours (vitesse de rotation, ex. \SI{1500}{\tr\per\minute})
\DeclareSIUnit\molar{M} % concentration molaire (ex. \SI{3}{\molar}, \SI{1}{\micro\molar})
\DeclareSIUnit\an{an} % année (le modèle confond parfois avec \year, un compteur TeX) — ex. \SI{5}{\kilo\meter\per\an}
\DeclareSIUnit\FCFA{FCFA} % franc CFA (ex. \SI{500}{\FCFA\per\kilo\gram})
\DeclareSIUnit\degreeBrix{\textdegree Brix} % teneur en sucre d'un sirop/jus (réfractométrie)
\DeclareSIUnit\month{mois} % le modèle utilise parfois \month, un compteur TeX, comme unité de durée
\DeclareSIUnit\microliter{\micro\liter} % le modèle écrit parfois \microliter en un seul token
\DeclareSIUnit\million{millions} % dénombrement (ex. \SI{15}{\million\per\mL} spermatozoïdes)
\usepackage{qrcode}
% \degree : commande fréquemment utilisée par les modèles pour un angle ou une
% température en dehors de \SI{}{} (non fournie par les paquets chargés).
\providecommand{\degree}{^{\circ}}
% \qed (fin de démonstration), utilisé par les modèles sans amsthm
\providecommand{\qed}{\hfill$\square$}
% \barre : double barre des valeurs interdites dans un tableau de variations (tkz-tab)
\providecommand{\barre}{\ensuremath{\|}}
% environnement proof (amsthm non chargé) utilisé par les modèles
\ifdefined\proof\else\newenvironment{proof}[1][Démonstration]{\par\noindent\textit{#1.}\ }{\qed\par}\fi
\usepackage{lastpage}
\usepackage{hyperref}
\hypersetup{hidelinks}

% ============================================================
% Palette « Pop » OnBuch+ — mêmes hex que la classe P (plus_ui.dart)
% ============================================================
\definecolor{poporange}{HTML}{F59321}
\definecolor{popdark}{HTML}{E07A0C}
\definecolor{popcream}{HTML}{FFF6E8}
\definecolor{popink}{HTML}{1C1714}
\definecolor{poppurple}{HTML}{7A5AE0}
\definecolor{popgreen}{HTML}{2E9E6A}
\definecolor{poppink}{HTML}{E0457B}
\definecolor{popblue}{HTML}{2F7DE1}
\definecolor{popgold}{HTML}{C98A12}
\definecolor{popmuted}{HTML}{8A7F76}
\definecolor{popline}{HTML}{EADFD2}
\definecolor{popgray}{HTML}{8A7F76}
\colorlet{popgrey}{popgray}
% Alias tolérants (le modèle nomme parfois les couleurs différemment)
\colorlet{popwhite}{white}
\colorlet{popblack}{popink}
\colorlet{popred}{poppink}
\colorlet{popyellow}{popgold}
% Teintes claires (fonds de tableaux/encadrés) — le modèle nomme parfois
% les couleurs avec un suffixe L (ex. popblueL) sans les avoir définies.
\colorlet{poporangeL}{poporange!20}
\colorlet{popdarkL}{popdark!20}
\colorlet{poppurpleL}{poppurple!20}
\colorlet{popgreenL}{popgreen!20}
\colorlet{poppinkL}{poppink!20}
\colorlet{popblueL}{popblue!20}
\colorlet{popgoldL}{popgold!20}
\colorlet{popredL}{popred!20}
\colorlet{popgrayL}{popgray!20}
% Teintes claires pour fonds de tableaux / zones
\colorlet{popblueL}{popblue!10!white}
\colorlet{poporangeL}{poporange!14!white}
\colorlet{popgreenL}{popgreen!12!white}
\colorlet{poppurpleL}{poppurple!12!white}
\colorlet{poppinkL}{poppink!12!white}
\colorlet{popgoldL}{popgold!14!white}

\pagecolor{popcream}

% ============================================================
% Typographie
% ============================================================
\defaultfontfeatures{Ligatures=TeX}
\setmainfont{PlusJakartaSans-Medium.ttf}[Path=../../../fonts/,BoldFont=PlusJakartaSans-Bold.ttf]
% Maths en sans-serif (Fira Math) avec chiffres et lettres droites de Plus
% Jakarta, pour que les formules se fondent dans le texte.
\usepackage[math-style=ISO,bold-style=ISO]{unicode-math}
\setmathfont{FiraMath-Regular.otf}[Path=../../../fonts/]
\setmathfont{PlusJakartaSans-Medium.ttf}[Path=../../../fonts/,range={up/num,up/{Latin,latin},\mathup}]
\usepackage{icomma} % virgule décimale sans espace en mode math
% Caractères mathématiques Unicode absents des polices de texte (Plus Jakarta,
% Archivo Black) : rendus par leur équivalent mathématique, en texte comme en
% formule — sinon ils s'affichent en carré vide (ex. « Arithmétique dans ℤ »).
\usepackage{newunicodechar}
\newunicodechar{ℝ}{\ensuremath{\mathbb{R}}}
\newunicodechar{ℤ}{\ensuremath{\mathbb{Z}}}
\newunicodechar{ℂ}{\ensuremath{\mathbb{C}}}
\newunicodechar{ℕ}{\ensuremath{\mathbb{N}}}
\newunicodechar{ℚ}{\ensuremath{\mathbb{Q}}}
\newunicodechar{𝔻}{\ensuremath{\mathbb{D}}}
\newunicodechar{α}{\ensuremath{\alpha}}
\newunicodechar{β}{\ensuremath{\beta}}
\newunicodechar{γ}{\ensuremath{\gamma}}
\newunicodechar{δ}{\ensuremath{\delta}}
\newunicodechar{ε}{\ensuremath{\varepsilon}}
\newunicodechar{θ}{\ensuremath{\theta}}
\newunicodechar{λ}{\ensuremath{\lambda}}
\newunicodechar{μ}{\ensuremath{\mu}}
\newunicodechar{σ}{\ensuremath{\sigma}}
\newunicodechar{τ}{\ensuremath{\tau}}
\newunicodechar{φ}{\ensuremath{\varphi}}
\newunicodechar{ψ}{\ensuremath{\psi}}
\newunicodechar{ω}{\ensuremath{\omega}}
\newunicodechar{Σ}{\ensuremath{\Sigma}}
% majuscules produites par \MakeUppercase dans les titres (α-aminés -> Α)
\newunicodechar{Α}{\ensuremath{\alpha}}
\newunicodechar{Β}{\ensuremath{\beta}}
\newunicodechar{Γ}{\ensuremath{\Gamma}}
\newunicodechar{Δ}{\ensuremath{\Delta}}
\newunicodechar{Ε}{\ensuremath{\varepsilon}}
\newunicodechar{Θ}{\ensuremath{\Theta}}
\newunicodechar{Λ}{\ensuremath{\Lambda}}
\newunicodechar{Μ}{\ensuremath{\mu}}
\newunicodechar{Τ}{\ensuremath{\tau}}
\newunicodechar{Φ}{\ensuremath{\Phi}}
\newunicodechar{Ψ}{\ensuremath{\Psi}}
\newunicodechar{Ω}{\ensuremath{\Omega}}
\newunicodechar{↦}{\ensuremath{\mapsto}}
\newunicodechar{∈}{\ensuremath{\in}}
\newunicodechar{∉}{\ensuremath{\notin}}
\newunicodechar{∧}{\ensuremath{\wedge}}
\newunicodechar{⇔}{\ensuremath{\Leftrightarrow}}
\newunicodechar{⟺}{\ensuremath{\Longleftrightarrow}}
\newunicodechar{⇒}{\ensuremath{\Rightarrow}}
\newunicodechar{⟹}{\ensuremath{\Longrightarrow}}
\newunicodechar{∘}{\ensuremath{\circ}}
\newunicodechar{∪}{\ensuremath{\cup}}
\newunicodechar{∩}{\ensuremath{\cap}}
\newunicodechar{≡}{\ensuremath{\equiv}}
\newunicodechar{⊂}{\ensuremath{\subset}}
\newunicodechar{⊥}{\ensuremath{\perp}}
\newunicodechar{∃}{\ensuremath{\exists}}
\newunicodechar{∀}{\ensuremath{\forall}}
\newunicodechar{∛}{\ensuremath{\sqrt[3]{\,}}}
\newunicodechar{∜}{\ensuremath{\sqrt[4]{\,}}}
\newunicodechar{⃗}{\ensuremath{{}^{\to}}}
\newunicodechar{ⁿ}{\ifmmode^{n}\else\textsuperscript{\ensuremath{n}}\fi}
\newunicodechar{ˣ}{\ifmmode^{x}\else\textsuperscript{\ensuremath{x}}\fi}
\newunicodechar{ᵇ}{\ifmmode^{b}\else\textsuperscript{\ensuremath{b}}\fi}
\newunicodechar{ʳ}{\ifmmode^{r}\else\textsuperscript{\ensuremath{r}}\fi}
\newunicodechar{ᵘ}{\ifmmode^{u}\else\textsuperscript{\ensuremath{u}}\fi}
\newunicodechar{ʸ}{\ifmmode^{y}\else\textsuperscript{\ensuremath{y}}\fi}
\newunicodechar{ᵃ}{\ifmmode^{a}\else\textsuperscript{\ensuremath{a}}\fi}
\newunicodechar{ᵉ}{\ifmmode^{e}\else\textsuperscript{\ensuremath{e}}\fi}
\newunicodechar{ᵏ}{\ifmmode^{k}\else\textsuperscript{\ensuremath{k}}\fi}
\newunicodechar{ᵖ}{\ifmmode^{p}\else\textsuperscript{\ensuremath{p}}\fi}
\newunicodechar{ᴺ}{\ifmmode^{N}\else\textsuperscript{\ensuremath{N}}\fi}
\newunicodechar{ᶜ}{\ifmmode^{c}\else\textsuperscript{\ensuremath{c}}\fi}
\newunicodechar{ʰ}{\ifmmode^{h}\else\textsuperscript{\ensuremath{h}}\fi}
\newunicodechar{ᵠ}{\ifmmode^{q}\else\textsuperscript{\ensuremath{q}}\fi}
\newunicodechar{ₙ}{\ifmmode_{n}\else\textsubscript{\ensuremath{n}}\fi}
\newunicodechar{ₐ}{\ifmmode_{a}\else\textsubscript{\ensuremath{a}}\fi}
\newunicodechar{ᵢ}{\ifmmode_{i}\else\textsubscript{\ensuremath{i}}\fi}
\newunicodechar{ᵣ}{\ifmmode_{r}\else\textsubscript{\ensuremath{r}}\fi}
\newunicodechar{ₖ}{\ifmmode_{k}\else\textsubscript{\ensuremath{k}}\fi}
\newunicodechar{ᵦ}{\ifmmode_{\beta}\else\textsubscript{\ensuremath{\beta}}\fi}
\newunicodechar{ₓ}{\ifmmode_{x}\else\textsubscript{\ensuremath{x}}\fi}
\newfontfamily\popdisplay{ArchivoBlack-Regular.ttf}[Path=../../../fonts/]
\newfontfamily\poplogo{SpaceGrotesk-Bold.ttf}[Path=../../../fonts/]
\newfontfamily\popsemi{PlusJakartaSans-SemiBold.ttf}[Path=../../../fonts/]
\newfontfamily\popxbold{PlusJakartaSans-ExtraBold.ttf}[Path=../../../fonts/]
\newfontfamily\popxboldls{PlusJakartaSans-ExtraBold.ttf}[Path=../../../fonts/,LetterSpace=3.0]

\setlist[enumerate]{leftmargin=1.7em,itemsep=5pt,topsep=4pt}
\setlist[itemize]{leftmargin=1.7em,itemsep=4pt,topsep=4pt,label={\color{poporange}\textbullet}}
\setlength{\parindent}{0pt}
\setlength{\parskip}{5pt}
\renewcommand{\arraystretch}{1.25}

% pgfplots : style Pop par défaut
\pgfplotsset{
  every axis/.append style={axis line style={popink,line width=1pt},tick style={popink},
    grid style={popline,line width=0.5pt},label style={font=\small},tick label style={font=\footnotesize}},
  popcurve/.style={poporange,line width=1.8pt,no markers,smooth},
}
% Même style disponible dans un \draw TikZ simple (hors pgfplots)
\tikzset{popcurve/.style={poporange,line width=1.8pt}}
\tikzset{popgrid/.style={popline,very thin}}
\colorlet{popcurve}{poporange} % le nom du style est parfois utilisé comme couleur

% ============================================================
% Pill / squiggle
% ============================================================
\newcommand{\poppill}[3]{%
  \tikz[baseline=-0.55ex]\node[draw=popink,line width=1.2pt,rounded corners=2.6pt,%
    fill=#2,inner xsep=6.5pt,inner ysep=3pt]{%
    \popxboldls\fontsize{7}{7}\selectfont\color{#3}\MakeUppercase{#1}};%
}
\newcommand{\popsquiggle}[1]{%
  \tikz[baseline=-0.4ex]{
    \draw[#1,line width=2.6pt,line cap=round]
      plot [smooth] coordinates {(0,0.09) (0.34,0.01) (0.68,0.21) (1.02,0.01) (1.36,0.10)};
  }%
}
% Mot-clé surligné (marqueur orange sous le texte)
\newcommand{\cle}[1]{{\popxbold\color{popdark}#1}}

% ============================================================
% En-tête / pied
% ============================================================
\pagestyle{fancy}
\fancyhf{}
\renewcommand{\headrulewidth}{0pt}
\renewcommand{\footrulewidth}{0pt}
\lhead{\raisebox{-0.15ex}{\poplogo\fontsize{13}{13}\selectfont\color{popink}OnBuch\color{poporange}+}}
\rhead{\poppill{\DOCMATIERE\ \textbullet\ \DOCNIVEAU\ \textbullet\ Leçon \DOCLECON}{white}{popink}}
\lfoot{\popxbold\fontsize{7.6}{7.6}\selectfont\color{popmuted}\MakeUppercase{Cours signé OnBuch+}\ \textbullet\ {\popsemi\DOCTITRE}}
\rfoot{\tikz[baseline=-0.6ex]\node[rounded corners=8.5pt,fill=popink,minimum height=17pt,minimum width=17pt,inner xsep=5pt,inner sep=0]{\popxbold\fontsize{8}{8}\selectfont\color{popcream}\thepage\,/\,\pageref*{LastPage}};}

% ============================================================
% Titres
% ============================================================
\newcommand{\chaptitre}[2]{%
  \begin{tcolorbox}[enhanced,colback=poporange,colframe=popink,boxrule=2.6pt,arc=11pt,
      drop shadow=popink,left=15pt,right=15pt,top=12pt,bottom=11pt,before skip=3pt,after skip=18pt]
    \poppill{#2}{popink}{popcream}\par\vspace{9pt}
    {\raggedright\hyphenpenalty=10000\exhyphenpenalty=10000\popdisplay\fontsize{22}{24}\selectfont\color{popcream}\MakeUppercase{#1}\par}
  \end{tcolorbox}
}
% Section numérotée : pastille orange avec le numéro + titre Archivo
\newcounter{popsec}
\newcommand{\coursec}[1]{%
  \stepcounter{popsec}\setcounter{popsub}{0}%
  \par\vspace{16pt}\noindent
  \begin{minipage}{\linewidth}
  \tikz[baseline=-0.7ex]\node[draw=popink,line width=1.6pt,fill=poporange,rounded corners=4pt,minimum size=22pt,inner sep=0]{\popdisplay\fontsize{12}{12}\selectfont\color{popcream}\thepopsec};%
  \hspace{8pt}{\popdisplay\fontsize{15}{17}\selectfont\color{popink}\MakeUppercase{#1}}\par
  \vspace{4pt}\noindent{\color{poporange}\rule{36pt}{3.6pt}}
  \end{minipage}\par\nopagebreak\vspace{8pt}
}
\newcounter{popsub}
\newcommand{\courssub}[1]{%
  \stepcounter{popsub}%
  \par\vspace{9pt}\noindent{\popxbold\fontsize{11.5}{13}\selectfont\color{popink}{\color{poporange}\thepopsec.\thepopsub}\hspace{6pt}#1}\par\nopagebreak\vspace{4pt}
}

% ============================================================
% Boîtes pédagogiques — même recette : fond blanc, bordure épaisse colorée,
% ombre dure encre, badge plein. Titre optionnel [#1].
% ============================================================
\tcbset{popbase/.style={breakable,enhanced,colback=white,boxrule=2.3pt,arc=9pt,
  shadow={3pt}{-3pt}{0pt}{popink},left=14pt,right=14pt,top=11pt,bottom=11pt,before skip=11pt,after skip=15pt,
  fontupper=\popsemi\fontsize{10.6}{14.5}\selectfont\color{popink},
  fontlower=\popsemi\fontsize{10.6}{14.5}\selectfont\color{popink}}}
% Coche dessinée (le glyphe ✓ n'existe pas dans les polices du document)
\newcommand{\popcheck}[1]{\tikz[baseline=-0.5ex]\draw[#1,line width=1.6pt,line cap=round,line join=round](-0.13,0.0)--(-0.04,-0.09)--(0.13,0.1);}
\providecommand{\coche}{\popcheck{popgreen}}
\providecommand{\resultat}[1]{\par\smallskip\noindent\fcolorbox{popgreen}{popgreenL}{\parbox{\dimexpr\linewidth-2\fboxsep-2\fboxrule\relax}{\textbf{Résultat.} #1}}\par\smallskip}
\newcommand{\popboxhead}[3]{\poppill{#1}{#2}{white}\ifx\relax#3\relax\else\hspace{6pt}{\popxbold\fontsize{10.6}{12}\selectfont\color{popink}#3}\fi\par\vspace{7pt}}

\newtcolorbox{aretenir}[1][]{popbase,colframe=popgreen,before upper={\popboxhead{À retenir}{popgreen}{#1}}}
\newtcolorbox{exemplebox}[1][]{popbase,colframe=poporange,before upper={\popboxhead{Exemple}{poporange}{#1}}}
\newtcolorbox{definition}[1][]{popbase,colframe=popblue,before upper={\popboxhead{Définition}{popblue}{#1}}}
\newtcolorbox{propriete}[1][]{popbase,colframe=poppurple,before upper={\popboxhead{Propriété}{poppurple}{#1}}}
\newtcolorbox{methode}[1][]{popbase,colframe=popgold,before upper={\popboxhead{Méthode}{popgold}{#1}}}
\newtcolorbox{attention}[1][]{popbase,colframe=poppink,before upper={\popboxhead{Attention}{poppink}{#1}}}
\newtcolorbox{experience}[1][]{popbase,colframe=popblue,colback=popblueL,before upper={\popboxhead{Expérience}{popblue}{#1}}}
% « Le savais-tu ? » : carte violette pleine (contexte, culture, Cameroun)
\newtcolorbox{savaistu}[1][]{popbase,colback=poppurple,colframe=popink,
  fontupper=\popsemi\fontsize{10.4}{14.2}\selectfont\color{white},
  before upper={\poppill{Le savais-tu\,?}{white}{poppurple}\ifx\relax#1\relax\else\hspace{6pt}{\popxbold\color{white}#1}\fi\par\vspace{7pt}}}
% Exercice résolu : énoncé en haut, \tcblower puis la solution sur fond crème
\newcounter{popexo}\newcounter{popexr}
\newtcolorbox{exoresolu}[1][]{popbase,colframe=poporange,colbacklower=popcream,
  segmentation style={popink,line width=1.2pt,dashed},
  before upper={\stepcounter{popexr}\popboxhead{Exercice résolu \thepopexr}{poporange}{#1}},
  before lower={\poppill{Solution}{popink}{popcream}\par\vspace{6pt}}}
% Exercice (non résolu en place) — la correction est dans la partie Corrigés
\newtcolorbox{exercice}[1][]{popbase,colframe=popink,
  before upper={\stepcounter{popexo}\popboxhead{Exercice \thepopexo}{popink}{#1}}}
% Corrigé
\newtcolorbox{corrige}[1][]{popbase,colframe=popgreen,colback=popgreenL,
  before upper={\popboxhead{Corrigé}{popgreen}{#1}}}
% Objectifs / prérequis / bilan
\newtcolorbox{objectifs}{popbase,colframe=popink,colback=poporangeL,
  before upper={\popboxhead{Objectifs de la leçon}{popink}{}}}
\newtcolorbox{prerequis}{popbase,colframe=popink,colback=popblueL,
  before upper={\popboxhead{Prérequis}{popblue}{}}}
\newtcolorbox{bilan}{popbase,colframe=popink,colback=white,boxrule=2.8pt,
  before upper={\popboxhead{Fiche bilan}{poporange}{L'essentiel à connaître pour l'examen}}}
% Figure encadrée avec légende
\newtcolorbox{popfigure}[1][]{enhanced,breakable=false,colback=white,colframe=popline,boxrule=1.4pt,arc=8pt,
  left=10pt,right=10pt,top=10pt,bottom=8pt,before skip=10pt,after skip=12pt,halign=center,
  lower separated=false,halign lower=center,
  fontlower=\popsemi\fontsize{9}{11.5}\selectfont\color{popmuted},
  #1}
\newcommand{\legende}[1]{\par\vspace{6pt}{\popsemi\fontsize{9}{11.5}\selectfont\color{popmuted}#1}}

% ============================================================
% Signature OnBuch+
% ============================================================
\newcommand{\poplogoinline}[1]{{\poplogo\fontsize{#1}{#1}\selectfont\color{popink}OnBuch\color{poporange}+}}
\newcommand{\signatureonbuch}{%
  \begin{tcolorbox}[enhanced,colback=popink,colframe=popink,boxrule=2.2pt,arc=10pt,
      drop shadow=poporange,left=14pt,right=14pt,top=10pt,bottom=10pt,before skip=0pt,after skip=0pt]
    \tikz[baseline=-0.7ex]\node[circle,fill=popgreen,draw=popcream,line width=1.3pt,minimum size=22pt,inner sep=0]{\popcheck{white}};%
    \hspace{9pt}\begin{minipage}[c]{0.62\linewidth}
      {\popxbold\fontsize{12}{13}\selectfont\color{popcream}Signé\ \textbullet\ {\poplogo OnBuch\color{poporange}+}}\par\vspace{2pt}
      {\raggedright\popsemi\fontsize{8}{10.5}\selectfont\color{popline}\MakeUppercase{Équipe pédagogique OnBuch+}\\\MakeUppercase{Conforme au programme officiel MINESEC}\par}
    \end{minipage}\hfill
    {\poplogo\fontsize{22}{22}\selectfont\color{popcream}OnBuch\color{poporange}+}
  \end{tcolorbox}%
}

% ============================================================
% Page de garde
% #1 titre · #2 matière · #3 niveau · #4 module · #5 n° leçon · #6 durée ·
% #7 description / objectifs (texte ou itemize)
% ============================================================
\newcommand{\pagegarde}[7]{%
  \thispagestyle{empty}
  \newgeometry{a4paper,margin=1.7cm,top=1.6cm,bottom=1.4cm}%
  \begin{tikzpicture}[remember picture,overlay]
    \fill[poporange!18] ([xshift=-3.2cm,yshift=-3.6cm]current page.north east) circle (4.6cm);
    \fill[poppurple!14] ([xshift=2.4cm,yshift=7.6cm]current page.south west) circle (3.8cm);
  \end{tikzpicture}%
  {\poplogo\fontsize{30}{30}\selectfont\color{popink}OnBuch\color{poporange}+}\hfill
  \raisebox{0.6ex}{\poppill{Cours signé \textbullet\ Premium}{popink}{popcream}}\par
  \vspace{1pt}\popsquiggle{poppurple}\par
  \vspace{8pt}
  \begin{tcolorbox}[enhanced,colback=poporange,colframe=popink,boxrule=2.8pt,arc=13pt,
      drop shadow=popink,left=18pt,right=18pt,top=12pt,bottom=13pt,after skip=10pt]
    \poppill{#2 \textbullet\ #3}{popink}{popcream}\hspace{5pt}\poppill{#4}{white}{popink}\par\vspace{8pt}
    {\popdisplay\fontsize{14}{14}\selectfont\color{popink}LEÇON #5}\par\vspace{4pt}
    {\raggedright\hyphenpenalty=10000\exhyphenpenalty=10000\popdisplay\fontsize{25}{27}\selectfont\color{popcream}\MakeUppercase{#1}\par}
    \vspace{8pt}
    {\popxbold\fontsize{9.5}{9.5}\selectfont\color{popink}\MakeUppercase{Durée conseillée : #6}}
  \end{tcolorbox}
  \begin{tcolorbox}[enhanced,colback=white,colframe=popink,boxrule=2.2pt,arc=10pt,
      drop shadow=popink,left=16pt,right=16pt,top=10pt,bottom=10pt,after skip=8pt]
    \poppill{Dans cette leçon}{poppurple}{white}\par\vspace{5pt}
    \popsemi\fontsize{9.3}{12.6}\selectfont\color{popink}#7
  \end{tcolorbox}
  \noindent\begin{minipage}[t]{0.487\linewidth}
    \begin{tcolorbox}[enhanced,colback=popgreen,colframe=popink,boxrule=2.2pt,arc=10pt,
        drop shadow=popink,left=11pt,right=11pt,top=10pt,bottom=10pt,height=3.1cm,valign=center]
      \tcbox[colback=white,colframe=popink,boxrule=1.3pt,arc=5pt,boxsep=0pt,nobeforeafter,
        left=3pt,right=3pt,top=3pt,bottom=3pt,valign=center]{\qrcode[height=1.9cm]{https://play.google.com/store/apps/details?id=com.luvvix.onbuch&hl=fr}}%
      \hspace{8pt}\begin{minipage}[c]{0.5\linewidth}
        {\popdisplay\fontsize{11}{12.5}\selectfont\color{white}\MakeUppercase{Télécharge OnBuch}}\par\vspace{4pt}
        {\raggedright\popsemi\fontsize{8.4}{11}\selectfont\color{white}Gratuit \textbullet\ Google Play\\Tuteur IA, annales, concours, résultats\par}
      \end{minipage}
    \end{tcolorbox}
  \end{minipage}\hfill
  \begin{minipage}[t]{0.487\linewidth}
    \begin{tcolorbox}[enhanced,colback=poppurple,colframe=popink,boxrule=2.2pt,arc=10pt,
        drop shadow=popink,left=11pt,right=11pt,top=10pt,bottom=10pt,height=3.1cm,valign=center]
      \tikz[baseline=-0.7ex]\node[circle,fill=white,draw=popink,line width=1.3pt,minimum size=1.6cm,inner sep=0]{\popdisplay\fontsize{24}{24}\selectfont\color{poppurple}+};%
      \hspace{8pt}\begin{minipage}[c]{0.6\linewidth}
        {\popdisplay\fontsize{11}{12.5}\selectfont\color{white}\MakeUppercase{Passe à OnBuch+}}\par\vspace{4pt}
        {\raggedright\popsemi\fontsize{8.4}{11}\selectfont\color{white}Tous les cours signés, Tuteur IA illimité, fiches PDF — onglet OnBuch+ de l'app\par}
      \end{minipage}
    \end{tcolorbox}
  \end{minipage}
  \vspace{8pt}
  \popcontenu
  \vfill
  \signatureonbuch
  \restoregeometry
  \newpage
}

% Rangée « Ce cours contient » (page de garde)
\newcommand{\poptile}[3]{% #1 couleur #2 gros texte #3 légende
  \begin{tcolorbox}[enhanced,colback=#1,colframe=popink,boxrule=2pt,arc=9pt,shadow={2.5pt}{-2.5pt}{0pt}{popink},
      left=6pt,right=6pt,top=9pt,bottom=9pt,halign=center,valign=center,height=1.75cm,width=\linewidth,nobeforeafter]
    {\popdisplay\fontsize{12.5}{13}\selectfont\color{white}\MakeUppercase{#2}}\par\vspace{3pt}
    {\centering\popsemi\fontsize{7.8}{9.5}\selectfont\color{white}#3\par}
  \end{tcolorbox}}
\newcommand{\popcontenu}{%
  \par\noindent{\popxboldls\fontsize{7.5}{7.5}\selectfont\color{popmuted}CE COURS SIGNÉ CONTIENT}\par\vspace{7pt}
  \noindent\begin{minipage}[t]{0.235\linewidth}\poptile{popblue}{Cours}{complet, illustré, pas à pas}\end{minipage}\hfill
  \begin{minipage}[t]{0.235\linewidth}\poptile{poporange}{Méthodes}{et exercices résolus}\end{minipage}\hfill
  \begin{minipage}[t]{0.235\linewidth}\poptile{poppink}{Exercices}{type examen + corrigés}\end{minipage}\hfill
  \begin{minipage}[t]{0.235\linewidth}\poptile{popgreen}{Fiche bilan}{l'essentiel à retenir}\end{minipage}\par}

% Sommaire
\newtcolorbox{sommaire}{enhanced,colback=white,colframe=popink,boxrule=2.2pt,arc=10pt,
  drop shadow=popink,left=18pt,right=18pt,top=13pt,bottom=13pt,before skip=6pt,after skip=20pt,
  before upper={\poppill{Sommaire}{poppurple}{white}\par\vspace{9pt}\popsemi\fontsize{10.6}{14.5}\selectfont\color{popink}}}

% Signature de fin de cours — poussée tout en bas de la dernière page
\newcommand{\findecours}{%
  \par\vfill\nopagebreak
  \begin{center}\popsquiggle{poporange}\end{center}\vspace{4pt}
  \signatureonbuch
}
\AtBeginDocument{\def\llbracket{\mathopen{[\mkern-3.5mu[}}\def\rrbracket{\mathclose{]\mkern-3.5mu]}}} % intervalles d'entiers (glyphe ⟦ absent de la police)
% Glyphes absents de Fira Math : redéfinis à partir de glyphes disponibles
\AtBeginDocument{%
  \def\setminus{\mathbin{\backslash}}\let\smallsetminus\setminus
  \def\vdots{\mathord{\vbox{\baselineskip3.2pt\lineskiplimit0pt\kern4pt\hbox{.}\hbox{.}\hbox{.}}}}%
  \def\ddots{\mathinner{\mkern1mu\raise6.5pt\hbox{.}\mkern2mu\raise3.6pt\hbox{.}\mkern2mu\raise0.7pt\hbox{.}\mkern1mu}}%
  \def\triangle{\mathord{\scalebox{0.9}{$\symup{\Delta}$}}}%
  \def\nabla{\mathord{\raisebox{\depth}{\scalebox{1}[-1]{$\symup{\Delta}$}}}}%
  \def\checkmark{\popcheck{popdark}}%
  \def\star{\mathbin{\ast}}\def\bigstar{\mathord{\ast}}%
  \def\diamond{\mathbin{\scalebox{0.6}{$\Diamond$}}}%
  \def\bigcup{\mathop{\mathchoice{\raisebox{-0.3em}{\scalebox{1.7}{$\cup$}}}{\scalebox{1.25}{$\cup$}}{\cup}{\cup}}\displaylimits}%
  \def\bigcap{\mathop{\mathchoice{\raisebox{-0.3em}{\scalebox{1.7}{$\cap$}}}{\scalebox{1.25}{$\cap$}}{\cap}{\cap}}\displaylimits}%
  \def\lesseqqgtr{\mathrel{\lessgtr}}\def\bigoplus{\mathop{\oplus}\displaylimits}%
}
\providecommand{\ssi}{si et seulement si} % abréviation fréquente
\AtBeginDocument{\providecommand{\wideparen}{\overparen}} % arc de cercle

%% FIGFIX-BEGIN v4 — correctifs globaux des figures (docs/onbuch-generation/tools/figfix)
\usepackage{adjustbox}\usepackage{etoolbox}
% toute figure TikZ/pgfplots plus large que la ligne est réduite à la largeur disponible
\BeforeBeginEnvironment{tikzpicture}{\begin{adjustbox}{max width=\linewidth}}
\AfterEndEnvironment{tikzpicture}{\end{adjustbox}}
% légendes pgfplots lisibles par-dessus les courbes
\pgfplotsset{every axis/.append style={legend style={fill=white,fill opacity=0.92,text opacity=1,draw=black!15,inner sep=3pt}}}
% étiquettes d'arêtes (syntaxe edge["..."]) avec fond blanc
\tikzset{every edge quotes/.append style={fill=white,inner sep=1.2pt}}
%% FIGFIX-END
\begin{document}

\pagegarde{Expression de l’information génétique}{SVTEEHB}{3ème}{SVTEEHB – classe de 3ème}{N03}{6 séances (fiche de progression harmonisée)}{Cette leçon explore la nature de l'information génétique humaine, portée par les chromosomes et les gènes, et montre comment elle s'exprime à travers un caractère concret : les groupes sanguins ABO et le facteur Rhésus. Elle culmine avec la réalisation d'un caryotype humain et la détermination expérimentale des groupes sanguins.
\vspace{4pt}
\begin{multicols}{2}\small
\begin{itemize}
\item À la fin de cette leçon, je sais décrire la structure et le nombre de chromosomes de l'espèce humaine.
\item À la fin de cette leçon, je sais distinguer autosomes et chromosomes sexuels et déterminer le sexe chromosomique d'un individu à partir de son caryotype.
\item À la fin de cette leçon, je sais construire une maquette de caryotype humain à partir d'une photographie de chromosomes en métaphase.
\item À la fin de cette leçon, je sais définir un gène et un allèle et localiser un gène sur un chromosome.
\item À la fin de cette leçon, je sais expliquer la diversité humaine à partir de la notion d'allèles multiples.
\item À la fin de cette leçon, je sais déterminer expérimentalement un groupe sanguin du système ABO et le facteur Rhésus à l'aide de sérums tests.
\end{itemize}\end{multicols}}

\begin{sommaire}
\begin{multicols}{2}
\begin{enumerate}
\item Les chromosomes de l'espèce humaine
\item TP N°3 : Conception d'une maquette de caryotype humain
\item Les gènes humains
\item Gènes et diversité humaine : le système ABO
\item TP N°4 : Recherche des groupes sanguins ABO et du facteur Rhésus
\item Activités d'intégration et synthèse
\item Activité d'intégration
\item Méthodes et exercices résolus
\item Exercices
\item Corrigés des exercices
\item Fiche bilan
\end{enumerate}
\end{multicols}
\end{sommaire}

\begin{objectifs}
\begin{itemize}
\item À la fin de cette leçon, je sais décrire la structure et le nombre de chromosomes de l'espèce humaine.
\item À la fin de cette leçon, je sais distinguer autosomes et chromosomes sexuels et déterminer le sexe chromosomique d'un individu à partir de son caryotype.
\item À la fin de cette leçon, je sais construire une maquette de caryotype humain à partir d'une photographie de chromosomes en métaphase.
\item À la fin de cette leçon, je sais définir un gène et un allèle et localiser un gène sur un chromosome.
\item À la fin de cette leçon, je sais expliquer la diversité humaine à partir de la notion d'allèles multiples.
\item À la fin de cette leçon, je sais déterminer expérimentalement un groupe sanguin du système ABO et le facteur Rhésus à l'aide de sérums tests.
\item À la fin de cette leçon, je sais interpréter les règles de compatibilité sanguine lors d'une transfusion.
\item À la fin de cette leçon, je sais prédire les groupes sanguins possibles des enfants à partir des génotypes des parents (échiquier de croisement).
\item À la fin de cette leçon, je sais rédiger et justifier une démarche scientifique, une réponse ou une conclusion.
\end{itemize}
\end{objectifs}

\begin{prerequis}
\begin{itemize}
\item Notion de cellule et de noyau cellulaire (4ème).
\item Notion de molécule d'ADN support de l'information génétique (leçon précédente de 3ème).
\item Notion de reproduction sexuée et de fécondation (4ème).
\item Lecture et interprétation de tableaux de résultats et de schémas.
\item Notion de caractère héréditaire observable (ressemblances et différences au sein d'une famille).
\end{itemize}
\end{prerequis}

\newpage

\coursec{Expression de l'information génétique}

\noindent Au centre de transfusion sanguine de l'Hôpital Central de Yaoundé, une urgence vitale arrive : une patiente anémique sévère (hémoglobine à 4,5 g/dL) post-partum a besoin d'une transfusion immédiate. Son groupe sanguin est \textbf{A Rhésus négatif (A-)}. Les poches disponibles sont \textbf{O-}, \textbf{A+}, \textbf{B-} et \textbf{AB-}. Laquelle choisir sans risque ? Pourquoi le groupe O- est-il appelé « donneur universel » et l'AB+ « receveur universel » ? Ces questions, qui se posent chaque jour dans nos hôpitaux camerounais, trouvent leur réponse dans l'\textbf{expression de l'information génétique}. Cette information, portée par les chromosomes, s'exprime par la synthèse de protéines (ici, des antigènes à la surface des globules rouges) qui déterminent nos caractères héréditaires. Commençons par le support physique de cette information : les chromosomes.

\coursec{Les chromosomes de l'espèce humaine}

\courssub{Mise en évidence des chromosomes : observation au microscope}

Les chromosomes ne sont pas visibles en permanence dans le noyau des cellules en interphase (chromatine décondensée). Pour les observer et les compter, il faut les mettre en évidence lors de la division cellulaire (mitose), plus précisément au stade de la \cle{métaphase}, où ils sont maximalement condensés (courts et épais).

\begin{experience}[Observation des chromosomes humains en métaphase (sur document photographique)]
\textbf{Matériel :} Photographies microscopiques de lymphocytes humains cultivés \emph{in vitro} (prélèvement sanguin veineux), colorés au Giemsa (technique de bandement G).
\textbf{Protocole (rappel technique) :}
\begin{enumerate}
    \item \textbf{Culture :} Lymphocytes mis en culture 72 h avec un mitogène (phytohémagglutinine, PHA) pour stimuler l'entrée en mitose.
    \item \textbf{Blocage métaphasique :} Ajout de colchicine (inhibiteur de la polymérisation des microtubules) qui détruit le fuseau mitotique $\rightarrow$ accumulation des cellules en métaphase.
    \item \textbf{Choc hypotonique :} Traitement par une solution hypotonique (eau distillée ou KCl 0,075 M) $\rightarrow$ entrée d'eau par osmose, gonflement cellulaire, écartement des chromosomes.
    \item \textbf{Fixation et étalement :} Fixateur (méthanol/acide acétique 3/1), goutte sur lame froide $\rightarrow$ éclatement de la membrane, étalement des chromosomes.
    \item \textbf{Coloration et observation :} Coloration Giemsa (bandes G : régions AT riches en adénine-thymine = bandes sombres ; régions GC riches en guanine-cytosine = bandes claires). Microscope optique $\times 1000$ (immersion d'huile).
\end{enumerate}
\textbf{Observations :} Structures filamenteuses, sombres, dispersées. Chaque chromosome présente une \cle{constriction primaire} (centromère) divisant le chromosome en deux bras. Morphologies variées : en V (métacentriques), en J (sous-métacentriques), en tige (acrocentriques).
\textbf{Interprétation :} La condensation de la chromatine (ADN + histones) en prophase rend les chromosomes individuels discernables. Le bandement G offre un « code-barres » unique pour chaque paire homologue.
\textbf{Conclusion :} Les chromosomes sont les supports physiques de l'information génétique, observables au microscope optique uniquement en métaphase de mitose.
\end{experience}

\begin{popfigure}
\begin{tikzpicture}[scale=1.3, >=Stealth]
    % Chromatide sœur 1 (haut)
    \draw[popdark, thick, fill=popblue!15] (0,0) .. controls (0.4,0.4) and (1.6,0.4) .. (2,0) -- (2,-0.35) .. controls (1.6,-0.35) and (0.4,-0.35) .. (0,-0.35) -- cycle;
    % Chromatide sœur 2 (bas)
    \draw[popdark, thick, fill=poporange!15] (0,-0.55) .. controls (0.4,-0.95) and (1.6,-0.95) .. (2,-0.55) -- (2,-0.9) .. controls (1.6,-0.9) and (0.4,-0.9) .. (0,-0.9) -- cycle;
    % Centromère
    \draw[poppurple, thick, fill=poppurple!40] (0.95,-0.35) rectangle (1.05,-0.55);
    % Bras courts (p) et longs (q)
    \node[popdark, font=\bfseries\small] at (0.4,0.25) {p};
    \node[popdark, font=\bfseries\small] at (0.4,-1.2) {p};
    \node[popdark, font=\bfseries\small] at (1.7,0.25) {q};
    \node[popdark, font=\bfseries\small] at (1.7,-1.2) {q};
    % Annotations
    \draw[<-, popdark, shorten >=2pt] (2.1,0.1) -- (2.7,0.1) node[right, font=\small] {Chromatide sœur 1 (réplique)};
    \draw[<-, popdark, shorten >=2pt] (2.1,-0.8) -- (2.7,-0.8) node[right, font=\small] {Chromatide sœur 2 (réplique)};
    \draw[<-, poppurple, shorten >=2pt] (1.0,-1.1) -- (1.0,-1.5) node[below, font=\small] {Centromère (constriction primaire)};
    \draw[<-, popdark, shorten >=2pt] (-0.4,0.1) -- (-1.2,0.1) node[left, font=\small] {Télomères};
    \draw[<-, popdark, shorten >=2pt] (-0.4,-0.8) -- (-1.2,-0.8) node[left, font=\small] {Télomères};
\end{tikzpicture}
\legende{Schéma d'un chromosome métaphasique humain : deux chromatides sœurs identiques (issues de la réplication de l'ADN en phase S) unies par le centromère. Bras court \textbf{p} (petit), bras long \textbf{q} (queue).}
\end{popfigure}

\begin{definition}[Chromosome]
Un \textbf{chromosome} est une structure filamenteuse, visible au microscope optique uniquement lors de la division cellulaire (métaphase), constituée d'une molécule d'ADN double brin très compactée (enroulée autour de protéines basiques : les \textbf{histones}, formant la \textbf{chromatine}). Il porte une partie de l'information génétique : les \textbf{gènes}.
\end{definition}

\begin{attention}[Chromatide $\neq$ Chromosome : piège classique au BEPC]
Ne confonds jamais :
\begin{itemize}
    \item \textbf{Chromosome} : entité génétique comptée \textbf{par centromère}. En métaphase, il est \textbf{double} (2 chromatides sœurs).
    \item \textbf{Chromatide} : l'un des deux brins d'ADN identiques d'un chromosome double. Après l'anaphase, chaque chromatide devient un chromosome \textbf{simple} (1 chromatide = 1 chromosome).
\end{itemize}
\textbf{Exemple :} En métaphase, une cellule humaine a \textbf{46 chromosomes} (46 centromères) et \textbf{92 chromatides} (92 molécules d'ADN). Dire « 92 chromosomes » est \textbf{FAUX}.
\end{attention}

\courssub{Nombre de chromosomes et notion de paires homologues}

Le comptage systématique sur de nombreuses cellules somatiques saines (fibroblastes, lymphocytes, cellules de moelle osseuse) donne un nombre constant, caractéristique de l'espèce.

\begin{propriete}[Nombre chromosomique de l'espèce humaine (admis)]
Toute cellule \textbf{somatique} (non reproductrice) humaine à noyau contient \textbf{46 chromosomes}, organisés en \textbf{23 paires}. On note \textbf{2n = 46} (caryotype \textbf{diploïde}).
Seuls les \textbf{gamètes} (spermatozoïdes et ovocytes) contiennent 23 chromosomes (n = 23, caryotype \textbf{haploïde}).
\end{propriete}

\begin{savaistu}[L'histoire du nombre 46 : une erreur de 30 ans !]
Jusqu'en 1956, les manuels enseignaient 48 chromosomes (comme le chimpanzé). C'est grâce à la technique du choc hypotonique (améliorant l'étalement) que \textbf{Joe Hin Tjio} et \textbf{Albert Levan} (Lund, Suède) ont établi le nombre exact : \textbf{46}. La science progresse par la technique !
\end{savaistu}

Ces 23 paires ne sont pas 23 chromosomes différents deux à deux. Elles forment des \cle{paires de chromosomes homologues}.

\begin{definition}[Chromosomes homologues]
Deux chromosomes forment une \textbf{paire homologue} s'ils partagent :
\begin{itemize}
    \item La \textbf{même taille} (longueur totale en µm),
    \item La \textbf{même forme} (position du centromère : métacentrique, sous-métacentrique, acrocentrique),
    \item Les \textbf{mêmes gènes} aux mêmes \textbf{loci} (positions), bien que les \textbf{allèles} (versions alléliques du gène) puissent différer.
\end{itemize}
L'un des homologues est d'origine \textbf{paternelle} (venant du spermatozoïde), l'autre d'origine \textbf{maternelle} (venant de l'ovocyte).
\end{definition}

\begin{exemplebox}[L'exemple du gène ABO (chromosome 9)]
Le gène \emph{ABO} se situe sur le bras long du chromosome 9 (locus 9q34.2). Sur la paire d'homologues (deux chromosomes 9) :
\begin{itemize}
    \item Chromosome 9 paternel : allèle $I^A$ (code pour antigène A).
    \item Chromosome 9 maternel : allèle $i$ (ne code pas d'antigène A/B).
\end{itemize}
Même gène, même locus, allèles différents $\rightarrow$ Génotype $I^A/i$ $\rightarrow$ Phénotype \textbf{Groupe A}.
\end{exemplebox}

\courssub{Autosomes et chromosomes sexuels : détermination du sexe}

Parmi les 23 paires, 22 sont identiques chez l'homme et la femme : ce sont les \cle{autosomes} (chromosomes somatiques), numérotés de 1 à 22 par taille décroissante. La 23\up{ème} paire détermine le sexe chromosomique : ce sont les \cle{chromosomes sexuels} (gonosomes).

\begin{itemize}
    \item \textbf{Femme (46,XX) :} Deux chromosomes X homologues (même taille, même bandement, mêmes gènes).
    \item \textbf{Homme (46,XY) :} Un chromosome X et un chromosome Y. Le Y est beaucoup plus petit ; ils sont \textbf{non homologues} (hétérologues), sauf au niveau des \textbf{régions pseudo-autosomiques} (PAR1 et PAR2) où ils s'apparient en méiose.
\end{itemize}

\begin{propriete}[Détermination du sexe chromosomique (admis)]
C'est la présence du \textbf{chromosome Y} qui détermine le développement mâle. Le chromosome Y porte le gène \textbf{SRY} (\emph{Sex-determining Region Y}), facteur de transcription déclenchant la différenciation de la gonade indifférenciée en \textbf{testicule} (vers 7\up{ème} semaine embryonnaire). Sans SRY (génotype XX), la gonade se différencie en \textbf{ovaire} par voie de développement par défaut.
\end{propriete}

\begin{popfigure}
\begin{tikzpicture}[scale=0.65, >=Stealth]
    \tikzset{chromo/.style={draw=popdark, thick, fill=popblue!10, minimum width=0.55cm, rounded corners=0.08cm}}
    \tikzset{chromoX/.style={chromo, fill=poppink!30}}
    \tikzset{chromoY/.style={draw=popdark, thick, fill=poporange!30, minimum width=0.35cm, rounded corners=0.08cm}}
    \tikzset{label/.style={font=\tiny, text=popdark, anchor=east, xshift=-0.2cm}}
    % Titre
    \node[font=\bfseries\small] at (0,9.5) {Caryotype humain schématisé (classement Denver modifié, bandes G simplifiées)};
    % Groupes A à G : on dessine UN représentant par paire (alignés par centromère simulé)
    % Coordonnées Y de base pour chaque paire
    \foreach \num/\y/\h/\g in {
        1/8.8/2.2/A, 2/8.3/2.1/A, 3/7.8/2.0/A,
        4/7.3/1.8/B, 5/6.9/1.7/B,
        6/6.4/1.5/C, 7/6.0/1.4/C, 8/5.6/1.3/C, 9/5.2/1.3/C, 10/4.8/1.2/C, 11/4.4/1.1/C, 12/4.0/1.1/C,
        13/3.5/1.0/D, 14/3.1/0.9/D, 15/2.7/0.9/D,
        16/2.2/0.8/E, 17/1.8/0.7/E, 18/1.4/0.7/E,
        19/0.9/0.6/F, 20/0.5/0.5/F,
        21/0.0/0.4/G, 22/-0.4/0.4/G} {
        \node[chromo, minimum height=\h cm] (c\num a) at (-2.2,\y) {};
        \node[chromo, minimum height=\h cm] (c\num b) at (-1.3,\y) {};
        \node[label] at (-2.7,\y) {Paire \num};
        % Simuler centromère par un trait fin
        \draw[poppurple, line width=0.5pt] (-2.475,\y) -- (-1.925,\y);
        \draw[poppurple, line width=0.5pt] (-1.575,\y) -- (-1.025,\y);
    }
    % Chromosomes sexuels Femme (XX) - Paire 23
    \node[chromoX, minimum height=1.5cm] (cX a) at (-2.2,-1.0) {};
    \node[chromoX, minimum height=1.5cm] (cX b) at (-1.3,-1.0) {};
    \node[label, text=poppink] at (-2.7,-1.0) {XX (Femme)};
    \draw[poppurple, line width=0.5pt] (-2.475,-1.0) -- (-1.925,-1.0);
    \draw[poppurple, line width=0.5pt] (-1.575,-1.0) -- (-1.025,-1.0);
    % Chromosomes sexuels Homme (XY) - Paire 23
    \node[chromoX, minimum height=1.5cm] (cX2) at (-2.2,-1.6) {};
    \node[chromoY, minimum height=0.5cm] (cY) at (-1.3,-1.6) {};
    \node[label, text=popblue] at (-2.7,-1.6) {XY (Homme)};
    \draw[poppurple, line width=0.5pt] (-2.475,-1.6) -- (-1.925,-1.6);
    \draw[poppurple, line width=0.5pt] (-1.575,-1.6) -- (-1.025,-1.6); % centromere Y simulé
    % Accolades groupes
    \draw[popmuted, decorate, decoration={brace, amplitude=4pt}] (-2.9,9.0) -- (-2.9,7.5) node[midway, left=6pt, rotate=90, font=\tiny] {Groupe A (1-3) : Grands métacentriques};
    \draw[popmuted, decorate, decoration={brace, amplitude=4pt}] (-2.9,7.5) -- (-2.9,3.8) node[midway, left=6pt, rotate=90, font=\tiny] {Groupes B, C (4-12) : Sous-métacentriques};
    \draw[popmuted, decorate, decoration={brace, amplitude=4pt}] (-2.9,3.8) -- (-2.9,0.2) node[midway, left=6pt, rotate=90, font=\tiny] {Groupes D, E (13-18) : Acrocentriques / Sous-métacentriques};
    \draw[popmuted, decorate, decoration={brace, amplitude=4pt}] (-2.9,0.2) -- (-2.9,-1.8) node[midway, left=6pt, rotate=90, font=\tiny] {Groupes F, G (19-22) : Petits métacentriques / Acrocentriques};
    % Légende
    \node[font=\small, align=left, anchor=north west] at (1.5, 9.0) {
        \textbf{Classification Denver (modifiée) :}\\
        \textcolor{popblue}{$\blacksquare$} Groupe A : 1-3 (Grands métacentriques)\\
        \textcolor{poporange}{$\blacksquare$} Groupe B : 4-5 (Grands sous-métacentriques)\\
        \textcolor{popgreen}{$\blacksquare$} Groupe C : 6-12 + X (Moyens sous-métacentriques)\\
        \textcolor{poppurple}{$\blacksquare$} Groupe D : 13-15 (Moyens acrocentriques)\\
        \textcolor{poppink}{$\blacksquare$} Groupe E : 16-18 (Petits sous-métacentriques)\\
        \textcolor{popgold}{$\blacksquare$} Groupe F : 19-20 (Petits métacentriques)\\
        \textcolor{popdark}{$\blacksquare$} Groupe G : 21-22 + Y (Très petits acrocentriques)
    };
\end{tikzpicture}
\legende{Caryotype humain schématisé. Les 22 paires d'autosomes classées par taille décroissante et position du centromère (groupes A à G). Paire 23 : XX (femme, rose) ou XY (homme, bleu/orange). Le bandement G (alternance bandes claires/sombres) permet l'identification certaine de chaque paire.}
\end{popfigure}

\courssub{Le caryotype : définition, notation et réalisation (TP N°3)}

Le \cle{caryotype} est la « carte d'identité » chromosomique d'un individu.

\begin{definition}[Caryotype]
Le \textbf{caryotype} est le classement ordonné de l'ensemble des chromosomes d'une cellule (généralement en métaphase) par \textbf{taille décroissante}, \textbf{forme} (position du centromère) et \textbf{profil de bandement} (bandes G). Il se note sous la forme standardisée (ISCN) : \textbf{46,XX} (femme) ou \textbf{46,XY} (homme).
\end{definition}

\begin{experience}[TP N°3 : Conception d'une maquette de caryotype humain]
\textbf{Objectif :} Réaliser le caryotype d'un individu à partir d'une photographie de chromosomes dispersés en métaphase.
\textbf{Matériel :} Photographie agrandie (A4) d'une métaphase humaine (bandes G), ciseaux, colle, feuille blanche format A3, règle, crayon.
\textbf{Protocole :}
\begin{enumerate}
    \item \textbf{Découpage :} Découper soigneusement chaque chromosome individuel en laissant une marge blanche.
    \item \textbf{Appariement (recherche des homologues) :} Regrouper les chromosomes deux par deux en comparant : (1) Longueur totale, (2) Position du centromère (rapport bras court/bras long), (3) Profil de bandement G (alternance bandes claires/sombres). $\rightarrow$ 23 paires.
    \item \textbf{Classement :} Ordonner les paires par taille décroissante (1 à 22). Respecter les groupes (A à G) selon la position du centromère.
    \item \textbf{Chromosomes sexuels :} Identifier la paire 23 (XX ou XY) et la placer à la fin (ou en position 23).
    \item \textbf{Collage et notation :} Coller les paires alignées par leur centromère. Écrire la formule caryotypique (ex: \textbf{46,XY}).
\end{enumerate}
\textbf{Résultat attendu :} Une maquette propre, 23 paires alignées, formule correcte.
\textbf{Compétences évaluées :} Rigueur observationnelle, organisation logique, motricité fine, respect de la nomenclature internationale.
\end{experience}

\begin{exoresolu}[Analyse de formules caryotypiques (type BEPC)]
\textbf{Énoncé :} On donne trois formules caryotypiques :
\begin{itemize}
    \item Individu 1 : 46,XX
    \item Individu 2 : 46,XY
    \item Individu 3 : 47,XY,+21
\end{itemize}
\textbf{Questions :}
\begin{enumerate}
    \item Préciser le sexe chromosomique de chaque individu.
    \item Interpréter l'anomalie de l'individu 3 (nom, mécanisme, conséquences).
    \item Calculer le nombre de chromatides en métaphase pour l'individu 3.
\end{enumerate}
\tcblower
\textbf{Solution détaillée :}
\begin{enumerate}
    \item \textbf{Individu 1 :} 46 chromosomes, deux X $\rightarrow$ \textbf{Femme}.\\
          \textbf{Individu 2 :} 46 chromosomes, un X et un Y $\rightarrow$ \textbf{Homme}.\\
          \textbf{Individu 3 :} Présence d'un Y $\rightarrow$ \textbf{Homme}.
    \item \textbf{Anomalie :} \textbf{Trisomie 21} (syndrome de Down). Notation \texttt{+21} = chromosome en excès sur la paire 21 (3 exemplaires au lieu de 2).
          \textbf{Mécanisme (admis) :} \textbf{Non-disjonction} de la paire 21 lors de la méiose I (séparation des homologues) ou II (séparation des chromatides) $\rightarrow$ gamète à 24 chromosomes (n+1) + gamète normal (n) $\rightarrow$ zygote 47 chromosomes (2n+1).
          \textbf{Conséquences :} Retard mental modéré, traits faciaux caractéristiques (pli épicanthique, fentes palpébrales obliques), hypotonie musculaire, cardiopathies congénitales fréquentes (communication interventriculaire).
    \item En métaphase, chaque chromosome est double (2 chromatides). Nombre total de chromosomes = 47. Nombre de chromatides = $47 \times 2 = \textbf{94 chromatides}$.
\end{enumerate}
\end{exoresolu}

\courssub{Structure fine du chromosome métaphasique}

\begin{propriete}[Structure du chromosome métaphasique]
Un chromosome métaphasique est un \textbf{chromosome double} constitué de :
\begin{itemize}
    \item \textbf{Deux chromatides sœurs} : identiques (issues de la réplication semi-conservative de l'ADN en phase S), chacune contenant \textbf{une molécule d'ADN double brin} continue.
    \item \textbf{Un centromère} (région constrictive) : point d'attache du kinétochore (protéines) où s'insèrent les microtubules du fuseau mitotique. Il porte l'origine de réplication et les séquences centromériques (ADN satellite $\alpha$).
    \item \textbf{Deux télomères} (extrémités) : séquences répétitives non codantes (TTAGGG)$_n$ chez l'homme, protectrices contre la dégradation et la fusion chromosomique.
\end{itemize}
\end{propriete}

\begin{attention}[Distinction cruciale : Paire d'homologues vs Chromosome double]
\begin{center}
\begin{tabularx}{\linewidth}{|l|X|X|}
\hline
\rowcolor{popblueL}
\textbf{Critère} & \textbf{Chromosome double (en métaphase)} & \textbf{Paire d'homologues (dans le noyau 2n)} \\
\hline
Nombre de centromères & 1 & 2 \\
\hline
Nombre de chromatides & 2 (sœurs, ADN identique) & 4 (2 paternelles + 2 maternelles) \\
\hline
Origine de l'ADN & Une seule origine parentale & Origine biparentale (paternel + maternel) \\
\hline
Allèles aux mêmes loci & Identiques (sauf mutation de novo) & Potentiellement différents (hétérozygotie) \\
\hline
Exemple visible & Un seul « objet » sur la photo de caryotype & Deux « objets » appariés sur le caryotype \\
\hline
\end{tabularx}
\end{center}
\end{attention}

\courssub{Anomalies numériques : la trisomie 21}

\begin{center}
\begin{tabularx}{\linewidth}{|l|X|X|}
\hline
\rowcolor{popblueL}
\textbf{Paramètre} & \textbf{Caryotype normal (46,XX / 46,XY)} & \textbf{Trisomie 21 libre (47,XX,+21 / 47,XY,+21)} \\
\hline
Nombre total chromosomes & 46 (2n) & 47 (2n+1) \\
\hline
Paire 21 (Groupe G) & 2 chromosomes homologues & 3 chromosomes (triplet) \\
\hline
Notation ISCN & 46,XX ou 46,XY & 47,XX,+21 ou 47,XY,+21 \\
\hline
Mécanisme principal & Disjonction normale en méiose & \textbf{Non-disjonction méiotique} (majoritairement méiose I maternelle) \\
\hline
Facteur de risque & — & \textbf{Âge maternel avancé} ($>35$ ans : risque $\uparrow$ significativement) \\
\hline
Fréquence à la naissance & — & $\approx$ 1 / 700 naissances vivantes (variable selon pays) \\
\hline
Phénotype principal & Développement standard & Syndrome de Down (cf. exoresolu) \\
\hline
\end{tabularx}
\end{center}

\begin{aretenir}[Synthèse : Les chromosomes de l'espèce humaine]
\begin{itemize}
    \item Support de l'information génétique : \textbf{chromosomes} (ADN + histones), visibles en \textbf{métaphase}.
    \item Espèce humaine : \textbf{2n = 46 chromosomes} (cellules somatiques), organisés en \textbf{23 paires d'homologues} (22 paires d'autosomes + 1 paire de gonosomes).
    \item Sexe chromosomique : \textbf{Femme 46,XX} | \textbf{Homme 46,XY}. Détermination par le gène \textbf{SRY} sur le chromosome Y.
    \item \textbf{Caryotype} : classement standardisé (taille, centromère, bandement G). Outil diagnostic (trisomies, translocations, délétions).
    \item En métaphase : chromosome \textbf{double} = 2 chromatides sœurs (ADN répliqué) unies au centromère.
    \item \textbf{Trisomie 21} : 3 chromosomes 21 (47 total), non-disjonction méiotique, risque maternel âge-dépendant.
\end{itemize}
\end{aretenir}

\coursec{Les gènes humains}

\courssub{Du chromosome au gène : définition et localisation}

\begin{definition}[Gène (niveau 3ème)]
Un \textbf{gène} est une \textbf{portion de la molécule d'ADN} (une séquence de nucléotides) située à un endroit précis (\textbf{locus}) sur un chromosome, qui contient l'information nécessaire à la synthèse d'un \textbf{produit fonctionnel} : le plus souvent une \textbf{protéine} (enzyme, protéine structurale, récepteur, hormone...), parfois un ARN fonctionnel (ARNt, ARNr, ARNm régulateur).
\end{definition}

\begin{propriete}[Organisation du génome humain (ordres de grandeur admis)]
\begin{itemize}
    \item Taille du génome haploïde : $\approx$ \textbf{3,2 milliards de paires de bases} (3,2 Gb).
    \item Nombre de gènes codant des protéines : $\approx$ \textbf{20 000 à 25 000}.
    \item Densité génique variable : chromosome 19 (riche), chromosome 13, 18, 21, Y (pauvres).
    \item ADN non codant (introns, séquences régulatrices, ADN répétitif) : $\approx$ 98\% du génome.
\end{itemize}
\end{propriete}

\begin{exemplebox}[Exemples de gènes humains localisés]
\begin{center}
\begin{tabularx}{\linewidth}{|l|l|X|}
\hline
\rowcolor{popblueL}
\textbf{Gène} & \textbf{Localisation (Locus)} & \textbf{Fonction / Pathologie si muté} \\
\hline
\emph{HBB} (bêta-globine) & 11p15.5 & Sous-unité de l'hémoglobine. Mutation $\rightarrow$ Drépanocytose (SS), Thalassémie. \\
\hline
\emph{CFTR} & 7q31.2 & Canal chlorure. Mutation $\rightarrow$ Mucoviscidose. \\
\hline
\emph{ABO} & 9q34.2 & Glycosyltransférase $\rightarrow$ Antigènes A/B (Groupes sanguins). \\
\hline
\emph{RHD} & 1p36.11 & Protéine RhD (Facteur Rhésus). Délétion $\rightarrow$ Rhésus négatif. \\
\hline
\emph{SRY} & Yp11.3 & Facteur de transcription déterminant le testicule. \\
\hline
\end{tabularx}
\end{center}
\end{exemplebox}

\courssub{Allèles, génotype et phénotype}

\begin{definition}[Allèle]
Un \textbf{allèle} est une \textbf{version alternative} d'un même gène, occupant le \textbf{même locus} sur des chromosomes homologues, différant par sa séquence nucléotidique (mutation). Exemple : pour le gène \emph{ABO}, les allèles sont $I^A$, $I^B$, $i$.
\end{definition}

\begin{definition}[Génotype et Phénotype]
\begin{itemize}
    \item \textbf{Génotype} : Composition allélique d'un individu pour un gène donné (ex: $I^A/i$). Se note entre crochets $[I^A//i]$ ou en ligne $I^A/i$.
    \item \textbf{Phénotype} : Expression observable du génotype (caractère visible, mesurable) dans un environnement donné (ex: Groupe sanguin A).
\end{itemize}
\end{definition}

\begin{propriete}[Relations de dominance / récessivité / codominance (admis)]
Pour un gène à deux allèles chez un individu diploïde :
\begin{itemize}
    \item \textbf{Dominance complète :} Allèle $A$ dominant sur $a$ récessif. $A/A$ et $A/a$ $\rightarrow$ même phénotype (A). $a/a$ $\rightarrow$ phénotype (a).
    \item \textbf{Codominance :} Allèles $A$ et $B$ codominants. $A/B$ $\rightarrow$ phénotype exprimant les deux produits (AB).
    \item \textbf{Dominance incomplète :} Phénotype intermédiaire chez l'hétérozygote (ex: drépanocytose AS $\rightarrow$ trait drépanocytaire, résistant au paludisme).
\end{itemize}
\end{propriete}

\coursec{Gènes et diversité humaine : le système ABO et le facteur Rhésus}

\courssub{Le système ABO : génétique à trois allèles}

Le système ABO est déterminé par un seul gène (\emph{ABO}, chromosome 9) qui possède \textbf{trois allèles} dans la population : $I^A$, $I^B$, $i$ (ou $I^O$).

\begin{propriete}[Biochimie des antigènes ABO (admis)]
\begin{itemize}
    \item Allèle $I^A$ : code une glycosyltransférase ajoutant un \textbf{N-acétylgalactosamine} (GalNAc) sur le précurseur H (antigène H) $\rightarrow$ \textbf{Antigène A}.
    \item Allèle $I^B$ : code une glycosyltransférase ajoutant un \textbf{galactose} sur l'antigène H $\rightarrow$ \textbf{Antigène B}.
    \item Allèle $i$ (récessif) : allèle nul (mutation sans fonction enzymatique) $\rightarrow$ pas de modification de l'antigène H $\rightarrow$ \textbf{Antigène H seul} (Groupe O).
\end{itemize}
$I^A$ et $I^B$ sont \textbf{codominants} l'un par rapport à l'autre. $I^A$ et $I^B$ sont \textbf{dominants} sur $i$.
\end{propriete}

\begin{center}
\begin{tabular}{|c|c|c|c|}
\hline
\rowcolor{popblueL}
\textbf{Phénotype (Groupe)} & \textbf{Génotypes possibles} & \textbf{Antigènes sur Globules Rouges (GR)} & \textbf{Anticorps (Ac) dans le sérum (anti-A, anti-B)} \\
\hline
\textbf{A} & $I^A I^A$ ; $I^A i$ & A (et H) & Anti-B \\
\hline
\textbf{B} & $I^B I^B$ ; $I^B i$ & B (et H) & Anti-A \\
\hline
\textbf{AB} & $I^A I^B$ & A et B (peu de H) & \textbf{Aucun} (ni anti-A, ni anti-B) \\
\hline
\textbf{O} & $i i$ & H seul (ni A, ni B) & Anti-A \textbf{et} Anti-B \\
\hline
\end{tabular}
\end{center}

\begin{attention}[Règle de la transfusion ABO (Beth-Vincent) : « Le receveur ne doit pas avoir d'anticorps contre les antigènes du donneur »]
\begin{itemize}
    \item \textbf{Donneur universel : Groupe O} (GR sans A ni B $\rightarrow$ pas d'agglutination par les anti-A/anti-B du receveur). \emph{Attention : sérum du donneur O contient anti-A + anti-B $\rightarrow$ risque si gros volume.}
    \item \textbf{Receveur universel : Groupe AB} (Sérum sans anti-A ni anti-B $\rightarrow$ n'agglutine pas les GR du donneur).
    \item \textbf{Compatibilité :} A donne à A et AB. B donne à B et AB. AB donne à AB seulement. O donne à tous.
\end{itemize}
\end{attention}

\courssub{Le système Rhésus (facteur Rh) : génétique à deux allèles}

Le système Rhésus est complexe (plus de 50 antigènes), mais le facteur principal est l'antigène \textbf{D} (RhD).

\begin{propriete}[Génétique du facteur Rhésus D (admis)]
\begin{itemize}
    \item Gène \emph{RHD} (chromosome 1) et gène voisin \emph{RHCE}.
    \item Allèle \textbf{$RHD^+$ (ou $D$)} : fonctionnel $\rightarrow$ synthèse protéine RhD $\rightarrow$ \textbf{Rhésus Positif (Rh+)}.
    \item Allèle \textbf{$RHD^-$ (ou $d$)} : le plus souvent \textbf{délétion complète du gène} $\rightarrow$ pas de protéine RhD $\rightarrow$ \textbf{Rhésus Négatif (Rh-)}.
    \item $D$ est \textbf{dominant} sur $d$.
\end{itemize}
\end{propriete}

\begin{center}
\begin{tabular}{|c|c|c|}
\hline
\rowcolor{popblueL}
\textbf{Phénotype Rhésus} & \textbf{Génotype} & \textbf{Fréquence approx. (Cameroun / Afrique)} \\
\hline
\textbf{Rhésus Positif (Rh+)} & $D/D$ ou $D/d$ & $\approx$ 95\% à 99\% (très majoritaire) \\
\hline
\textbf{Rhésus Négatif (Rh-)} & $d/d$ & $\approx$ 1\% à 5\% (rare) \\
\hline
\end{tabular}
\end{center}

\begin{savaistu}[Contexte camerounais : Rhésus négatif et grossesse]
Au Cameroun, comme en Afrique sub-saharienne, le phénotype Rh- est rare ($<5\%$). Une femme Rh- enceinte d'un fœtus Rh+ (père Rh+) risque une \textbf{immunisation anti-D} lors de l'accouchement (passage de GR fœtaux dans la circulation maternelle). Cela peut provoquer une \textbf{maladie hémolytique du nouveau-né} (ictère grave, anémie) lors des grossesses suivantes. La prophylaxie par \textbf{immunoglobulines anti-D} (RhoGAM) dans les 72h post-partum (et à 28 SA) est essentielle et disponible dans nos maternités (ex: Hôpital Laquintinie Douala, HGOPY Yaoundé).
\end{savaistu}

\courssub{Compatibilité transfusionnelle complète (ABO + Rhésus)}

La règle combine les deux systèmes : \textbf{Donneur compatible = Antigènes du donneur absents des anticorps du receveur}.

\begin{center}
\begin{tabularx}{\linewidth}{|c|X|X|}
\hline
\rowcolor{popblueL}
\textbf{Groupe Receveur} & \textbf{Groupes Donneurs compatibles (GR concentrés)} & \textbf{Groupes Receveurs compatibles (Don de GR)} \\
\hline
\textbf{O-} & O- & \textbf{Tous (O-, O+, A-, A+, B-, B+, AB-, AB+)} $\rightarrow$ \textbf{Donneur universel GR} \\
\hline
\textbf{O+} & O-, O+ & O+, A+, B+, AB+ \\
\hline
\textbf{A-} & O-, A- & A-, A+, AB-, AB+ \\
\hline
\textbf{A+} & O-, O+, A-, A+ & A+, AB+ \\
\hline
\textbf{B-} & O-, B- & B-, B+, AB-, AB+ \\
\hline
\textbf{B+} & O-, O+, B-, B+ & B+, AB+ \\
\hline
\textbf{AB-} & O-, A-, B-, AB- & AB-, AB+ \\
\hline
\textbf{AB+} & \textbf{Tous} $\rightarrow$ \textbf{Receveur universel GR} & AB+ seulement \\
\hline
\end{tabularx}
\end{center}

\begin{exemplebox}[Retour au cas clinique de Yaoundé]
Patiente \textbf{A-} (Antigènes A ; Ac anti-B, anti-A ? non, anti-B ; Rh- : pas Ag D, Ac anti-D possible si immunisée).
Pochoirs disponibles : O-, A+, B-, AB-.
\begin{itemize}
    \item \textbf{O-} : Pas Ag A, pas Ag B, pas Ag D $\rightarrow$ \textbf{COMPATIBLE} (Choix idéal).
    \item \textbf{A+} : Ag A OK, mais Ag D présent $\rightarrow$ \textbf{INCOMPATIBLE} (Risque immunisation anti-D ou réaction si déjà immunisée).
    \item \textbf{B-} : Ag B présent $\rightarrow$ Agglutination par anti-B du receveur $\rightarrow$ \textbf{INCOMPATIBLE}.
    \item \textbf{AB-} : Ag A + Ag B $\rightarrow$ Agglutination par anti-B $\rightarrow$ \textbf{INCOMPATIBLE}.
\end{itemize}
\textbf{Décision : Transfuser O- (Donneur universel).}
\end{exemplebox}

\coursec{TP N°4 : Recherche des groupes sanguins (ABO) et du facteur Rhésus}

\begin{experience}[TP N°4 : Détermination du phénotype sanguin par hémagglutination (Méthode Beth-Vincent / Test sur plaque)]
\textbf{Objectif :} Identifier le groupe ABO et le Rhésus d'un échantillon de sang total (ex: sang de l'élève, ou échantillons témoins).
\textbf{Matériel :}
\begin{itemize}
    \item Sang total frais (prélèvement au doigt, hépateine) ou GR lavés.
    \item Sérums de groupage (réactifs) : \textbf{Anti-A} (sérum de groupe B), \textbf{Anti-B} (sérum de groupe A), \textbf{Anti-A,B} (contrôle), \textbf{Anti-D} (monoclonal IgM/IgG).
    \item Plaque de groupage (tuile en verre ou plaque à puits), baguettes de mélange (stylo), chronomètre.
    \item Microscope (optionnel, pour confirmer agglutination faible).
\end{itemize}
\textbf{Protocole (sur plaque, température ambiante) :}
\begin{enumerate}
    \item Identifier les puits : \textbf{A}, \textbf{B}, \textbf{AB} (ou témoin), \textbf{D (Rh)}.
    \item Déposer \textbf{1 goutte} de chaque réactif dans son puits respectif (Anti-A dans A, Anti-B dans B, Anti-D dans D).
    \item Déposer \textbf{1 goutte de sang} (ou suspension de GR à 30-50\%) à côté de chaque réactif (sans toucher le flacon).
    \item \textbf{Mélanger} intimement réactif + sang avec une baguette propre par puits (étaler sur $\approx$ 2 cm).
    \item \textbf{Incuber} 2 à 5 min (selon notice fabricant). Incliner doucement la plaque (rocking) pour favoriser la réaction.
    \item \textbf{Lire} à l'œil nu (ou loupe $\times 10$) sur fond blanc : \textbf{Agglutination} = grumeaux rouges visibles (réaction +) ; \textbf{Non agglutination} = tache rouge lisse homogène (réaction -).
    \item \textbf{Noter} les résultats dans un tableau.
\end{enumerate}
\end{experience}

\textbf{Interprétation des résultats (Tableau de lecture) :}

\begin{center}
\begin{tabular}{|c|c|c|c|c|}
\hline
\rowcolor{popblueL}
\textbf{Réaction Anti-A} & \textbf{Réaction Anti-B} & \textbf{Réaction Anti-D} & \textbf{Groupe ABO} & \textbf{Rhésus} \\
\hline
+ & - & + & A & Rh+ \\
\hline
+ & - & - & A & Rh- \\
\hline
- & + & + & B & Rh+ \\
\hline
- & + & - & B & Rh- \\
\hline
+ & + & + & AB & Rh+ \\
\hline
+ & + & - & AB & Rh- \\
\hline
- & - & + & O & Rh+ \\
\hline
- & - & - & O & Rh- \\
\hline
\end{tabular}
\end{center}

\textbf{Contrôles de qualité (obligatoires) :}
\begin{itemize}
    \item \textbf{Auto-contrôle :} Puits « Témoin » (Sérum AB du patient + ses propres GR) $\rightarrow$ Doit être \textbf{Négatif} (exclut auto-anticorps / groupe sanguin rare type Bombay).
    \item \textbf{Contrôle réactifs :} Vérifier réactifs sur GR témoins connus (A1, B, O, Rh+).
\end{itemize}

\begin{attention}[Pièges techniques fréquents au TP]
\begin{itemize}
    \item \textbf{Faux positif :} Fibrine (caillot) confondue avec agglutination $\rightarrow$ Bien mélanger, utiliser sang hépariné/EDTA, observer au microscope (agrégats fibrinaires = fils ; agglutination = grumeaux ronds).
    \item \textbf{Faux négatif (Rh) :} Variant \textbf{Du (D faible)} $\rightarrow$ Non réactif à l'anti-D standard (IgM) mais réactif au test indirect (Coombs) ou anti-D IgG. Au programme 3ème : on note Rh- si test direct négatif, mais on signale l'existence du Du.
    \item \textbf{Prozone :} Excès d'anticorps $\rightarrow$ pas d'agglutination visible. Rare avec réactifs calibrés.
\end{itemize}
\end{attention}

\coursec{Activités d'intégration et synthèse}

\courssub{Exercices résolus (type BEPC)}

\begin{exoresolu}[Exercice 1 : Génétique des groupes sanguins - Déduction parentale]
\textbf{Énoncé :} Un couple a trois enfants de groupes sanguins : Enfant 1 (A+), Enfant 2 (O-), Enfant 3 (B+). Le père est de groupe A+, la mère de groupe B+.
\begin{enumerate}
    \item Déterminer les génotypes probables des parents pour les systèmes ABO et Rhésus.
    \item Justifier la possibilité d'avoir un enfant O-.
    \item Quelle est la probabilité d'avoir un enfant de groupe AB+ ?
\end{enumerate}
\tcblower
\textbf{Solution détaillée :}
\begin{enumerate}
    \item \textbf{Père (A+) :} Phénotype A $\rightarrow$ Génotype ABO : $I^A I^A$ ou $I^A i$. Phénotype Rh+ $\rightarrow$ Génotype Rh : $D/D$ ou $D/d$.
          \textbf{Mère (B+) :} Phénotype B $\rightarrow$ Génotype ABO : $I^B I^B$ ou $I^B i$. Phénotype Rh+ $\rightarrow$ Génotype Rh : $D/D$ ou $D/d$.
    \item \textbf{Enfant O- ($ii$, $dd$) :} Pour être $ii$, l'enfant doit recevoir $i$ du père ET $i$ de la mère $\rightarrow$ Père \textbf{$I^A i$}, Mère \textbf{$I^B i$} (obligatoirement hétérozygotes). Pour être $dd$, l'enfant doit recevoir $d$ du père ET $d$ de la mère $\rightarrow$ Père \textbf{$D/d$}, Mère \textbf{$D/d$} (obligatoirement hétérozygotes).
          \textbf{Génotypes parentaux déduits :} Père $[I^A//i ; D//d]$ ; Mère $[I^B//i ; D//d]$.
    \item \textbf{Croisement ABO :} $I^A i \times I^B i$ $\rightarrow$ Gamètes paternels : $I^A$, $i$ ; Gamètes maternels : $I^B$, $i$.
          \begin{center}
          \begin{tabular}{|c|c|c|}\hline
          & $I^B$ & $i$ \\ \hline
          $I^A$ & $I^A I^B$ (AB) & $I^A i$ (A) \\ \hline
          $i$ & $I^B i$ (B) & $i i$ (O) \\ \hline
          \end{tabular}
          \end{center}
          Probabilité AB = 1/4.
          \textbf{Croisement Rhésus :} $D/d \times D/d$ $\rightarrow$ Probabilité Rh+ ($D/-$) = 3/4.
          \textbf{Indépendance des systèmes (chromosomes différents : 9 et 1) :} $P(AB+) = P(AB) \times P(Rh+) = \frac{1}{4} \times \frac{3}{4} = \frac{3}{16}$.
\end{enumerate}
\end{exoresolu}

\begin{exoresolu}[Exercice 2 : Transfusion et urgence - Cas pratique]
\textbf{Énoncé :} Au service d'urgence de l'Hôpital de District de Bertoua, un accidenté de la route (homme, 30 ans) arrive en choc hémorragique. Son groupe sanguin n'est pas connu. Le stock de sang disponible : 2 poches O-, 1 poche O+, 1 poche A-, 1 poche B+.
\begin{enumerate}
    \item Quelle(s) poche(s) peut-on transfuser \textbf{immédiatement} sans attendre le groupage (protocole d'urgence vitale) ? Justifier.
    \item Une fois le groupage effectué, l'accidenté est \textbf{AB+}. Quelles poches deviennent utilisables ?
    \item Pourquoi ne transfère-t-on généralement pas du sang total mais des \textbf{culots globulaires} (GR concentrés) ?
\end{enumerate}
\tcblower
\textbf{Solution détaillée :}
\begin{enumerate}
    \item \textbf{Urgence vitale (groupage inconnu) :} Règle = \textbf{Donneur universel de GR = O-}.
          Seules les \textbf{2 poches O-} sont transfusables sans risque immédiat (pas d'Ag A, B, D sur GR donneur $\rightarrow$ pas de réaction hémolytique aiguë par anti-A, anti-B, anti-D du receveur inconnu).
          \emph{Note :} O+ interdit (risque anti-D si receveur Rh-). A- interdit (risque anti-A si receveur B ou O). B+ interdit (double risque).
    \item \textbf{Receveur identifié AB+ (Receveur universel GR) :} N'a ni anti-A, ni anti-B, ni anti-D (sauf immunisation préexistante rare). Peut recevoir \textbf{TOUTES les poches restantes} : O+, A-, B+ (et O- déjà utilisées). Les GR du donneur (quel que soit leur antigène) ne seront pas agressés par les anticorps du receveur AB+.
    \item \textbf{Culots globulaires vs Sang total :}
          \begin{itemize}
              \item Volume réduit (hématocrate $\approx$ 60-70\% vs 40-45\%) $\rightarrow$ Moins de surcharge volémique (cœur).
              \item Plasma retiré $\rightarrow$ Élimine les anticorps du donneur (anti-A, anti-B) qui pourraient hémolyser les GR du receveur (surtout si donneur O donne à non-O). Élimine facteurs de coagulation instables, cytokines, potassium extracellulaire.
              \item Conservation plus longue (42 jours en CPD/SAGM vs 35 jours sang total).
          \end{itemize}
\end{enumerate}
\end{exoresolu}

\begin{exercice}[Exercice 3 : Analyse de caryotype et conseil génétique]
\textbf{Énoncé :} Un couple consulte en génétique médicale à l'Hôpital Gynéco-Obstétrique et Pédiatrique de Yaoundé (HGOPY) après deux fausses couches précoces. Le caryotype de la femme est normal (46,XX). Celui de l'homme révèle une \textbf{translocation robertsonienne équilibrée} entre les chromosomes 14 et 21 : \textbf{45,XY,der(14;21)(q10;q10)}.
\begin{enumerate}
    \item Expliquer pourquoi cet homme a 45 chromosomes au lieu de 46 mais un phénotype normal.
    \item Déterminer les types de gamètes (spermatozoïdes) que cet homme peut produire par ségrégation lors de la méiose (schématiser le trivalent 14-21).
    \item Calculer le risque théorique de trisomie 21 pour une grossesse à venir (non-disjonction adjacente vs alternée).
    \item Quel examen prénatal proposer au couple ? (Nommer la technique et le délai).
\end{enumerate}
\end{exercice}

\begin{corrige}[Exercice 3]
\textbf{1. Caryotype 45,XY,der(14;21) :} Fusion des bras longs de deux chromosomes acrocentriques (14 et 21) au niveau du centromère. Perte des bras courts (p) qui ne contiennent que des gènes d'ARNr (redondants, présents sur d'autres acrocentriques). $\rightarrow$ Pas de perte de matériel génétique essentiel $\rightarrow$ \textbf{Phénotype normal}. Mais nombre de chromosomes = 45 (1 chromosome fusionné remplace 2).
\textbf{2. Méiose :} Formation d'un \textbf{trivalent} (chromosome 14 normal + chromosome 21 normal + chromosome fusionné der(14;21)).
\textbf{Ségrégations possibles :}
\begin{itemize}
    \item \textbf{Alternée (équilibrée) :} 14 + 21 vont dans un pôle ; der(14;21) dans l'autre.
          $\rightarrow$ Gamète normal (23, 14+21) OU Gamète porteur translocation équilibrée (22, der(14;21)). $\rightarrow$ Enfant normal ou porteur sain comme le père.
    \item \textbf{Adjacente-1 (déséquilibrée) :} 14 + der(14;21) / 21 seul. $\rightarrow$ Trisomie 14 (létale) ou Monosomie 21 (létale).
    \item \textbf{Adjacente-2 (rare) :} 21 + der(14;21) / 14 seul. $\rightarrow$ Trisomie 21 (viable) ou Monosomie 14 (létale).
    \item \textbf{3:0 (tout dans un pôle) :} Létal.
\end{itemize}
\textbf{3. Risque Trisomie 21 :} Seule la ségrégation Adjacente-2 donnant (21 + der(14;21)) conduit à une trisomie 21 viable (46,XX,+21,der(14;21) ou 46,XY,+21,der(14;21)). Probabilité théorique $\approx$ 1/6 (si toutes ségrégations équiprobes) mais risque réel estimé \textbf{10-15\%} si le père est le porteur (moins de sélection contre spermatozoïdes anormaux que pour ovocytes). Risque plus élevé si la mère est porteuse.
\textbf{4. Diagnostic Prénatal (DPN) :} \textbf{Amniocentèse} (ponction de liquide amniotique, culture amniocytes, caryotype) vers \textbf{15-17 SA} (Semaines d'Aménorrhée). Ou \textbf{Prélèvement de Villosités Choriales (PVC)} vers \textbf{11-13 SA} (caryotype direct sur trophoblaste + culture). Confirmation par FISH (sonde chr 21) rapide en 24-48h. Conseil génétique obligatoire.
\end{corrige}

\courssub{Synthèse de l'unité : Expression de l'information génétique}

\begin{aretenir}[Bilan complet pour le BEPC]
\textbf{1. Support de l'information : Les Chromosomes}
\begin{itemize}
    \item Visibles en \textbf{métaphase} (mitose). Condensation chromatine (ADN + histones).
    \item \textbf{Espèce humaine : 2n = 46} (23 paires homologues). 22 paires \textbf{autosomes} + 1 paire \textbf{gonosomes} (XX / XY).
    \item \textbf{Caryotype} : Classement standard (taille, centromère, bandement G). Notation : 46,XX ou 46,XY.
    \item \textbf{Structure métaphase :} Chromosome double = 2 chromatides sœurs (ADN répliqué) + centromère + télomères.
    \item \textbf{Anomalies numériques :} Trisomie 21 (47,+21) par non-disjonction méiotique. Risque âge maternel.
\end{itemize}

\textbf{2. Unité d'information : Le Gène}
\begin{itemize}
    \item Portion d'ADN à un \textbf{locus} précis. Code une protéine (ou ARN fonctionnel).
    \item \textbf{Allèles} : Versions alternatives (même locus, séquence différente).
    \item \textbf{Génotype} (allèles présents) $\rightarrow$ \textbf{Phénotype} (caractère exprimé).
    \item Relations alléliques : Dominance / Récessivité / Codominance.
\end{itemize}

\textbf{3. Diversité humaine : Groupes sanguins (Exemple d'expression génique)}
\begin{itemize}
    \item \textbf{Système ABO (Chr 9) :} 3 allèles ($I^A, I^B, i$). Codominance $I^A/I^B$ ; Dominance sur $i$. 4 phénotypes (A, B, AB, O). Antigènes = sucres (glucides) sur GR. Anticorps naturels (anti-A, anti-B) dans le sérum (règle de Landsteiner).
    \item \textbf{Système Rhésus (Chr 1) :} Gène \emph{RHD}. Allèle $D$ (Ag D présent, dominant) / $d$ (délétion, pas Ag D, récessif). 2 phénotypes (Rh+, Rh-). Pas d'anticorps naturels (immunisation requise).
    \item \textbf{Transfusion :} Compatibilité ABO + Rhésus. \textbf{O- = Donneur universel GR} ; \textbf{AB+ = Receveur universel GR}. Culots globulaires préférés au sang total.
\end{itemize}

\textbf{4. Applications concrètes au Cameroun}
\begin{itemize}
    \item Dépistage néonatal drépanocytose (HbS $\rightarrow$ gène \emph{HBB} chr 11).
    \item Conseil génétique prénuptial (drépanocytose, thalassémie).
    \item Transfusion sécurisée (Centres de transfusion : Yaoundé, Douala, Garoua, etc.).
    \item Diagnostic prénatal (Amniocentèse, PVC) pour trisomies et translocations.
    \item Identification judiciaire / Tests de paternité (Marqueurs STR / ADN).
\end{itemize}
\end{aretenir}

\begin{methode}[Méthode type BEPC : Résoudre un problème de génétique formelle (ABO / Rhésus)]
\begin{enumerate}
    \item \textbf{Identifier les systèmes} concernés (ABO ? Rhésus ? Liés ou indépendants ? $\rightarrow$ Chr 9 et Chr 1 = Indépendants).
    \item \textbf{Noter les allèles} (symboles officiels : $I^A, I^B, i$ ; $D, d$).
    \item \textbf{Déduire les génotypes parentaux} à partir des phénotypes donnés + phénotypes des enfants (déduction rétrograde).
    \item \textbf{Construire les carrés de Punnett} (croissements) séparément pour chaque système.
    \item \textbf{Calculer les fréquences} génotypiques et phénotypiques pour chaque système.
    \item \textbf{Combiner les probabilités} (produit) si systèmes indépendants (loi de Mendel 2).
    \item \textbf{Répondre à la question} (probabilité, risque, conseil) en phrase complète, avec unités (\% ou fraction).
\end{enumerate}
\end{methode}

\coursec{TP N°3 : Conception d'une maquette de caryotype humain}

Dans la section précédente, nous avons découvert que le noyau de chaque cellule humaine contient \cle{46 chromosomes}, répartis en \cle{23 paires} : 22 paires de chromosomes ordinaires (les \cle{autosomes}) et une paire de \cle{chromosomes sexuels} (XX chez la femme, XY chez l'homme). Mais comment les scientifiques ont-ils établi ces faits ? Comment un médecin peut-il, à partir d'une simple photographie de chromosomes, déterminer le sexe d'un individu ou dépister une anomalie chromosomique ? C'est exactement ce que tu vas apprendre à faire dans ce travail pratique : construire toi-même un \cle{caryotype}, c'est-à-dire le portrait ordonné des chromosomes d'une cellule.

\courssub{Problème et hypothèses}

\begin{experience}[Situation-problème]
Au laboratoire d'analyses médicales de l'hôpital de district, un technicien observe au microscope une cellule humaine en \cle{métaphase} (stade de la division cellulaire où les chromosomes sont bien visibles, courts et épais). Il photographie les chromosomes, mais sur le cliché, ils sont \textbf{mélangés, désordonnés}.

\textbf{Problème posé :} comment organiser les chromosomes de cette cellule humaine pour :
\begin{itemize}
\item vérifier le nombre total de chromosomes ;
\item identifier le sexe chromosomique de l'individu ;
\item dépister une éventuelle anomalie (chromosome manquant ou surnuméraire) ?
\end{itemize}

\textbf{Hypothèses possibles :}
\begin{enumerate}
\item On peut ranger les chromosomes par taille, du plus grand au plus petit.
\item On peut regrouper les chromosomes deux par deux, car ils existent en paires.
\item On peut reconnaître le sexe en observant la dernière paire de chromosomes.
\end{enumerate}
\end{experience}

Pour trancher entre ces hypothèses, nous allons réaliser une maquette de caryotype à partir d'une photographie de chromosomes mélangés.

\courssub{Matériel et protocole expérimental}

\textbf{Matériel nécessaire :}
\begin{itemize}
\item une photographie (ou une planche fournie par le professeur) de chromosomes humains en métaphase, mélangés ;
\item une paire de ciseaux ;
\item de la colle ;
\item une feuille blanche (format A4) ;
\item un crayon, une règle ;
\item les critères de classement (longueur, position du centromère, bandes de coloration).
\end{itemize}

\begin{methode}[Protocole : les étapes de construction du caryotype]
\begin{enumerate}
\item \textbf{Observer} attentivement la photographie : chaque chromosome en métaphase a la forme d'un X (il est constitué de deux \cle{chromatides} réunies par le \cle{centromère}).
\item \textbf{Découper} soigneusement chaque chromosome en suivant son contour.
\item \textbf{Compter} les chromosomes découpés : on doit en obtenir 46.
\item \textbf{Apparier} les chromosomes : chercher pour chaque chromosome son \cle{homologue}, c'est-à-dire le chromosome qui lui ressemble (même longueur, même position du centromère).
\item \textbf{Classer} les paires par ordre de taille décroissante et les \textbf{numéroter de 1 à 22} (les autosomes).
\item \textbf{Identifier la paire sexuelle} : la 23\textsuperscript{e} paire. Si les deux chromosomes sont identiques et de taille moyenne, ce sont deux chromosomes X (individu de sexe féminin). Si l'un est grand (X) et l'autre très petit (Y), l'individu est de sexe masculin.
\item \textbf{Coller} les paires sur la feuille en respectant la grille de classement, puis vérifier le nombre total.
\end{enumerate}
\end{methode}

\begin{methode}[Critères d'appariement des chromosomes]
Pour reconnaître deux chromosomes homologues, on utilise trois critères :
\begin{enumerate}
\item \textbf{La longueur} : deux chromosomes homologues ont la même taille. La paire 1 est la plus longue, la paire 22 l'une des plus petites.
\item \textbf{La position du centromère} : le centromère divise le chromosome en deux bras. Il peut être au milieu (bras égaux), décalé (un bras long, un bras court) ou presque au bout. Deux homologues ont le centromère à la même place.
\item \textbf{Les bandes de coloration} : après coloration spéciale, chaque chromosome présente une alternance de bandes claires et foncées, comme un code-barres. Deux chromosomes homologues ont le même code-barres.
\end{enumerate}
\end{methode}

\begin{attention}[Piège fréquent : la paire sexuelle]
Ne cherche pas à apparier le chromosome X avec le chromosome Y chez un garçon : ils sont de tailles très différentes ! La paire sexuelle est la seule paire où les deux chromosomes peuvent être \textbf{différents}. C'est justement cette différence qui permet d'identifier le sexe masculin (XY). De plus, vérifie \textbf{toujours le nombre total de chromosomes (46) avant de conclure} à une anomalie : un chromosome mal découpé ou collé deux fois peut fausser le diagnostic.
\end{attention}

\courssub{Schéma du matériel de départ : les chromosomes mélangés}

\begin{popfigure}
\begin{center}
\begin{tikzpicture}[scale=0.85]
% fond photo
\fill[popcream] (-0.5,-0.5) rectangle (10.5,6.5);
\draw[popline, thick] (-0.5,-0.5) rectangle (10.5,6.5);
% chromosomes mélangés (forme en X simplifiée) : x / y / rotation / echelle
\foreach \x/\y/\r/\s in {1.2/5.4/20/1.1, 2.8/5.6/-35/0.8, 4.5/5.3/70/1.0, 6.2/5.6/-15/0.7, 8.1/5.3/40/0.9, 9.6/5.6/-60/0.75, 1.0/3.8/-50/0.85, 2.5/3.6/15/1.05, 4.2/3.9/-80/0.7, 5.9/3.6/55/0.95, 7.6/3.9/-25/0.8, 9.4/3.6/10/1.0, 1.4/2.1/65/0.75, 3.1/2.3/-10/0.9, 4.8/2.0/30/0.8, 6.5/2.3/-55/1.0, 8.3/2.0/85/0.7, 9.8/2.2/-30/0.85, 1.8/0.7/-70/0.9, 3.6/0.8/25/0.75, 5.4/0.6/-40/0.95, 7.2/0.8/60/0.8, 8.9/0.6/-20/0.7}{
\begin{scope}[shift={(\x,\y)}, rotate=\r, scale=\s]
\draw[line width=2.2pt, popblue, rounded corners=2pt] (-0.12,-0.5) -- (0.05,0) -- (-0.12,0.5);
\draw[line width=2.2pt, popblue, rounded corners=2pt] (0.12,-0.5) -- (-0.05,0) -- (0.12,0.5);
\end{scope}}
% chromosome Y (petit, orange)
\begin{scope}[shift={(9.9,0.8)}, rotate=45, scale=0.4]
\draw[line width=2.2pt, poporange, rounded corners=2pt] (-0.12,-0.5) -- (0.05,0) -- (-0.12,0.5);
\draw[line width=2.2pt, poporange, rounded corners=2pt] (0.12,-0.5) -- (-0.05,0) -- (0.12,0.5);
\end{scope}
\end{tikzpicture}
\end{center}
\legende{Photographie-schéma d'une cellule humaine en métaphase : les chromosomes sont mélangés et désordonnés (seuls quelques-uns sont représentés ici pour la clarté). Chaque chromosome a une forme de X (deux chromatides jointes par le centromère). En bleu : autosomes et chromosome X ; en orange : le petit chromosome Y. Le travail de l'élève consiste à découper puis à classer tous ces chromosomes.}
\end{popfigure}

\courssub{Grille de classement et résultats attendus}

Après découpage et appariement, les chromosomes sont collés dans une grille ordonnée : les paires d'autosomes sont numérotées de 1 à 22 (de la plus grande à la plus petite), et la paire sexuelle occupe la dernière case.

\begin{popfigure}
\begin{center}
\begin{tikzpicture}[scale=1.0]
% ligne 1 : paires 1 à 5
\foreach \i/\n in {0/1, 1/2, 2/3, 3/4, 4/5}{
\draw[popline, thick] (\i*2.1,6) rectangle (\i*2.1+1.9,7.4);
\node[popdark, font=\small\bfseries] at (\i*2.1+0.95,7.15) {Paire \n};
\draw[line width=2pt, popblue, rounded corners=1.5pt] (\i*2.1+0.55,6.25) -- (\i*2.1+0.75,6.7) -- (\i*2.1+0.55,6.95);
\draw[line width=2pt, popblue, rounded corners=1.5pt] (\i*2.1+0.85,6.25) -- (\i*2.1+0.65,6.7) -- (\i*2.1+0.85,6.95);
\draw[line width=2pt, popblue, rounded corners=1.5pt] (\i*2.1+1.15,6.25) -- (\i*2.1+1.35,6.7) -- (\i*2.1+1.15,6.95);
\draw[line width=2pt, popblue, rounded corners=1.5pt] (\i*2.1+1.45,6.25) -- (\i*2.1+1.25,6.7) -- (\i*2.1+1.45,6.95);
}
% ligne 2 : paires 6 à 12
\foreach \i/\n in {0/6, 1/7, 2/8, 3/9, 4/10, 5/11, 6/12}{
\draw[popline, thick] (\i*1.5,4.4) rectangle (\i*1.5+1.35,5.7);
\node[popdark, font=\footnotesize\bfseries] at (\i*1.5+0.675,5.5) {\n};
\draw[line width=1.6pt, popblue, rounded corners=1pt] (\i*1.5+0.35,4.6) -- (\i*1.5+0.5,4.95) -- (\i*1.5+0.35,5.2);
\draw[line width=1.6pt, popblue, rounded corners=1pt] (\i*1.5+0.6,4.6) -- (\i*1.5+0.45,4.95) -- (\i*1.5+0.6,5.2);
\draw[line width=1.6pt, popblue, rounded corners=1pt] (\i*1.5+0.85,4.6) -- (\i*1.5+1.0,4.95) -- (\i*1.5+0.85,5.2);
\draw[line width=1.6pt, popblue, rounded corners=1pt] (\i*1.5+1.1,4.6) -- (\i*1.5+0.95,4.95) -- (\i*1.5+1.1,5.2);
}
% ligne 3 : paires 13 à 18
\foreach \i/\n in {0/13, 1/14, 2/15, 3/16, 4/17, 5/18}{
\draw[popline, thick] (\i*1.75,2.9) rectangle (\i*1.75+1.6,4.1);
\node[popdark, font=\footnotesize\bfseries] at (\i*1.75+0.8,3.92) {\n};
\draw[line width=1.5pt, popblue, rounded corners=1pt] (\i*1.75+0.45,3.05) -- (\i*1.75+0.6,3.35) -- (\i*1.75+0.45,3.6);
\draw[line width=1.5pt, popblue, rounded corners=1pt] (\i*1.75+0.7,3.05) -- (\i*1.75+0.55,3.35) -- (\i*1.75+0.7,3.6);
\draw[line width=1.5pt, popblue, rounded corners=1pt] (\i*1.75+1.0,3.05) -- (\i*1.75+1.15,3.35) -- (\i*1.75+1.0,3.6);
\draw[line width=1.5pt, popblue, rounded corners=1pt] (\i*1.75+1.25,3.05) -- (\i*1.75+1.1,3.35) -- (\i*1.75+1.25,3.6);
}
% ligne 4 : paires 19 à 22 + case sexuelle
\foreach \i/\n in {0/19, 1/20, 2/21, 3/22}{
\draw[popline, thick] (\i*1.75,1.4) rectangle (\i*1.75+1.6,2.6);
\node[popdark, font=\footnotesize\bfseries] at (\i*1.75+0.8,2.42) {\n};
\draw[line width=1.4pt, popblue, rounded corners=1pt] (\i*1.75+0.5,1.55) -- (\i*1.75+0.62,1.8) -- (\i*1.75+0.5,2.05);
\draw[line width=1.4pt, popblue, rounded corners=1pt] (\i*1.75+0.72,1.55) -- (\i*1.75+0.6,1.8) -- (\i*1.75+0.72,2.05);
\draw[line width=1.4pt, popblue, rounded corners=1pt] (\i*1.75+0.95,1.55) -- (\i*1.75+1.07,1.8) -- (\i*1.75+0.95,2.05);
\draw[line width=1.4pt, popblue, rounded corners=1pt] (\i*1.75+1.17,1.55) -- (\i*1.75+1.05,1.8) -- (\i*1.75+1.17,2.05);
}
% case sexuelle
\draw[poporange, thick] (7.0,1.4) rectangle (9.6,2.6);
\node[popdark, font=\footnotesize\bfseries] at (8.3,2.42) {Paire sexuelle};
% X
\draw[line width=1.6pt, popblue, rounded corners=1pt] (7.6,1.55) -- (7.75,1.8) -- (7.6,2.05);
\draw[line width=1.6pt, popblue, rounded corners=1pt] (7.9,1.55) -- (7.75,1.8) -- (7.9,2.05);
\node[popblue, font=\footnotesize] at (7.75,1.48) {X};
% Y petit
\draw[line width=1.4pt, poporange, rounded corners=1pt] (8.7,1.6) -- (8.8,1.75) -- (8.7,1.9);
\draw[line width=1.4pt, poporange, rounded corners=1pt] (8.9,1.6) -- (8.8,1.75) -- (8.9,1.9);
\node[poporange, font=\footnotesize] at (8.8,1.5) {Y};
% annotations
\node[popmuted, font=\footnotesize, align=center] at (5.25,0.7) {22 paires d'autosomes (n\textsuperscript{o}1 à 22, taille décroissante) + 1 paire de chromosomes sexuels.};
\node[popmuted, font=\footnotesize, align=center] at (5.25,0.25) {Ici : XY $\Rightarrow$ individu de sexe masculin. Total : 46 chromosomes.};
\end{tikzpicture}
\end{center}
\legende{Grille de classement du caryotype humain. Les chromosomes sont collés par paires d'homologues, numérotées de 1 à 22 (autosomes, classés par taille décroissante), la 23\textsuperscript{e} case recevant la paire sexuelle. Dans l'exemple représenté, la paire sexuelle est XY : le caryotype s'écrit 46, XY (homme).}
\end{popfigure}

\courssub{Exploitation des résultats}

Une fois la maquette terminée, l'exploitation des résultats se fait en trois temps :

\begin{enumerate}
\item \textbf{Vérification du nombre total :} on compte les chromosomes collés. Un caryotype humain normal comporte \textbf{46 chromosomes}, soit 23 paires complètes.
\item \textbf{Détermination du sexe chromosomique :} on observe la 23\textsuperscript{e} paire. Deux chromosomes X identiques $\Rightarrow$ individu de sexe féminin (caryotype noté \cle{46, XX}). Un chromosome X et un petit chromosome Y $\Rightarrow$ individu de sexe masculin (caryotype noté \cle{46, XY}).
\item \textbf{Recherche d'anomalie :} on vérifie que chaque paire comporte exactement deux chromosomes. Une paire à trois chromosomes (ou à un seul) signale une anomalie du nombre de chromosomes.
\end{enumerate}

\begin{exemplebox}[Lecture de caryotypes]
\begin{itemize}
\item \textbf{46, XY} : 46 chromosomes, paire sexuelle XY $\Rightarrow$ homme sans anomalie chromosomique.
\item \textbf{46, XX} : 46 chromosomes, paire sexuelle XX $\Rightarrow$ femme sans anomalie chromosomique.
\item \textbf{47, XX, +21} : 47 chromosomes, paire sexuelle XX, avec un chromosome 21 \textbf{surnuméraire} (trois chromosomes 21 au lieu de deux) $\Rightarrow$ femme atteinte de \cle{trisomie 21}.
\end{itemize}
La formule du caryotype se lit donc toujours dans l'ordre : nombre total, paire sexuelle, anomalie éventuelle.
\end{exemplebox}

\courssub{Comparaison avec un caryotype de trisomie 21}

Le professeur distribue maintenant une seconde photographie de chromosomes, provenant cette fois d'un enfant atteint de \cle{trisomie 21}. En appliquant exactement la même démarche (découpage, appariement, classement, comptage), on obtient les résultats suivants :

\begin{center}
\begin{tabularx}{\linewidth}{|l|X|X|}
\hline
\rowcolor{popblueL} \textbf{Critère observé} & \textbf{Caryotype normal} & \textbf{Caryotype trisomie 21} \\
\hline
Nombre total de chromosomes & 46 & 47 \\
\hline
Paire 21 & 2 chromosomes & 3 chromosomes (un chromosome surnuméraire) \\
\hline
Paire sexuelle & XX ou XY & XX ou XY (non modifiée) \\
\hline
Formule du caryotype & 46, XX ou 46, XY & 47, XX, +21 ou 47, XY, +21 \\
\hline
\end{tabularx}
\end{center}

Le repérage du chromosome surnuméraire est facile sur la maquette : la case de la paire 21 contient \textbf{trois} chromosomes au lieu de deux, et le comptage final donne 47 au lieu de 46. Cette anomalie, présente dès la naissance, est à l'origine des signes de la trisomie 21.

\begin{attention}[Toujours compter avant de conclure]
Une erreur de découpage (chromosome coupé en deux) ou de collage (chromosome perdu) peut faire croire à tort à une anomalie. La règle d'or du diagnostic : \textbf{on ne conclut à une anomalie que si le comptage est fiable et vérifié deux fois}. Au laboratoire, le technicien analyse toujours plusieurs cellules avant de valider un caryotype.
\end{attention}

\begin{savaistu}[L'amniocentèse : le caryotype avant la naissance]
Il est possible d'établir le caryotype d'un bébé \textbf{avant sa naissance} : c'est le \cle{diagnostic prénatal}. Lors de l'\cle{amniocentèse}, le médecin prélève, sous contrôle de l'échographie, un peu de liquide amniotique (le liquide qui entoure le fœtus dans le ventre de la mère). Ce liquide contient des cellules du fœtus. Mises en culture puis photographiées en métaphase, ces cellules permettent de réaliser le caryotype du bébé : on connaît alors son sexe chromosomique et on peut dépister une anomalie comme la trisomie 21. Cet examen est proposé dans les maternités équipées, notamment lorsque la mère est âgée de plus de 35 ans, car le risque de trisomie 21 augmente avec l'âge de la mère.
\end{savaistu}

\courssub{Rédaction du compte rendu de TP}

Le compte rendu est la trace écrite de ta démarche scientifique. Il doit suivre le plan de la démarche d'investigation.

\begin{exoresolu}[Compte rendu type du TP N°3]
Rédige le compte rendu complet de ce TP à partir de la maquette réalisée (photographie d'une cellule humaine en métaphase).
\tcblower
\textbf{1. Problème :} comment organiser les chromosomes d'une cellule humaine pour identifier le sexe de l'individu et dépister une anomalie chromosomique ?

\textbf{2. Hypothèse :} les chromosomes existent par paires d'homologues ; on peut donc les apparier par taille et position du centromère, puis reconnaître le sexe grâce à la dernière paire.

\textbf{3. Démarche (protocole) :} nous avons découpé les 46 chromosomes de la photographie, apparié les homologues selon trois critères (longueur, position du centromère, bandes de coloration), numéroté les paires de 1 à 22 par taille décroissante, puis identifié la paire sexuelle avant de coller l'ensemble sur la grille.

\textbf{4. Résultats :}
\begin{itemize}
\item nombre total de chromosomes : 46 (vérifié deux fois) ;
\item 22 paires d'autosomes complètes, sans chromosome surnuméraire ;
\item paire sexuelle : un chromosome X et un chromosome Y (petit).
\end{itemize}

\textbf{5. Interprétation :} la présence de la paire XY indique un individu de sexe masculin. Le total de 46 chromosomes, avec toutes les paires complètes, indique l'absence d'anomalie du nombre de chromosomes.

\textbf{6. Conclusion :} le caryotype de l'individu étudié s'écrit \textbf{46, XY} : c'est un homme sans anomalie chromosomique. L'hypothèse est validée : le classement par paires d'homologues permet bien d'identifier le sexe et de dépister les anomalies. (Si la case 21 avait contenu trois chromosomes, le caryotype aurait été 47, XY, +21 : garçon atteint de trisomie 21.)
\end{exoresolu}

\begin{exercice}[Pour t'entraîner]
Sur une maquette de caryotype, un élève compte 45 chromosomes et observe que la paire 13 ne contient qu'un seul chromosome.
\begin{enumerate}
\item Quelle vérification l'élève doit-il faire avant de conclure à une anomalie ?
\item Si le résultat est confirmé, écris la formule de ce caryotype sachant que la paire sexuelle est XX.
\end{enumerate}
\end{exercice}

\begin{corrige}[Exercice]
\begin{enumerate}
\item Avant de conclure, l'élève doit vérifier qu'aucun chromosome n'a été perdu lors du découpage ou du collage, et refaire le comptage (au laboratoire, on analyserait plusieurs cellules). On ne conclut à une anomalie que sur un comptage fiable.
\item Si l'anomalie est confirmée, il manque un chromosome 13 : le total est 45 et la paire sexuelle est XX. Le caryotype s'écrit \textbf{45, XX, $-$13}.
\end{enumerate}
\end{corrige}

\begin{aretenir}[L'essentiel du TP N°3]
\begin{itemize}
\item Le \cle{caryotype} est le portrait ordonné des chromosomes d'une cellule, réalisé à partir d'une photographie en métaphase.
\item On apparie les chromosomes homologues grâce à trois critères : \textbf{longueur}, \textbf{position du centromère}, \textbf{bandes de coloration}.
\item Un caryotype humain normal comporte \textbf{46 chromosomes} : 22 paires d'autosomes (n\textsuperscript{o}1 à 22) et 1 paire sexuelle (XX = femme, XY = homme).
\item La formule du caryotype s'écrit : nombre total, paire sexuelle, anomalie éventuelle (ex. : 47, XX, +21).
\item Toujours vérifier le comptage avant de conclure à une anomalie.
\end{itemize}
\end{aretenir}

Maintenant que nous savons lire et construire un caryotype, une question se pose : que portent exactement ces chromosomes ? Chacun d'eux renferme de l'information génétique sous forme de \cle{gènes} : c'est ce que nous étudierons dans la section suivante, « Les gènes humains ».

\textbf{Observation : la nature chimique du chromosome}\par
Depuis les travaux d'Avery, MacLeod et McCarty (1944), puis de Hershey et Chase (1952), on sait que l'acide désoxyribonucléique (\textbf{ADN}) est le support de l'information héréditaire. Au sein du noyau, l'ADN n'est pas nu : il est enroulé autour de protéines appelées \textbf{histones}, formant une structure appelée \textbf{chromatine}. Au moment de la division cellulaire, la chromatine se condense fortement pour former les chromosomes visibles.

\begin{popfigure}
\begin{tikzpicture}[scale=0.9, every node/.style={font=\small}]
    % Couleurs
    \colorlet{cellcolor}{popblue!10}
    \colorlet{nuccolor}{popblue!30}
    \colorlet{chromcolor}{poporange!60}
    \colorlet{dnacolor}{poppink!70}
    \colorlet{genecolor}{popgreen!70}
    
    % 1. Cellule
    \begin{scope}[xshift=0cm, yshift=0cm]
        \draw[popblue, thick, fill=cellcolor] (0,0) ellipse (2.5 and 1.8);
        \draw[popblue, thick, fill=nuccolor] (0,0.2) ellipse (1.2 and 0.8);
        \foreach \i in {1,...,6} {
            \pgfmathsetmacro{\ang}{rand*360}
            \pgfmathsetmacro{\r}{0.3+rand*0.7}
            \draw[chromcolor, thick] ({0.2*cos(\ang)+\r*cos(\ang)},{0.2*sin(\ang)+\r*sin(\ang)}) -- ++({0.3*cos(\ang+90)},{0.3*sin(\ang+90)});
        }
        \node[popdark, align=center] at (0,-2.2) {\textbf{1. Cellule eucaryote}\\(~10--30 \textmu m)};
        \draw[->, thick, popdark] (3,0) -- (4,0);
    \end{scope}
    
    % 2. Noyau
    \begin{scope}[xshift=5.5cm, yshift=0cm]
        \draw[popblue, thick, fill=nuccolor] (0,0) ellipse (2 and 1.5);
        \foreach \i in {1,...,12} {
            \pgfmathsetmacro{\ang}{rand*360}
            \pgfmathsetmacro{\r}{0.2+rand*1.2}
            \draw[chromcolor, thick, line width=1.5pt] ({0.1*cos(\ang)+\r*cos(\ang)},{0.1*sin(\ang)+\r*sin(\ang)}) -- ++({0.5*cos(\ang+90)},{0.5*sin(\ang+90)});
        }
        \node[popdark, align=center] at (0,-2.2) {\textbf{2. Noyau}\\(~5--10 \textmu m)};
        \draw[->, thick, popdark] (3,0) -- (4,0);
    \end{scope}
    
    % 3. Chromosome (condensé)
    \begin{scope}[xshift=11cm, yshift=0cm]
        \draw[chromcolor, thick, fill=chromcolor!20, rounded corners=2pt] (-0.3,-1.2) -- (-0.3,1.2) .. controls (-0.3,1.5) and (0.3,1.5) .. (0.3,1.2) -- (0.3,-1.2) .. controls (0.3,-1.5) and (-0.3,-1.5) .. (-0.3,-1.2);
        \draw[popdark, thick] (0,0) ellipse (0.15 and 0.05); % centromère
        \node[popdark, align=center] at (0,-2.2) {\textbf{3. Chromosome (mitose)}\\(~1--10 \textmu m)};
        \draw[->, thick, popdark] (3,0) -- (4,0);
    \end{scope}
    
    % 4. Fibre de chromatine / ADN
    \begin{scope}[xshift=16.5cm, yshift=0cm]
        \draw[dnacolor, thick, decorate, decoration={coil, aspect=0.3, segment length=3mm, amplitude=1.5mm}] (-1.5,0) -- (1.5,0);
        \draw[popdark, thick] (-1.5,0.3) -- (-1.5,-0.3);
        \draw[popdark, thick] (1.5,0.3) -- (1.5,-0.3);
        \node[popdark] at (-1.5,0.6) {Histones};
        \node[popdark] at (1.5,0.6) {Histones};
        \foreach \x in {-1.2,-0.6,0,0.6,1.2} {
            \draw[popdark, fill=popdark] (\x,0) circle (0.1);
        }
        \node[popdark, align=center] at (0,-2.2) {\textbf{4. Fibre de chromatine (ADN + histones)}\\(~30 nm)};
        \draw[->, thick, popdark] (3,0) -- (4,0);
    \end{scope}
    
    % 5. Double hélice ADN -> Gène
    \begin{scope}[xshift=22cm, yshift=0cm]
        % Double hélice simplifiée
        \foreach \y in {-1.5,-1.2,...,1.5} {
            \pgfmathsetmacro{\xoff}{0.3*sin(\y*180)}
            \draw[dnacolor, thick] (\xoff-0.3,\y) -- (\xoff+0.3,\y);
            \draw[popdark, very thin] (\xoff-0.3,\y) -- (\xoff+0.3,\y);
        }
        % Brins
        \draw[dnacolor, thick] (-0.3,-1.5) .. controls (-0.3,0) .. (-0.3,1.5);
        \draw[dnacolor, thick] (0.3,-1.5) .. controls (0.3,0) .. (0.3,1.5);
        % Gène surligné
        \draw[genecolor, line width=4pt, opacity=0.6] (-0.3,0.2) -- (-0.3,-0.8);
        \draw[genecolor, line width=4pt, opacity=0.6] (0.3,0.2) -- (0.3,-0.8);
        \node[genecolor, font=\bfseries] at (0,1.8) {GÈNE};
        \node[popdark, align=center] at (0,-2.2) {\textbf{5. Molécule d'ADN (double hélice)}\\(~2 nm de large)};
    \end{scope}
\end{tikzpicture}
\legende{Zoom successif : de la cellule au gène. L'information génétique est portée par la séquence des nucléotides le long de la molécule d'ADN. Un gène est un segment précis de cette molécule.}
\end{popfigure}

\begin{definition}[Gène]
Un \textbf{gène} est un \textbf{segment d'ADN} (une portion de la molécule d'ADN d'un chromosome) qui porte l'information nécessaire à la synthèse d'un produit fonctionnel (généralement une \textbf{protéine}, parfois un ARN fonctionnel) et qui détermine un \textbf{caractère héréditaire} (ex. : groupe sanguin ABO, couleur des yeux, type d'hémoglobine).
\end{definition}

\begin{aretenir}[Du chromosome au gène]
\begin{itemize}
    \item Le chromosome est constitué \textbf{d'une molécule d'ADN linéaire} (chez l'eucaryote) \textbf{associée à des protéines} (histones).
    \item L'ADN est une double hélice dont les « marches » sont formées de paires de bases azotées (A-T, C-G). La \textbf{séquence de ces bases} code l'information.
    \item Un \textbf{gène} est une \textbf{portion définie} de cette séquence. Le génome humain (ensemble de l'ADN) contient environ \textbf{20 000 à 25 000 gènes} (résultat du Projet Génome Humain, achevé en 2003).
\end{itemize}
\end{aretenir}

\courssub{Localisation des gènes : locus et allèles}

Sur votre maquette de caryotype, vous avez associé les chromosomes deux par deux : ce sont les \textbf{chromosomes homologues} (même taille, même forme, même centromère, même bandeau de coloration). Ils portent les \textbf{mêmes gènes} aux \textbf{mêmes positions}.

\begin{definition}[Locus]
Le \textbf{locus} (pluriel : loci) est la \textbf{position précise} qu'occupe un gène sur un chromosome. Les deux chromosomes homologues portent le même gène \textbf{au même locus}.
\end{definition}

\begin{definition}[Allèle]
Un \textbf{allèle} est une \textbf{version différente} d'un même gène. Les allèles d'un gène occupent le même locus sur les chromosomes homologues. Ils diffèrent par leur \textbf{séquence de nucléotides} (mutations), ce qui peut modifier la protéine produite et donc le caractère observable.
\end{definition}

\begin{popfigure}
\begin{tikzpicture}[scale=1.1, every node/.style={font=\small}]
    % Chromosome homologue 1 (paternel)
    \draw[popblue, thick, fill=popblue!10, rounded corners=1pt] (0,1.5) -- (0,0.5) .. controls (0,0) and (1,0) .. (1,0.5) -- (1,1.5) .. controls (1,2) and (0,2) .. (0,1.5);
    \draw[popdark, thick] (0.5,1) ellipse (0.15 and 0.05); % centromère
    % Locus sur chr 1
    \draw[popgreen, line width=3pt] (0.5,1.35) -- (0.5,1.2);
    \node[popgreen, right, font=\bfseries] at (0.55,1.27) {Allèle $A$};
    
    % Chromosome homologue 2 (maternel)
    \draw[poporange, thick, fill=poporange!10, rounded corners=1pt] (3,1.5) -- (3,0.5) .. controls (3,0) and (4,0) .. (4,0.5) -- (4,1.5) .. controls (4,2) and (3,2) .. (3,1.5);
    \draw[popdark, thick] (3.5,1) ellipse (0.15 and 0.05); % centromère
    % Locus sur chr 2
    \draw[poppink, line width=3pt] (3.5,1.35) -- (3.5,1.2);
    \node[poppink, right, font=\bfseries] at (3.55,1.27) {Allèle $a$};
    
    % Flèches locus
    \draw[<->, popdark, thick] (0.5,1.6) -- (3.5,1.6) node[midway, above] {\textbf{Même locus}};
    
    % Légendes
    \node[popblue, left, font=\bfseries] at (-0.5,1) {Chr. paternel};
    \node[poporange, right, font=\bfseries] at (4.5,1) {Chr. maternel};
    
    % Zoom ADN au niveau du locus
    \begin{scope}[xshift=7cm, yshift=0cm]
        \node[align=center, popdark] at (0,1.5) {\textbf{Zoom au niveau du locus}};
        % Allèle A
        \draw[popgreen, thick] (-1.5,0.5) -- (-1.5,-0.5);
        \draw[popgreen, thick] (-1.1,0.5) -- (-1.1,-0.5);
        \foreach \y in {0.4,0.1,-0.2} {
            \draw[popgreen, thick] (-1.5,\y) -- (-1.1,\y);
        }
        \node[popgreen, left] at (-1.6,0) {Séquence : ...ATG \textbf{GCT} TAC...};
        \node[popgreen, below, font=\bfseries] at (-1.3,-0.8) {Allèle $A$};
        
        % Allèle a
        \draw[poppink, thick] (1.1,0.5) -- (1.1,-0.5);
        \draw[poppink, thick] (1.5,0.5) -- (1.5,-0.5);
        \foreach \y in {0.4,0.1,-0.2} {
            \draw[poppink, thick] (1.1,\y) -- (1.5,\y);
        }
        \node[poppink, right] at (1.6,0) {Séquence : ...ATG \textbf{GTT} TAC...};
        \node[poppink, below, font=\bfseries] at (1.3,-0.8) {Allèle $a$ (mutation)};
        
        \draw[<->, popmuted] (-1.3,-0.5) -- (1.3,-0.5) node[midway, below] {Différence de séquence (mutation)};
    \end{scope}
\end{tikzpicture}
\legende{Paire de chromosomes homologues. Le gène pour un caractère donné (ex. groupe sanguin) se trouve au même locus sur les deux homologues. Les deux versions (allèles) peuvent être identiques ou différentes (ici $A$ et $a$).}
\end{popfigure}

\textbf{Notation des génotypes}\par
On note le génotype (les allèles portés) entre crochets ou avec une barre double.
\begin{itemize}
    \item \textbf{Homozygote} : les deux allèles sont \textbf{identiques}. Ex. : $[A//A]$ ou $[a//a]$.
    \item \textbf{Hétérozygote} : les deux allèles sont \textbf{différents}. Ex. : $[A//a]$.
\end{itemize}

\begin{definition}[Génotype et Phénotype]
\begin{itemize}
    \item Le \textbf{génotype} est l'\textbf{ensemble des allèles} portés par un individu pour un ou plusieurs gènes (ex. : $[A//a]$).
    \item Le \textbf{phénotype} est l'\textbf{expression observable} du caractère (ex. : « Groupe sanguin A »).
\end{itemize}
\end{definition}

\courssub{Relations entre allèles : dominance et récessivité}

Tous les allèles ne s'expriment pas de la même manière chez l'hétérozygote. C'est la notion de \textbf{dominance}.

\begin{experience}[Mise en évidence de la dominance : le système ABO (rappel historique)]
\textbf{Observation :} Au début du XX\textsuperscript{e} siècle, Karl Landsteiner découvre que le sang de certains individus agglutine les globules rouges d'autres. Il identifie 4 groupes : A, B, AB, O.
\textbf{Interprétation génétique (Bernstein, 1925) :} Un seul gène ($I$), trois allèles : $I^A$, $I^B$, $i$ (ou $I^O$).
\begin{itemize}
    \item $I^A$ et $I^B$ sont \textbf{codominants} : chez l'hétérozygote $[I^A//I^B]$, les deux s'expriment $\rightarrow$ Phénotype AB.
    \item $I^A$ et $I^B$ sont \textbf{dominants} sur $i$ : $[I^A//i]$ donne phénotype A ; $[I^B//i]$ donne phénotype B.
    \item $i$ est \textbf{récessif} : il ne s'exprime qu'à l'état \textbf{homozygote} $[i//i]$ $\rightarrow$ Phénotype O.
\end{itemize}
\end{experience}

\begin{propriete}[Dominance et récessivité]
\begin{itemize}
    \item Un \textbf{allèle dominant} s'exprime au niveau du phénotype \textbf{tant qu'il est présent} (états homozygote $[D//D]$ et hétérozygote $[D//r]$).
    \item Un \textbf{allèle récessif} ne s'exprime au niveau du phénotype \textbf{que s'il est présent en double dose} (état homozygote $[r//r]$). Il est « masqué » par l'allèle dominant chez l'hétérozygote.
\end{itemize}
\end{propriete}

\begin{attention}[Pièges classiques sur la dominance]
\begin{itemize}
    \item \textbf{« Dominant = fréquent » : FAUX.} L'allèle $i$ (récessif, groupe O) est le \textbf{plus fréquent} au Cameroun et dans le monde (environ 60 à 70\% de fréquence allélique), alors qu'il est récessif.
    \item \textbf{« Dominant = meilleur / normal » : FAUX.} La dominance est une relation d'expression entre deux allèles, pas une valeur qualitative. L'allèle de la drépanocytose ($S$) est codominant avec l'allèle normal ($A$) pour le phénotype hémoglobine, mais récessif pour la maladie.
    \item \textbf{Notation :} Utilisez des lettres \emph{italiques} pour les allèles ($A$, $a$, $I^A$, $i$). Le génotype s'écrit $[A//a]$ ou $A/a$. Ne confondez pas le gène ($I$) et l'allèle ($I^A$).
\end{itemize}
\end{attention}

\begin{exemplebox}[Système ABO : correspondance Génotype $\rightarrow$ Phénotype]
\begin{center}
\begin{tabularx}{\linewidth}{|c|X|X|}\hline
\textbf{Génotype} & \textbf{Allèles présents} & \textbf{Phénotype (Groupe sanguin)} \\ \hline
$[I^A//I^A]$ & $I^A$ homozygote & A \\ \hline
$[I^A//i]$ & $I^A$ dominant sur $i$ & A \\ \hline
$[I^B//I^B]$ & $I^B$ homozygote & B \\ \hline
$[I^B//i]$ & $I^B$ dominant sur $i$ & B \\ \hline
$[I^A//I^B]$ & $I^A$ et $I^B$ codominants & AB \\ \hline
$[i//i]$ & $i$ récessif homozygote & O \\ \hline
\end{tabularx}
\end{center}
\emph{Remarque :} Les génotypes $[I^A//I^A]$ et $[I^A//i]$ donnent le \textbf{même phénotype A}. On ne peut pas distinguer un homozygote d'un hétérozygote par le seul phénotype (sauf étude de descendance ou génotypage moléculaire).
\end{exemplebox}

\courssub{De l'allèle au caractère : l'expression du gène}

\begin{propriete}[De l'allèle au caractère]
Un gène (segment d'ADN) contient les instructions pour fabriquer une \textbf{protéine}. 
La séquence des bases de l'ADN détermine la séquence des acides aminés de la protéine. 
La structure de la protéine détermine sa fonction, ce qui produit le \textbf{phénotype} (caractère observable).
\textit{Exemple :} Les allèles $I^A$ et $I^B$ codent pour des enzymes fonctionnelles (glycosyltransférases) qui fabriquent les antigènes A et B à la surface des globules rouges. L'allèle $i$ code pour une enzyme non fonctionnelle $\rightarrow$ absence d'antigènes A et B (groupe O).
\end{propriete}

\begin{savaistu}[Diversité humaine et similarité génomique]
Deux êtres humains pris au hasard partagent \textbf{environ 99,9\% de leur séquence d'ADN}. Toute notre diversité (taille, couleur de peau, groupes sanguins, susceptibilité aux maladies) repose sur les différences présentes dans les \textbf{0,1\% restants} (polymorphismes mononucléotidiques - SNP, variations de nombre de copies, etc.). C'est dire si la précision de l'expression des gènes est cruciale ! Au Cameroun, la grande diversité des groupes sanguins et des phénotypes hémoglobiniques (AA, AS, SS, AC, CC) illustre cette richesse génétique façonnée par l'histoire des populations et la pression de sélection (paludisme).
\end{savaistu}

\courssub{Exercices résolus et mise en application}

\begin{exoresolu}[Détermination des génotypes possibles à partir du phénotype]
\textbf{Énoncé :} Une femme de groupe sanguin A a un enfant de groupe sanguin O. Le père de l'enfant est de groupe sanguin B. Déterminer les génotypes probables de la mère, du père et de l'enfant. Justifier.

\textbf{Solution détaillée :}
\begin{enumerate}
    \item \textbf{Analyse du phénotype de l'enfant (Groupe O) :} Le groupe O correspond \textbf{uniquement} au génotype homozygote récessif $[i//i]$. L'enfant a donc \textbf{obligatoirement} reçu un allèle $i$ de sa mère \textbf{et} un allèle $i$ de son père.
    \item \textbf{Déduction pour le père (Groupe B) :} Son phénotype est B. Ses génotypes possibles sont $[I^B//I^B]$ ou $[I^B//i]$. Comme il a transmis un allèle $i$ à l'enfant, il est \textbf{nécessairement hétérozygote} : \textbf{Génotype père = $[I^B//i]$}.
    \item \textbf{Déduction pour la mère (Groupe A) :} Son phénotype est A. Ses génotypes possibles sont $[I^A//I^A]$ ou $[I^A//i]$. Comme elle a transmis un allèle $i$ à l'enfant, elle est \textbf{nécessairement hétérozygote} : \textbf{Génotype mère = $[I^A//i]$}.
    \item \textbf{Génotype de l'enfant :} \textbf{$[i//i]$}.
\end{enumerate}
\textbf{Conclusion :} Mère $[I^A//i]$, Père $[I^B//i]$, Enfant $[i//i]$. Ce couple peut avoir des enfants de groupes A ($[I^A//i]$), B ($[I^B//i]$), AB ($[I^A//I^B]$) ou O ($[i//i]$).
\end{exoresolu}

\begin{exoresolu}[Vocabulaire et relations alléliques]
\textbf{Énoncé :} Compléter le tableau suivant en utilisant les termes : \emph{locus, allèle, gène, chromosome, ADN, protéine, génotype, phénotype, dominant, récessif, homozygote, hétérozygote}.

\begin{center}
\begin{tabularx}{\linewidth}{|l|X|}\hline
\textbf{Définition / Situation} & \textbf{Terme} \\ \hline
Position d'un gène sur un chromosome & \ldots \\ \hline
Version alternative d'un gène & \ldots \\ \hline
Ensemble des allèles d'un individu & \ldots \\ \hline
Caractère observable & \ldots \\ \hline
Individu portant deux allèles identiques & \ldots \\ \hline
Allèle qui s'exprime chez l'hétérozygote & \ldots \\ \hline
Allèle masqué chez l'hétérozygote & \ldots \\ \hline
Support chimique de l'information génétique & \ldots \\ \hline
Produit final de l'expression d'un gène (souvent) & \ldots \\ \hline
\end{tabularx}
\end{center}

\textbf{Solution détaillée :}
\begin{center}
\begin{tabularx}{\linewidth}{|l|X|}\hline
\textbf{Définition / Situation} & \textbf{Terme} \\ \hline
Position d'un gène sur un chromosome & \textbf{Locus} \\ \hline
Version alternative d'un gène & \textbf{Allèle} \\ \hline
Ensemble des allèles d'un individu & \textbf{Génotype} \\ \hline
Caractère observable & \textbf{Phénotype} \\ \hline
Individu portant deux allèles identiques & \textbf{Homozygote} \\ \hline
Allèle qui s'exprime chez l'hétérozygote & \textbf{Dominant} \\ \hline
Allèle masqué chez l'hétérozygote & \textbf{Récéssif} \\ \hline
Support chimique de l'information génétique & \textbf{ADN} \\ \hline
Produit final de l'expression d'un gène (souvent) & \textbf{Protéine} \\ \hline
\end{tabularx}
\end{center}
\end{exoresolu}

\begin{exercice}[Application : Système ABO et famille]
Dans une famille, le père est de groupe A, la mère de groupe B. Ils ont 4 enfants : Paul (groupe A), Pierre (groupe B), Jacques (groupe AB) et Marie (groupe O).
\begin{enumerate}
    \item Déterminer le génotype du père et de la mère. Justifier.
    \item Déterminer le génotype de chacun des 4 enfants.
    \item Quelle est la probabilité pour qu'un 5\textsuperscript{e} enfant soit de groupe O ?
\end{enumerate}
\end{exercice}

\begin{corrige}[Exercice : Application Système ABO et famille]
\begin{enumerate}
    \item \textbf{Enfant Marie (Groupe O) :} Génotype $[i//i]$. Elle a reçu un $i$ du père et un $i$ de la mère.
    \item \textbf{Père (Groupe A) :} A transmis $i$ $\rightarrow$ Génotype $[I^A//i]$.
    \item \textbf{Mère (Groupe B) :} A transmis $i$ $\rightarrow$ Génotype $[I^B//i]$.
    \item \textbf{Enfants :}
        \begin{itemize}
            \item Paul (A) : A reçu $I^A$ du père et $i$ de la mère $\rightarrow$ $[I^A//i]$.
            \item Pierre (B) : A reçu $i$ du père et $I^B$ de la mère $\rightarrow$ $[I^B//i]$.
            \item Jacques (AB) : A reçu $I^A$ du père et $I^B$ de la mère $\rightarrow$ $[I^A//I^B]$.
            \item Marie (O) : $[i//i]$.
        \end{itemize}
    \item \textbf{Croisement :} $[I^A//i] \times [I^B//i]$.
        \begin{center}
        \begin{tabular}{|c|c|c|}\hline
        \multicolumn{1}{|c|}{} & \multicolumn{2}{c|}{\textbf{Gamètes Père}} \\ \cline{2-3}
        \textbf{Gamètes Mère} & $I^A$ & $i$ \\ \hline
        $I^B$ & $[I^A//I^B]$ (AB) & $[I^B//i]$ (B) \\ \hline
        $i$ & $[I^A//i]$ (A) & $[i//i]$ (O) \\ \hline
        \end{tabular}
        \end{center}
        4 combinaisons équiprobables. Probabilité Groupe O ($[i//i]$) = \textbf{1/4 (soit 25\%)}.
\end{enumerate}
\end{corrige}

\begin{aretenir}[Synthèse : Les gènes humains]
\begin{itemize}
    \item L'information génétique est portée par l'\textbf{ADN} organisé en \textbf{chromosomes}.
    \item Un \textbf{gène} est un segment d'ADN à un \textbf{locus} précis, codant pour une \textbf{protéine} (souvent) et déterminant un \textbf{caractère}.
    \item Un individu possède deux allèles par gène (un sur chaque chromosome homologue) : \textbf{Génotype} = $[Allèle_1 // Allèle_2]$.
    \item \textbf{Phénotype} = expression observable. Relation \textbf{Dominant / Récessif} (ou codominance).
    \item \textbf{Homozygote} (allèles identiques) vs \textbf{Hétérozygote} (allèles différents).
    \item Nombre estimé de gènes chez l'humain : \textbf{20 000 -- 25 000}.
\end{itemize}
\end{aretenir}

\textbf{Transition vers la section suivante}\par
Nous venons de voir comment les gènes s'organisent sur les chromosomes et comment leurs allèles interagissent pour produire un phénotype. Le système ABO est l'exemple type. Mais comment \textbf{prouver} concrètement le groupe sanguin d'un individu au laboratoire ? C'est l'objet du \textbf{TP N°4 : Recherche des groupes sanguins du système ABO et du facteur Rhésus}.

\coursec{Gènes et diversité humaine : le système ABO}

Nous venons de voir que les gènes, portés par les chromosomes, sont les supports de l'information génétique et que chaque gène peut exister sous plusieurs versions appelées \cle{allèles}. Mais comment cette diversité des allèles se traduit-elle concrètement dans la population humaine ? Pour le comprendre, nous allons étudier un caractère que tu connais peut-être déjà : le \cle{groupe sanguin}. À l'hôpital de district ou lors d'une campagne de don du sang, on entend souvent dire : « je suis O positif » ou « elle est AB ». D'où viennent ces groupes ? Pourquoi sont-ils différents d'une personne à l'autre ? C'est ce que nous allons découvrir grâce à une démarche d'investigation scientifique.

\courssub{Le point de départ historique : les expériences de Landsteiner}

\begin{experience}[Les mélanges de sangs de Landsteiner (1900-1901)]
En 1900, le médecin autrichien Karl \cle{Landsteiner} réalise une expérience simple mais géniale dans son laboratoire de Vienne. Il prélève le sang sur six personnes (cinq collègues et lui-même), sépare pour chacun les \cle{hématies} (globules rouges) du \cle{sérum} (le liquide jaunâtre du sang), puis mélange goutte à goutte, sur des lames de verre, le sérum de chaque personne avec les hématies de toutes les autres.

\textbf{Observations :} selon les mélanges, deux résultats sont possibles :
\begin{itemize}
\item soit les hématies restent bien dispersées : il n'y a \textbf{pas d'agglutination} ;
\item soit les hématies se collent les unes aux autres et forment des amas visibles à l'œil nu : il y a \textbf{agglutination}.
\end{itemize}

\textbf{Interprétation :} l'agglutination révèle une réaction entre une substance portée par les hématies et une substance présente dans le sérum. En classant les résultats, Landsteiner montre que les êtres humains se répartissent en plusieurs groupes, qu'il nomme A, B et O (le groupe AB sera identifié en 1902 par ses élèves Decastello et Sturli).

\textbf{Conclusion :} tous les sangs humains ne sont pas identiques : il existe au moins quatre \cle{groupes sanguins} : A, B, AB et O. C'est le \cle{système ABO}.
\end{experience}

\begin{savaistu}[Un Prix Nobel qui a sauvé des millions de vies]
Avant Landsteiner, les transfusions sanguines étaient une loterie dangereuse : parfois elles sauvaient le malade, parfois elles le tuaient brutalement. Grâce à la découverte des groupes sanguins, on a compris qu'il faut transfuser un sang \emph{compatible} avec celui du receveur. Karl Landsteiner a reçu le \textbf{Prix Nobel de médecine en 1930}. Aujourd'hui encore, dans les hôpitaux camerounais, aucune transfusion n'est faite sans déterminer au préalable le groupe sanguin du donneur et du receveur.
\end{savaistu}

\courssub{Interprétation : agglutinogènes et agglutinines}

Pourquoi certains mélanges agglutinent-ils et d'autres non ? L'analyse des résultats de Landsteiner conduit à l'hypothèse suivante, vérifiée depuis par de nombreuses expériences : les hématies portent à leur surface des substances particulières, et le sérum contient des substances capables de s'y fixer.

\begin{definition}[Agglutinogène et agglutinine]
\begin{itemize}
\item Un \cle{agglutinogène} (ou \cle{antigène}) est une substance portée à la surface des hématies. Il en existe deux types dans le système ABO : l'agglutinogène \textbf{A} et l'agglutinogène \textbf{B}.
\item Une \cle{agglutinine} (ou \cle{anticorps}) est une substance présente dans le sérum, capable d'agglutiner les hématies portant l'agglutinogène correspondant. Il en existe deux types : l'agglutinine \textbf{anti-A} et l'agglutinine \textbf{anti-B}.
\end{itemize}
L'agglutination se produit lorsqu'une agglutinine rencontre l'agglutinogène correspondant : anti-A agglutine les hématies portant A, anti-B agglutine les hématies portant B.
\end{definition}

Une règle essentielle se dégage : \textbf{un individu ne possède jamais dans son sérum l'agglutinine qui détruirait ses propres hématies}. Un sujet de groupe A possède l'agglutinogène A sur ses hématies et l'agglutinine anti-B dans son sérum, mais jamais d'anti-A. On peut alors dresser le tableau complet du système ABO.

\begin{popfigure}
\begin{center}
\begin{tikzpicture}[scale=0.98, every node/.style={font=\small}]
% En-têtes
\fill[popblue] (0,0) rectangle (3.2,0.9); \node[white,font=\small\bfseries] at (1.6,0.45) {Groupe};
\fill[popblue] (3.2,0) rectangle (7.2,0.9); \node[white,font=\small\bfseries] at (5.2,0.45) {Agglutinogènes (hématies)};
\fill[popblue] (7.2,0) rectangle (11.2,0.9); \node[white,font=\small\bfseries] at (9.2,0.45) {Agglutinines (sérum)};
\fill[popblue] (11.2,0) rectangle (15.2,0.9); \node[white,font=\small\bfseries] at (13.2,0.45) {Génotypes possibles};
% Ligne A
\fill[popblueL] (0,-1) rectangle (3.2,0); \node[font=\small\bfseries] at (1.6,-0.5) {A};
\fill[popcream] (3.2,-1) rectangle (7.2,0); \node at (5.2,-0.5) {A};
\fill[popcream] (7.2,-1) rectangle (11.2,0); \node at (9.2,-0.5) {anti-B};
\fill[popcream] (11.2,-1) rectangle (15.2,0); \node at (13.2,-0.5) {$A//A$ ou $A//O$};
% Ligne B
\fill[poporangeL] (0,-2) rectangle (3.2,-1); \node[font=\small\bfseries] at (1.6,-1.5) {B};
\fill[popcream] (3.2,-2) rectangle (7.2,-1); \node at (5.2,-1.5) {B};
\fill[popcream] (7.2,-2) rectangle (11.2,-1); \node at (9.2,-1.5) {anti-A};
\fill[popcream] (11.2,-2) rectangle (15.2,-1); \node at (13.2,-1.5) {$B//B$ ou $B//O$};
% Ligne AB
\fill[popgreenL] (0,-3) rectangle (3.2,-2); \node[font=\small\bfseries] at (1.6,-2.5) {AB};
\fill[popcream] (3.2,-3) rectangle (7.2,-2); \node at (5.2,-2.5) {A et B};
\fill[popcream] (7.2,-3) rectangle (11.2,-2); \node at (9.2,-2.5) {aucune};
\fill[popcream] (11.2,-3) rectangle (15.2,-2); \node at (13.2,-2.5) {$A//B$};
% Ligne O
\fill[poppinkL] (0,-4) rectangle (3.2,-3); \node[font=\small\bfseries] at (1.6,-3.5) {O};
\fill[popcream] (3.2,-4) rectangle (7.2,-3); \node at (5.2,-3.5) {aucun};
\fill[popcream] (7.2,-4) rectangle (11.2,-3); \node at (9.2,-3.5) {anti-A et anti-B};
\fill[popcream] (11.2,-4) rectangle (15.2,-3); \node at (13.2,-3.5) {$O//O$};
% Contours
\draw[popdark, thick] (0,0.9) rectangle (15.2,-4);
\draw[popdark] (3.2,0.9) -- (3.2,-4); \draw[popdark] (7.2,0.9) -- (7.2,-4);
\draw[popdark] (11.2,0.9) -- (11.2,-4);
\draw[popdark] (0,0) -- (15.2,0); \draw[popdark] (0,-1) -- (15.2,-1);
\draw[popdark] (0,-2) -- (15.2,-2); \draw[popdark] (0,-3) -- (15.2,-3);
\end{tikzpicture}
\end{center}
\legende{Tableau récapitulatif du système ABO : pour chaque groupe sanguin, agglutinogènes portés par les hématies, agglutinines présentes dans le sérum et génotypes possibles.}
\end{popfigure}

\begin{attention}[Le groupe O n'est pas « vide » d'anticorps]
Une erreur fréquente consiste à croire qu'un sujet de groupe O n'a « rien » dans son sang. C'est faux : ses hématies ne portent \textbf{aucun agglutinogène} (ni A ni B), mais son sérum contient \textbf{les deux agglutinines} anti-A et anti-B. C'est d'ailleurs pour cela qu'on dit que le sujet O est « donneur universel » pour les hématies : ses globules rouges, sans antigène, ne peuvent être agglutinés par le sérum d'aucun receveur.
\end{attention}

\courssub{La génétique du système ABO : un gène, trois allèles}

Pourquoi certains individus sont-ils de groupe A et d'autres de groupe O ? L'étude des familles montre que le groupe sanguin est un caractère \cle{héréditaire} : il se transmet des parents aux enfants selon des règles précises. Ce caractère est commandé par \textbf{un seul gène}, le gène ABO, qui existe dans la population sous \textbf{trois allèles} : l'allèle $A$, l'allèle $B$ et l'allèle $O$.

Rappelons qu'un individu possède \textbf{deux exemplaires} de chaque chromosome (l'un reçu du père, l'autre de la mère) : il porte donc \textbf{deux allèles} du gène ABO, qui constituent son \cle{génotype}. Mais comment ces deux allèles s'expriment-ils ?

\begin{definition}[Codominance]
Il y a \cle{codominance} lorsque deux allèles différents s'expriment \textbf{tous les deux} chez l'individu qui les porte. Dans le système ABO, les allèles $A$ et $B$ sont codominants : un individu de génotype $A//B$ fabrique à la fois l'agglutinogène A et l'agglutinogène B : il est de groupe AB.
\end{definition}

\begin{propriete}[Expression des allèles du gène ABO]
Le gène du groupe sanguin ABO possède \textbf{trois allèles} : $A$, $B$ et $O$.
\begin{itemize}
\item Les allèles $A$ et $B$ sont \textbf{codominants} : présents ensemble, ils s'expriment tous les deux (groupe AB).
\item L'allèle $O$ est \textbf{récessif} : il ne s'exprime que s'il est présent en deux exemplaires ($O//O$). En présence de $A$ ou de $B$, il est « masqué » : un individu $A//O$ est de groupe A.
\end{itemize}
On dit que $A$ et $B$ sont \textbf{dominants} par rapport à $O$.
\end{propriete}

Avec trois allèles pris deux par deux, on peut former \textbf{six génotypes} différents, qui ne donnent que \textbf{quatre phénotypes} (les quatre groupes sanguins) :

\begin{center}
\begin{tabularx}{\linewidth}{|l|X|}
\hline
\rowcolor{popblueL} \textbf{Phénotype (groupe)} & \textbf{Génotypes possibles} \\
\hline
A & $A//A$ (homozygote) ou $A//O$ (hétérozygote) \\
\hline
B & $B//B$ (homozygote) ou $B//O$ (hétérozygote) \\
\hline
AB & $A//B$ (hétérozygote) \\
\hline
O & $O//O$ (homozygote) \\
\hline
\end{tabularx}
\end{center}

\begin{attention}[Le phénotype ne révèle pas toujours le génotype]
Un individu de groupe A peut être de génotype $A//A$ \textbf{ou} $A//O$ : l'allèle $O$, récessif, ne se voit pas. De même, un individu de groupe B peut être $B//B$ ou $B//O$. Le phénotype seul ne suffit donc pas à connaître le génotype. Seuls les groupes AB ($A//B$) et O ($O//O$) correspondent à un génotype unique. Pour deviner le génotype d'un sujet A ou B, il faut étudier sa famille (ses parents ou ses enfants).
\end{attention}

\courssub{Prédire les groupes sanguins des enfants : l'échiquier de croisement}

Lors de la formation des gamètes (spermatozoïdes et ovules), les deux allèles d'un gène se séparent : chaque gamète ne reçoit qu'\textbf{un seul} allèle. Un parent de génotype $A//O$ produit donc deux sortes de gamètes : la moitié porte $A$, l'autre moitié porte $O$. À la fécondation, les allèles du père et de la mère se réunissent au hasard. L'\cle{échiquier de croisement} permet de prévoir toutes les combinaisons possibles.

\begin{methode}[Construire et exploiter un échiquier de croisement]
\begin{enumerate}
\item \textbf{Identifier les génotypes des parents} (ex. : père $A//O$, mère $B//O$). Si seuls les phénotypes sont connus, lister les génotypes possibles.
\item \textbf{Déterminer les gamètes} produits par chacun : le père donne $A$ ou $O$ ; la mère donne $B$ ou $O$.
\item \textbf{Tracer un tableau à double entrée} : les gamètes du père en colonnes, ceux de la mère en lignes.
\item \textbf{Remplir chaque case} en réunissant les deux allèles : c'est le génotype possible d'un enfant.
\item \textbf{Déterminer le phénotype} de chaque case et \textbf{calculer les proportions} : chaque case a la même probabilité, soit 1 chance sur 4 (25 \%).
\end{enumerate}
\end{methode}

\begin{popfigure}
\begin{center}
\begin{tikzpicture}[scale=1.05, every node/.style={font=\small}]
% Gamètes père (colonnes)
\node[font=\small\bfseries, popdark] at (2.5,1.6) {Gamètes du père ($A//O$)};
\fill[popblue] (1,0.5) rectangle (3,1.3); \node[white,font=\small\bfseries] at (2,0.9) {$A$};
\fill[popblue] (3,0.5) rectangle (5,1.3); \node[white,font=\small\bfseries] at (4,0.9) {$O$};
% Gamètes mère (lignes)
\node[font=\small\bfseries, popdark, rotate=90] at (-1.3,-1.5) {Gamètes de la mère ($B//O$)};
\fill[poporange] (-1,0.5) rectangle (0,-0.5); \node[white,font=\small\bfseries] at (-0.5,0) {$B$};
\fill[poporange] (-1,-0.5) rectangle (0,-2.5); \node[white,font=\small\bfseries] at (-0.5,-1.5) {$O$};
% Cases
\fill[popgreenL] (1,0.5) rectangle (3,-0.5); \node at (2,0) {$A//B$ : \textbf{AB}};
\fill[popblueL] (3,0.5) rectangle (5,-0.5); \node at (4,0) {$A//O$ : \textbf{A}};
\fill[poporangeL] (1,-0.5) rectangle (3,-2.5); \node at (2,-1.5) {$B//O$ : \textbf{B}};
\fill[poppinkL] (3,-0.5) rectangle (5,-2.5); \node at (4,-1.5) {$O//O$ : \textbf{O}};
% Contours
\draw[popdark, thick] (1,1.3) rectangle (5,-2.5);
\draw[popdark] (3,1.3) -- (3,-2.5);
\draw[popdark] (1,0.5) -- (5,0.5); \draw[popdark] (1,-0.5) -- (5,-0.5);
\draw[popdark, thick] (-1,1.3) -- (-1,-2.5) -- (0,-2.5);
\draw[popdark] (-1,0.5) -- (0,0.5); \draw[popdark] (-1,-0.5) -- (0,-0.5);
\end{tikzpicture}
\end{center}
\legende{Échiquier de croisement entre un père de génotype $A//O$ (groupe A) et une mère de génotype $B//O$ (groupe B). Chaque case représente un génotype possible pour un enfant, avec la même probabilité de 25 \%.}
\end{popfigure}

\begin{exemplebox}[Croisement $A//O \times B//O$ : quatre groupes possibles]
Un père de groupe A (génotype $A//O$) et une mère de groupe B (génotype $B//O$) peuvent avoir des enfants de \textbf{quatre groupes différents} :
\begin{itemize}
\item 25 \% de groupe AB (génotype $A//B$) ;
\item 25 \% de groupe A (génotype $A//O$) ;
\item 25 \% de groupe B (génotype $B//O$) ;
\item 25 \% de groupe O (génotype $O//O$).
\end{itemize}
Ce résultat est remarquable : un enfant peut être d'un groupe sanguin \emph{différent} de celui de ses deux parents !
\end{exemplebox}

\begin{exoresolu}[Une mère de groupe O peut-elle avoir un enfant de groupe A ?]
\textbf{Énoncé.} À la maternité, une maman de groupe O s'étonne : « Mon bébé est de groupe A, est-ce possible ? » Le père est de groupe A. Réponds à sa question en justifiant ta réponse par un raisonnement génétique.
\tcblower
\textbf{Solution détaillée.}
\begin{enumerate}
\item \textbf{Génotype de la mère :} le groupe O correspond à un seul génotype possible : $O//O$. La mère ne produit donc que des ovules portant l'allèle $O$.
\item \textbf{Génotype du père :} le groupe A correspond à deux génotypes possibles : $A//A$ ou $A//O$. Dans les deux cas, le père produit des spermatozoïdes portant l'allèle $A$ (tous s'il est $A//A$, la moitié s'il est $A//O$).
\item \textbf{Fécondation :} si un spermatozoïde portant $A$ féconde un ovule portant $O$, l'enfant reçoit le génotype $A//O$.
\item \textbf{Phénotype de l'enfant :} l'allèle $A$ étant dominant sur $O$, l'enfant $A//O$ est de \textbf{groupe A}.
\end{enumerate}
\textbf{Conclusion :} oui, une mère de groupe O peut tout à fait avoir un enfant de groupe A, si le père transmet l'allèle $A$. Le groupe sanguin de l'enfant dépend des allèles reçus des \emph{deux} parents, et non seulement de la mère.
\end{exoresolu}

\courssub{Les groupes sanguins, un reflet de la diversité humaine}

Les quatre groupes sanguins ne sont pas également répartis dans la population : c'est un bel exemple de la \cle{diversité humaine} liée à la diversité des allèles. Dans la population camerounaise, comme dans la plupart des populations africaines, le groupe O est le plus fréquent, suivi des groupes A et B, tandis que le groupe AB reste rare.

\begin{popfigure}
\begin{center}
\begin{tikzpicture}
\begin{axis}[
    ybar,
    width=0.85\linewidth,
    height=6.5cm,
    bar width=1.1cm,
    ymin=0, ymax=60,
    ylabel={Fréquence (\%)},
    xlabel={Groupe sanguin},
    symbolic x coords={O, A, B, AB},
    xtick=data,
    ytick={0,10,20,30,40,50,60},
    axis lines=left,
    nodes near coords,
    every node near coord/.append style={font=\small\bfseries, color=popdark},
    ylabel style={font=\small},
    xlabel style={font=\small},
]
\addplot[fill=popblue, draw=popdark] coordinates {(O,48) (A,26) (B,21) (AB,5)};
\end{axis}
\end{tikzpicture}
\end{center}
\legende{Fréquences indicatives des groupes sanguins ABO dans une population d'Afrique centrale (valeurs approchées, exprimées en pourcentage). Le groupe O est le plus répandu, le groupe AB le plus rare.}
\end{popfigure}

\begin{aretenir}[L'essentiel sur le système ABO]
\begin{itemize}
\item Le système ABO comprend quatre groupes sanguins : A, B, AB et O, découverts par Landsteiner (1900-1901) grâce à des mélanges sang + sérum (agglutination ou non).
\item Les hématies portent des \textbf{agglutinogènes} (A et/ou B) ; le sérum contient des \textbf{agglutinines} (anti-A et/ou anti-B), jamais dirigées contre ses propres hématies.
\item Le groupe sanguin est commandé par \textbf{un gène à trois allèles} : $A$ et $B$ codominants, $O$ récessif. Il existe \textbf{six génotypes} pour \textbf{quatre phénotypes}.
\item L'\textbf{échiquier de croisement} permet de prévoir les groupes possibles des enfants à partir des génotypes des parents.
\item La répartition inégale des groupes dans les populations illustre la \textbf{diversité génétique} de l'espèce humaine.
\end{itemize}
\end{aretenir}

Ce caractère du groupe sanguin, si facile à observer au laboratoire, va maintenant nous servir de terrain d'expérimentation : dans le TP N°4, tu détermineras toi-même les groupes sanguins ABO et le facteur Rhésus à l'aide de sérums tests, en reproduisant la démarche historique de Landsteiner.

\coursec{TP N°4 : Recherche des groupes sanguins ABO et du facteur Rhésus}

\courssub{Problématique et rappel des acquis}

Dans la section précédente, nous avons établi que le groupe sanguin ABO est déterminé par un gène à trois allèles ($I^A$, $I^B$, $i$) et que le facteur Rhésus dépend de la présence ou de l'absence d'une substance (l'antigène D, appelé aussi facteur Rh) à la surface des globules rouges. Nous savons désormais \emph{quel} génotype produit \emph{quels} caractères. Mais face à un malade qui a besoin d'une transfusion d'urgence à l'hôpital central de Yaoundé ou à l'hôpital Laquintinie de Douala, comment le laborantin détermine-t-il concrètement et rapidement son groupe sanguin ? C'est tout l'objet de ce TP : passer du génotype (invisible) au phénotype (observable) grâce à une méthode simple et standardisée : la \cle{réaction d'agglutination}.

\begin{definition}[Antigènes et anticorps du système ABO-Rh]
Les \textbf{antigènes} (appelés aussi \cle{agglutinogènes}) sont des substances portées à la surface de la membrane des globules rouges (hématies). Ils sont caractéristiques de chaque individu (sauf les vrais jumeaux).
Les \textbf{anticorps} (appelés aussi \cle{agglutinines}) sont des protéines présentes dans le plasma sanguin, capables de reconnaître et de se fixer sur un antigène étranger.
\end{definition}

\begin{propriete}[Règle fondamentale du groupage sanguin]
Un individu \textbf{ne possède jamais dans son plasma les anticorps dirigés contre ses propres antigènes}. En revanche, il possède naturellement les anticorps dirigés contre les antigènes ABO qu'il ne porte pas.
\begin{itemize}
    \item Groupe A : antigène A sur les hématies $\rightarrow$ anticorps anti-B dans le plasma.
    \item Groupe B : antigène B sur les hématies $\rightarrow$ anticorps anti-A dans le plasma.
    \item Groupe AB : antigènes A et B $\rightarrow$ \textbf{aucun} anticorps anti-A ni anti-B.
    \item Groupe O : \textbf{aucun} antigène A ni B $\rightarrow$ anticorps anti-A \textbf{et} anti-B.
    \item Rhésus + (Rh+) : antigène D présent sur les hématies.
    \item Rhésus - (Rh-) : antigène D absent ; l'individu peut fabriquer des anticorps anti-D s'il est mis en contact avec du sang Rh+.
\end{itemize}
\end{propriete}

\courssub{Principe de la méthode : la réaction d'agglutination}

\begin{experience}[Principe du test de groupage sanguin]
\textbf{Objectif :} Mettre en évidence la présence des antigènes A, B et D à la surface des hématies d'un sujet.

\textbf{Principe :} On utilise des \textbf{sérums tests} contenant des anticorps connus (anti-A, anti-B, anti-D). On les met en contact avec le sang du sujet. Si l'antigène correspondant est présent sur les hématies, les anticorps du sérum test se fixent sur plusieurs hématies à la fois et les collent entre elles : on observe des petits amas visibles à l'œil nu. C'est l'\textbf{agglutination}.

\textbf{Matériel :}
\begin{itemize}
    \item Sang du sujet (goutte prélevée au bout du doigt).
    \item Trois sérums tests : \textbf{anti-A}, \textbf{anti-B}, \textbf{anti-D (anti-Rhésus)}.
    \item Plaque de groupage (ou lame de verre à trois cases délimitées).
    \item Baguettes de mélange propres (une par case), chronomètre.
\end{itemize}

\textbf{Protocole :}
\begin{enumerate}
    \item Identifier les trois cases : \textbf{case 1 (anti-A)}, \textbf{case 2 (anti-B)}, \textbf{case 3 (anti-D)}.
    \item Déposer \textbf{une goutte de chaque sérum test} dans sa case.
    \item Ajouter \textbf{une goutte de sang du sujet} dans \textbf{chacune} des trois cases.
    \item Mélanger délicatement chaque case avec une baguette propre (changer de baguette à chaque case pour éviter les contaminations).
    \item Agiter doucement la plaque pendant \textbf{2 à 3 minutes} à température ambiante.
    \item Observer à l'œil nu (éventuellement avec une loupe).
\end{enumerate}

\textbf{Observations possibles :}
\begin{itemize}
    \item \textbf{Résultat positif (+) :} formation de petits grumeaux rouges (amas d'hématies) dans un liquide devenu clair : il y a \textbf{agglutination}.
    \item \textbf{Résultat négatif (-) :} le mélange reste homogène, rouge et liquide, sans grumeaux : \textbf{pas d'agglutination}.
\end{itemize}
\end{experience}

\begin{popfigure}
\begin{tikzpicture}[scale=0.95, every node/.style={font=\small}]
    % Case 1 : Anti-A (positif)
    \begin{scope}[xshift=-5cm]
        \draw[thick, fill=white] (-1.6,-1.6) rectangle (1.6,1.6);
        \node[align=center] at (0, 2.0) {\textbf{Case 1 : sérum anti-A}};
        \fill[popblue] (0,0.6) circle (0.18);
        \node[font=\tiny, text=white] at (0,0.6) {anti-A};
        \fill[popink] (0,-0.5) circle (0.18);
        \node[font=\tiny, text=white] at (0,-0.5) {sang};
        \draw[->, thick, popblue] (0,-1.7) -- (0,-2.2) node[midway, right, font=\scriptsize] {mélange, 3 min};
        \node[align=center] at (0,-2.7) {\textbf{Résultat : agglutination (+)}};
        \foreach \x/\y in {-0.7/-3.4, -0.2/-3.3, 0.3/-3.5, 0.8/-3.3, -0.5/-3.9, 0.1/-3.8, 0.6/-4.0} {
            \fill[popink] (\x,\y) circle (0.13);
            \fill[popink] (\x+0.12,\y+0.1) circle (0.1);
        }
    \end{scope}
    % Case 2 : Anti-B (négatif)
    \begin{scope}[xshift=0cm]
        \draw[thick, fill=white] (-1.6,-1.6) rectangle (1.6,1.6);
        \node[align=center] at (0, 2.0) {\textbf{Case 2 : sérum anti-B}};
        \fill[popgold] (0,0.6) circle (0.18);
        \node[font=\tiny] at (0,0.6) {anti-B};
        \fill[popink] (0,-0.5) circle (0.18);
        \node[font=\tiny, text=white] at (0,-0.5) {sang};
        \draw[->, thick, popblue] (0,-1.7) -- (0,-2.2) node[midway, right, font=\scriptsize] {mélange, 3 min};
        \node[align=center] at (0,-2.7) {\textbf{Résultat : négatif (-)}};
        \fill[popink!35] (-1.2,-4.3) rectangle (1.2,-3.2);
        \node[font=\scriptsize] at (0,-3.75) {mélange homogène};
    \end{scope}
    % Case 3 : Anti-D (positif)
    \begin{scope}[xshift=5cm]
        \draw[thick, fill=white] (-1.6,-1.6) rectangle (1.6,1.6);
        \node[align=center] at (0, 2.0) {\textbf{Case 3 : sérum anti-D}};
        \fill[popmuted] (0,0.6) circle (0.18);
        \node[font=\tiny, text=white] at (0,0.6) {anti-D};
        \fill[popink] (0,-0.5) circle (0.18);
        \node[font=\tiny, text=white] at (0,-0.5) {sang};
        \draw[->, thick, popblue] (0,-1.7) -- (0,-2.2) node[midway, right, font=\scriptsize] {mélange, 3 min};
        \node[align=center] at (0,-2.7) {\textbf{Résultat : agglutination (+)}};
        \foreach \x/\y in {-0.7/-3.4, -0.2/-3.3, 0.3/-3.5, 0.8/-3.3, -0.5/-3.9, 0.1/-3.8, 0.6/-4.0} {
            \fill[popink] (\x,\y) circle (0.13);
            \fill[popink] (\x+0.12,\y+0.1) circle (0.1);
        }
    \end{scope}
    % Interprétation
    \node[align=left, anchor=north west, font=\small] at (-7.2,-5.0) {
        \textbf{Interprétation de cet exemple :}\\
        Case 1 (+) : antigène A présent ; Case 2 (-) : antigène B absent $\Rightarrow$ groupe \textbf{A}.\\
        Case 3 (+) : antigène D présent $\Rightarrow$ Rhésus \textbf{+}.\\
        \textbf{Conclusion : le sujet est de groupe A Rh+ (A+).}
    };
\end{tikzpicture}
\legende{Principe du test de groupage sur plaque : on dépose dans chaque case une goutte de sérum test (anti-A, anti-B, anti-D) et une goutte de sang du sujet, on mélange, puis on lit le résultat après 2 à 3 minutes. L'exemple illustre un sujet de groupe A Rh+ (agglutination en cases 1 et 3 uniquement).}
\end{popfigure}

\courssub{Interprétation des résultats : tableau de lecture}

La lecture du test repose sur une démarche déductive simple : « \emph{quel sérum a provoqué une agglutination ?} ».

\begin{methode}[Lire un test de groupage sanguin (ABO-Rh)]
\begin{enumerate}
    \item Noter le résultat de chaque case : \textbf{+} (agglutination) ou \textbf{-} (pas d'agglutination).
    \item En déduire les antigènes présents sur les hématies :
    \begin{itemize}
        \item agglutination avec anti-A $\Rightarrow$ antigène A présent ;
        \item agglutination avec anti-B $\Rightarrow$ antigène B présent ;
        \item agglutination avec anti-D $\Rightarrow$ antigène D (Rhésus) présent.
    \end{itemize}
    \item En déduire le groupe ABO :
    \begin{itemize}
        \item A présent, B absent $\Rightarrow$ groupe \textbf{A} ;
        \item A absent, B présent $\Rightarrow$ groupe \textbf{B} ;
        \item A et B présents $\Rightarrow$ groupe \textbf{AB} ;
        \item A et B absents $\Rightarrow$ groupe \textbf{O}.
    \end{itemize}
    \item Ajouter le facteur Rhésus : anti-D (+) $\Rightarrow$ \textbf{Rh+} ; anti-D (-) $\Rightarrow$ \textbf{Rh-}.
    \item Écrire le phénotype complet, par exemple \textbf{A+} ou \textbf{O-}.
\end{enumerate}
\end{methode}

\begin{center}
\begin{tabularx}{\linewidth}{|X|c|c|c|c|}\hline
\rowcolor{popblueL}
\textbf{Groupe recherché} & \textbf{Anti-A} & \textbf{Anti-B} & \textbf{Anti-D} & \textbf{Phénotype} \\ \hline
A Rh+ & + & - & + & A+ \\ \hline
A Rh- & + & - & - & A- \\ \hline
B Rh+ & - & + & + & B+ \\ \hline
B Rh- & - & + & - & B- \\ \hline
AB Rh+ & + & + & + & AB+ \\ \hline
AB Rh- & + & + & - & AB- \\ \hline
O Rh+ & - & - & + & O+ \\ \hline
O Rh- & - & - & - & O- \\ \hline
\end{tabularx}
\end{center}

\begin{exemplebox}[Exemple d'interprétation]
Un technicien du Centre National de Transfusion Sanguine (CNTS) de Yaoundé obtient pour un donneur :
\begin{itemize}
    \item case anti-A : \textbf{agglutination (+)} ;
    \item case anti-B : \textbf{agglutination (+)} ;
    \item case anti-D : \textbf{pas d'agglutination (-)}.
\end{itemize}
\textbf{Raisonnement :} les hématies portent l'antigène A (réaction avec l'anti-A) \textbf{et} l'antigène B (réaction avec l'anti-B). Elles ne portent pas l'antigène D.
\textbf{Conclusion :} le donneur est de groupe \textbf{AB Rh- (AB-)}.
\end{exemplebox}

\courssub{Pourquoi les hématies s'agglutinent-elles ?}

Un anticorps possède \textbf{deux sites de fixation} identiques (sa molécule a la forme d'un Y). Chaque site peut se fixer sur l'antigène d'une hématie différente. L'anticorps forme ainsi un \textbf{pont} entre deux globules rouges. En se répétant des milliers de fois, ce pontage assemble les hématies en amas visibles à l'œil nu : c'est l'agglutination. Si l'antigène correspondant est absent, aucun pont ne se forme et le mélange reste homogène.

\begin{popfigure}
\begin{tikzpicture}[scale=1.0, every node/.style={font=\small}]
    % Hématie 1
    \draw[thick, fill=popink!25] (-2.2,0) ellipse (1.0 and 0.55);
    \node at (-2.2,0) {hématie};
    % Antigènes hématie 1
    \foreach \p in {(-2.9,0.45),(-2.2,0.55),(-1.5,0.45)} {
        \fill[poporange] \p circle (0.09);
    }
    % Hématie 2
    \draw[thick, fill=popink!25] (2.2,0) ellipse (1.0 and 0.55);
    \node at (2.2,0) {hématie};
    \foreach \p in {(1.5,0.45),(2.2,0.55),(2.9,0.45)} {
        \fill[poporange] \p circle (0.09);
    }
    % Anticorps en Y reliant les deux hématies
    \draw[very thick, popblue] (-1.5,0.5) -- (0,1.4);
    \draw[very thick, popblue] (1.5,0.5) -- (0,1.4);
    \draw[very thick, popblue] (0,1.4) -- (0,2.2);
    \node[popblue, align=center] at (0,2.6) {anticorps (forme de Y)};
    % Légendes
    \node[align=center, font=\scriptsize] at (-3.6,1.0) {antigène};
    \draw[->, poporange, thick] (-3.3,0.85) -- (-2.95,0.5);
    \node[align=center] at (0,-1.3) {L'anticorps relie deux hématies : formation d'un amas (agglutination).};
\end{tikzpicture}
\legende{Mécanisme de l'agglutination : l'anticorps, grâce à ses deux sites de fixation, forme un pont entre les antigènes de deux hématies. La répétition de ces ponts produit des amas de globules rouges visibles à l'œil nu.}
\end{popfigure}

\courssub{Compatibilité transfusionnelle : qui peut donner à qui ?}

Connaître son groupe ne suffit pas ; il faut savoir \emph{quel sang peut être transfusé à quel malade}. La règle d'or est la suivante : \textbf{le receveur ne doit jamais recevoir d'hématies portant un antigène contre lequel son plasma possède des anticorps}. Sinon, les anticorps du receveur agglutinent puis détruisent les hématies du donneur : c'est l'\cle{accident transfusionnel}, une urgence grave pouvant entraîner la mort.

\begin{propriete}[Règles de compatibilité ABO-Rhésus (donneur $\rightarrow$ receveur)]
\begin{itemize}
    \item \textbf{Donneur universel : O-}. Ses hématies ne portent \textbf{aucun} antigène A, B ni D : elles ne peuvent être attaquées par aucun anticorps du receveur. Le sang O- peut donc être transfusé à \textbf{tous} les groupes (en urgence).
    \item \textbf{Receveur universel : AB+}. Son plasma ne contient \textbf{aucun} anticorps anti-A, anti-B ni anti-D : il peut recevoir les hématies de \textbf{tous} les groupes.
    \item \textbf{Règle du Rhésus :} un receveur Rh- ne doit pas recevoir de sang Rh+ ; un receveur Rh+ peut recevoir du sang Rh+ ou Rh-.
    \item En pratique, on transfuse toujours, quand c'est possible, du sang \textbf{de même groupe} que le receveur (transfusion isogroupe).
\end{itemize}
\end{propriete}

\begin{center}
\begin{tabularx}{\linewidth}{|l|X|}\hline
\rowcolor{popblueL}
\textbf{Groupe du receveur} & \textbf{Groupes de donneurs acceptés (globules rouges)} \\ \hline
O- & O- \\ \hline
O+ & O-, O+ \\ \hline
A- & O-, A- \\ \hline
A+ & O-, O+, A-, A+ \\ \hline
B- & O-, B- \\ \hline
B+ & O-, O+, B-, B+ \\ \hline
AB- & O-, A-, B-, AB- \\ \hline
AB+ & tous les groupes (receveur universel) \\ \hline
\end{tabularx}
\end{center}

\begin{attention}[Erreurs fréquentes et pièges]
\begin{enumerate}
    \item \textbf{Confondre donneur et receveur universel :} O- est le \emph{donneur} universel (ses hématies n'ont aucun antigène) ; AB+ est le \emph{receveur} universel (son plasma n'a aucun anticorps). Ne jamais les inverser.
    \item \textbf{Croire qu'un AB+ peut donner à tout le monde :} c'est faux ! Ses hématies portent les antigènes A, B et D ; elles seraient détruites chez un receveur O, A ou B.
    \item \textbf{Contamination croisée au TP :} utiliser la même baguette pour mélanger les trois cases fausse tous les résultats. Une baguette par case.
    \item \textbf{Lire trop tard :} au-delà de quelques minutes, le mélange sèche et peut simuler de faux amas. Lire le résultat après 2 à 3 minutes.
\end{enumerate}
\end{attention}

\courssub{Incompatibilité Rhésus entre la mère et le fœtus}

\begin{definition}[Incompatibilité fœto-maternelle Rhésus]
Situation à risque : une mère \textbf{Rh-} porte un enfant \textbf{Rh+} (le caractère Rh+ étant hérité du père).
\begin{enumerate}
    \item Lors de l'accouchement, de petites quantités de sang du bébé Rh+ peuvent passer dans le sang de la mère.
    \item La mère, qui ne possède pas l'antigène D, fabrique alors des \textbf{anticorps anti-D} : on dit qu'elle est \emph{immunisée}.
    \item Lors d'une \textbf{grossesse suivante}, si le fœtus est à nouveau Rh+, ces anticorps maternels traversent le placenta et détruisent les globules rouges du fœtus : c'est la \cle{maladie hémolytique du nouveau-né}, qui peut être grave (anémie sévère, jaunisse intense, voire mort du fœtus).
\end{enumerate}
\end{definition}

\begin{propriete}[Prévention]
On évite cette maladie en injectant à la mère Rh-, juste après l'accouchement d'un enfant Rh+, des \textbf{anticorps anti-D} qui détruisent immédiatement les hématies fœtales passées dans son sang, avant que son organisme ne fabrique ses propres anticorps. C'est pourquoi le groupage sanguin (avec le facteur Rhésus) fait partie des examens obligatoires de la femme enceinte lors des consultations prénatales.
\end{propriete}

\begin{attention}[Ne pas confondre]
L'incompatibilité \textbf{ABO} entre mère et enfant (par exemple mère O, enfant A) est fréquente mais le plus souvent \textbf{bénigne}. L'incompatibilité \textbf{Rhésus} (mère Rh-, enfant Rh+) est plus rare mais peut être \textbf{grave} à partir de la deuxième grossesse : elle impose une surveillance médicale et une prévention systématique.
\end{attention}

\courssub{Enjeux de santé publique au Cameroun}

\begin{savaistu}[Le nom « Rhésus » et le don de sang au Cameroun]
Le facteur Rhésus a été découvert en \textbf{1940} par Karl Landsteiner (déjà découvreur des groupes ABO) et Alexander Wiener, grâce à des expériences réalisées avec le sang d'un singe, le \textbf{macaque rhésus} : d'où son nom.
Au Cameroun, le \textbf{Centre National de Transfusion Sanguine (CNTS)} et les banques de sang des hôpitaux font face à une pénurie chronique de sang. Le groupe \textbf{O+} est le plus fréquent dans la population, tandis que le \textbf{O-} (donneur universel) est rare et très recherché pour les urgences. \textbf{Donner son sang est un acte citoyen qui sauve des vies}, notamment lors des hémorragies de l'accouchement, des accidents de la route et des anémies graves dues au paludisme.
\end{savaistu}

\begin{exoresolu}[Détermination du phénotype et compatibilité]
\textbf{Énoncé :} Un blessé d'un accident de la route arrive aux urgences de l'hôpital régional de Bafoussam. Le groupage rapide donne les résultats suivants :
\begin{itemize}
    \item case anti-A : agglutination (+) ;
    \item case anti-B : pas d'agglutination (-) ;
    \item case anti-D : agglutination (+).
\end{itemize}
Le stock disponible contient des poches de tous les groupes : O-, O+, A-, A+, B-, B+, AB-, AB+.
\begin{enumerate}
    \item Déterminer le groupe sanguin du blessé.
    \item Quelles poches peuvent lui être transfusées sans risque ? Justifier.
    \item Peut-on utiliser du sang AB+ ? Justifier.
\end{enumerate}
\tcblower
\textbf{Solution détaillée :}
\begin{enumerate}
    \item \textbf{Interprétation du groupage :}
        \begin{itemize}
            \item anti-A (+) $\Rightarrow$ antigène A présent ;
            \item anti-B (-) $\Rightarrow$ antigène B absent ;
            \item donc groupe \textbf{A} ;
            \item anti-D (+) $\Rightarrow$ antigène D présent, donc Rhésus \textbf{+}.
        \end{itemize}
        \textbf{Phénotype du blessé : A+ (A Rh+).}
    \item \textbf{Recherche des donneurs compatibles :} le receveur A+ possède dans son plasma des anticorps \textbf{anti-B}. Il ne doit donc recevoir \textbf{aucune hématie portant l'antigène B}. Étant Rh+, il accepte le sang Rh+ ou Rh-.
        Les poches utilisables sont donc : \textbf{O-, O+, A-, A+}.
        Les poches \textbf{interdites} (portant l'antigène B) sont : B-, B+, AB-, AB+ ; elles seraient agglutinées puis détruites par les anticorps anti-B du malade (accident transfusionnel).
    \item \textbf{Sang AB+ : interdit.} Les hématies AB+ portent les antigènes A, B \textbf{et} D. Or le malade ne possède pas l'antigène B et son plasma contient des anticorps anti-B : les hématies transfusées seraient massivement détruites. C'est le groupe AB+ qui est \emph{receveur} universel, et non pas donneur universel.
\end{enumerate}
\end{exoresolu}

\courssub{Rédaction du compte rendu de TP (modèle attendu au BEPC)}

\begin{exemplebox}[Modèle de compte rendu : TP N°4 — Groupage sanguin]
\textbf{Titre :} Détermination du groupe sanguin ABO et du facteur Rhésus d'un sujet par la méthode d'agglutination sur plaque.

\textbf{I. Problématique :} avant toute transfusion, il faut connaître le groupe sanguin du malade pour éviter les accidents transfusionnels. Comment déterminer le groupe sanguin d'un individu ?

\textbf{Hypothèse :} la mise en contact du sang du sujet avec des sérums tests (anti-A, anti-B, anti-D) révèle, par agglutination, les antigènes présents sur ses hématies.

\textbf{II. Matériel et protocole :} sang du sujet ; sérums anti-A, anti-B, anti-D ; plaque à trois cases ; baguettes ; chronomètre. Une goutte de sérum et une goutte de sang par case ; mélange ; lecture après 2 à 3 minutes.

\textbf{III. Résultats :}
\begin{center}
\begin{tabular}{|l|c|c|}\hline
\textbf{Sérum test} & \textbf{Observation} & \textbf{Antigène} \\ \hline
Anti-A & agglutination (+) & A présent \\ \hline
Anti-B & pas d'agglutination (-) & B absent \\ \hline
Anti-D & agglutination (+) & D présent \\ \hline
\end{tabular}
\end{center}

\textbf{IV. Interprétation et conclusion :} l'agglutination avec l'anti-A prouve la présence de l'antigène A ; l'absence d'agglutination avec l'anti-B prouve l'absence de l'antigène B : le sujet est donc de groupe \textbf{A}. L'agglutination avec l'anti-D prouve la présence de l'antigène D : le sujet est \textbf{Rh+}.
\textbf{Conclusion : le sujet est de groupe A Rh+ (A+).} Il peut recevoir du sang O-, O+, A- ou A+ et donner aux sujets A+ et AB+.
\end{exemplebox}

\courssub{Synthèse des compétences évaluées}

\begin{aretenir}[Bilan de la leçon — Expression de l'information génétique : les groupes sanguins]
\begin{enumerate}
    \item \textbf{Savoir :} le groupe sanguin (phénotype) est l'expression des allèles des gènes ABO et Rhésus portés par les chromosomes. Il se traduit par la présence ou l'absence des antigènes A, B et D à la surface des hématies.
    \item \textbf{Savoir-faire (TP N°3) :} construire une maquette de caryotype humain (46 chromosomes, 23 paires, dont les chromosomes sexuels XX ou XY).
    \item \textbf{Savoir-faire (TP N°4) :} réaliser et interpréter un test de groupage par agglutination (anti-A, anti-B, anti-D), remplir un tableau de lecture et rédiger un compte rendu structuré (problématique, hypothèse, protocole, résultats, conclusion).
    \item \textbf{Savoir-faire (transfusion) :} utiliser les règles de compatibilité ; identifier le donneur universel (O-) et le receveur universel (AB+) ; expliquer l'accident transfusionnel par la rencontre antigène-anticorps.
    \item \textbf{Savoir-être :} comprendre l'importance du don de sang bénévole au Cameroun et du dépistage de l'incompatibilité Rhésus chez la femme enceinte.
\end{enumerate}
\end{aretenir}

\begin{exercice}[Entraînement BEPC]
\textbf{Sujet :} Une femme de groupe B Rh- accouche d'un enfant de groupe A Rh+. Le père est de groupe AB Rh+.
\begin{enumerate}
    \item Indiquer les antigènes portés par les hématies de chacun des trois membres de cette famille.
    \item Cette mère risque-t-elle une incompatibilité Rhésus lors d'une grossesse future ? Justifier.
    \item Le père peut-il donner son sang à l'enfant ? Justifier à l'aide des règles de compatibilité.
\end{enumerate}
\emph{(Corrigé disponible dans l'espace « Corrections » de l'application OnBuch+.)}
\end{exercice}

\coursec{Activités d'intégration et synthèse}

Au cours du TP N°4, tu as déterminé expérimentalement les groupes sanguins de plusieurs échantillons à l'aide des sérums testeurs anti-A, anti-B et anti-D. Tu as observé que certains échantillons agglutinent, d'autres non, et tu as conclu à un groupe : A, B, AB ou O, Rhésus positif ou négatif. Mais une question demeure : \emph{comment ces résultats s'expliquent-ils au niveau des chromosomes et des gènes ?} Cette dernière séance relie tous les savoirs de la leçon — chromosomes, gènes, allèles, groupes sanguins — en une chaîne logique unique, puis te propose des exercices de synthèse et des applications concrètes à la vie courante : don de sang, transfusion, test de paternité, mariage. C'est aussi l'occasion d'apprendre à rédiger une conclusion scientifique rigoureuse, compétence exigée au BEPC.

\courssub{Bilan de la leçon : de l'ADN à la diversité humaine}

\courssub{La chaîne logique de l'information génétique}

Toute la leçon tient en une phrase : \textbf{l'information génétique, portée par les gènes des chromosomes, s'exprime en caractères observables qui varient d'un individu à l'autre}. Décomposons cette chaîne maillon par maillon :

\begin{enumerate}
\item \cle{ADN} : la molécule qui porte l'information génétique, contenue dans le noyau de chaque cellule.
\item \cle{Chromosome} : un filament d'ADN condensé, visible au microscope pendant la division cellulaire. L'espèce humaine possède \textbf{46 chromosomes} (23 paires) : 22 paires d'autosomes et 1 paire de chromosomes sexuels (XX chez la femme, XY chez l'homme).
\item \cle{Gène} : un segment de chromosome qui commande la réalisation d'un caractère (par exemple le groupe sanguin).
\item \cle{Allèle} : une version possible d'un gène. Pour le système ABO, le gène existe en trois allèles : $A$, $B$ et $O$.
\item \cle{Génotype} : l'ensemble des allèles d'un individu pour un caractère (exemple : $A/O$).
\item \cle{Phénotype} : le caractère observable qui résulte de l'expression du génotype (exemple : groupe [A]).
\item \cle{Diversité humaine} : comme les allèles se combinent différemment chez chaque individu, les phénotypes varient au sein de l'espèce humaine.
\end{enumerate}

\begin{aretenir}[La chaîne à retenir]
\[\text{ADN} \rightarrow \text{Chromosome} \rightarrow \text{Gène} \rightarrow \text{Allèles} \rightarrow \text{Génotype} \rightarrow \text{Phénotype} \rightarrow \text{Diversité humaine}\]
Chaque flèche signifie « est porté par » ou « s'exprime en ».
\end{aretenir}

\begin{popfigure}
\begin{center}
\begin{tikzpicture}[mindmap, grow cyclic, every node/.style={concept}, concept color=popblueL, text=popdark,
  root concept/.append style={concept color=popblue, text=white, font=\bfseries},
  level 1/.append style={level distance=3.4cm, sibling angle=60},
  level 2/.append style={level distance=2.4cm, sibling angle=45, concept color=popgreenL, font=\small}]
\node[root concept]{Information génétique}
  child { node {Chromosomes} child { node {46 = 23 paires} } child { node {XX ou XY} } }
  child { node {Gène} child { node {segment de chromosome} } }
  child { node {Allèles} child { node {$A$, $B$, $O$} } }
  child { node {Génotype} child { node {ex. $A/O$} } }
  child { node {Phénotype} child { node {ex. groupe [A]} } }
  child { node {Diversité} child { node {ABO + Rhésus} } };
\end{tikzpicture}
\end{center}
\legende{Carte mentale récapitulative de la leçon : l'information génétique, portée par les chromosomes, s'exprime par les gènes et leurs allèles en phénotypes variés (groupes sanguins), source de la diversité humaine.}
\end{popfigure}

\courssub{Rappels essentiels sur le système ABO et le facteur Rhésus}

\begin{definition}[Dominance et codominance]
Un allèle est \cle{dominant} s'il s'exprime même en un seul exemplaire ; il est \cle{récessif} s'il ne s'exprime qu'en deux exemplaires. Deux allèles sont \cle{codominants} s'ils s'expriment tous les deux quand ils sont présents ensemble.
\end{definition}

Pour le système ABO : les allèles $A$ et $B$ sont \textbf{codominants} entre eux et \textbf{dominants} sur l'allèle $O$, qui est récessif. D'où le tableau des génotypes :

\begin{center}
\begin{tabularx}{\linewidth}{|l|X|X|}
\hline
\rowcolor{popblueL} \textbf{Phénotype (groupe)} & \textbf{Génotypes possibles} & \textbf{Agglutinogènes sur les hématies} \\
\hline
[A] & $A/A$ ou $A/O$ & A \\
\hline
[B] & $B/B$ ou $B/O$ & B \\
\hline
[AB] & $A/B$ & A et B \\
\hline
[O] & $O/O$ & aucun \\
\hline
\end{tabularx}
\end{center}

Pour le facteur Rhésus : l'allèle $Rh^+$ (noté $D$) est dominant sur l'allèle $Rh^-$ (noté $d$). Un individu Rhésus positif est de génotype $D/D$ ou $D/d$ ; un individu Rhésus négatif est obligatoirement $d/d$.

\begin{attention}[Ne confonds pas génotype et phénotype]
Le phénotype est ce que l'on \emph{observe} (groupe [A] au groupage) ; le génotype est ce que l'individu \emph{possède} dans ses chromosomes ($A/A$ ou $A/O$). Deux individus de même phénotype peuvent avoir des génotypes différents : c'est pourquoi on ne peut pas toujours déterminer le génotype exact à partir du seul groupage.
\end{attention}

\courssub{Exercices de synthèse}

\begin{exoresolu}[Exercice 1 : lecture d'un caryotype]
Sur une maquette de caryotype humain réalisée au TP N°3, un élève compte 46 chromosomes rangés en 23 paires. La dernière paire est formée de deux chromosomes de tailles différentes : un grand chromosome X et un petit chromosome Y.
\begin{enumerate}
\item Quel est le sexe de cet individu ? Justifie.
\item Comment appelle-t-on les 22 premières paires ?
\item Quel est le nombre de chromosomes contenus dans un spermatozoïde de cet individu ?
\end{enumerate}
\tcblower
\textbf{Solution détaillée.}
\begin{enumerate}
\item La paire sexuelle est XY : les deux chromosomes sont différents. Or la formule chromosomique de l'homme est 44 autosomes + XY, celle de la femme 44 autosomes + XX. Cet individu est donc de \textbf{sexe masculin}.
\item Les 22 premières paires sont les \textbf{autosomes} (chromosomes non sexuels), communs aux deux sexes.
\item Les gamètes ne contiennent qu'un chromosome de chaque paire : un spermatozoïde contient donc $\frac{46}{2} = 23$ chromosomes, soit 22 autosomes + un chromosome sexuel (X ou Y).
\end{enumerate}
\end{exoresolu}

\begin{exoresolu}[Exercice 2 : génotypes et échiquier de croisement]
Une femme de groupe [A], de génotype $A/O$, épouse un homme de groupe [B], de génotype $B/O$.
\begin{enumerate}
\item Quels sont les groupes sanguins possibles de leurs enfants ? Construis l'échiquier de croisement.
\item Quelle est la probabilité pour qu'un enfant soit de groupe [O] ?
\end{enumerate}
\tcblower
\textbf{Solution détaillée.}
\begin{enumerate}
\item La mère $A/O$ produit deux types d'ovules : porteurs de $A$ ou de $O$. Le père $B/O$ produit deux types de spermatozoïdes : porteurs de $B$ ou de $O$. L'échiquier de croisement donne :
\begin{center}
\begin{tabular}{|c|c|c|}
\hline
\rowcolor{popblueL} \textbf{Ovules $\backslash$ Spermatozoïdes} & $B$ & $O$ \\
\hline
$A$ & $A/B$ [AB] & $A/O$ [A] \\
\hline
$O$ & $B/O$ [B] & $O/O$ [O] \\
\hline
\end{tabular}
\end{center}
Les enfants peuvent donc être des quatre groupes : [AB], [A], [B] ou [O], chacun avec la même probabilité $\frac{1}{4}$.
\item La probabilité d'avoir un enfant de groupe [O] (génotype $O/O$) est de $\frac{1}{4}$, soit 25 \%.
\end{enumerate}
\textbf{Conclusion :} deux parents de groupes [A] et [B] peuvent avoir un enfant de groupe [O] si tous deux sont hétérozygotes ($A/O$ et $B/O$).
\end{exoresolu}

\begin{exemplebox}[Un enfant AB ne peut pas avoir deux parents O]
Un couple affirme être tous les deux de groupe [O], mais leur enfant est de groupe [AB]. Est-ce possible ? Raisonnons par échiquier : deux parents [O] sont obligatoirement de génotype $O/O$ (l'allèle $O$ est récessif, il ne s'exprime qu'en deux exemplaires).
\begin{center}
\begin{tabular}{|c|c|c|}
\hline
\rowcolor{popblueL} \textbf{Ovules $\backslash$ Spermatozoïdes} & $O$ & $O$ \\
\hline
$O$ & $O/O$ [O] & $O/O$ [O] \\
\hline
$O$ & $O/O$ [O] & $O/O$ [O] \\
\hline
\end{tabular}
\end{center}
Tous les enfants possibles sont $O/O$, donc de groupe [O]. Un enfant [AB] possède les allèles $A$ et $B$, qu'aucun des deux parents ne peut transmettre. \textbf{Conclusion : un enfant [AB] ne peut pas avoir deux parents de groupe [O].} Il y a eu soit une erreur de groupage, soit une erreur d'attribution de l'enfant.
\end{exemplebox}

\begin{exoresolu}[Exercice 3 : interprétation d'un groupage]
Au laboratoire du centre de santé, on dépose le sang de quatre patients sur des lames, en présence de trois sérums testeurs. Résultats (+ = agglutination ; $-$ = pas d'agglutination) :
\begin{center}
\begin{tabular}{|l|c|c|c|}
\hline
\rowcolor{popblueL} \textbf{Patient} & \textbf{Sérum anti-A} & \textbf{Sérum anti-B} & \textbf{Sérum anti-D} \\
\hline
Mme Ngo & + & $-$ & + \\
\hline
M. Etoundi & $-$ & $-$ & + \\
\hline
Mme Abena & + & + & $-$ \\
\hline
M. Kamga & $-$ & + & + \\
\hline
\end{tabular}
\end{center}
Détermine le groupe sanguin complet (ABO et Rhésus) de chaque patient.
\tcblower
\textbf{Solution détaillée.} Le principe : l'agglutination avec un sérum révèle la présence de l'agglutinogène correspondant sur les hématies.
\begin{itemize}
\item \textbf{Mme Ngo} : agglutination avec anti-A seulement $\Rightarrow$ agglutinogène A $\Rightarrow$ groupe [A] ; agglutination avec anti-D $\Rightarrow$ Rhésus positif. \textbf{Groupe : A Rh$^+$.}
\item \textbf{M. Etoundi} : aucune agglutination avec anti-A ni anti-B $\Rightarrow$ groupe [O] ; agglutination avec anti-D $\Rightarrow$ \textbf{groupe O Rh$^+$}.
\item \textbf{Mme Abena} : agglutination avec anti-A et anti-B $\Rightarrow$ groupe [AB] ; pas d'agglutination avec anti-D $\Rightarrow$ \textbf{groupe AB Rh$^-$}.
\item \textbf{M. Kamga} : agglutination avec anti-B seulement $\Rightarrow$ groupe [B] ; anti-D positif $\Rightarrow$ \textbf{groupe B Rh$^+$}.
\end{itemize}
\end{exoresolu}

\courssub{Applications à la vie courante}

\courssub{Don de sang et transfusions}

Lors d'une transfusion, les hématies du donneur ne doivent pas être agglutinées par les agglutinines du receveur : une telle agglutination obstrue les vaisseaux et peut entraîner la mort. On retient :

\begin{itemize}
\item L'individu de groupe [O] ne possède aucun agglutinogène : c'est le \cle{donneur universel} (pour les hématies).
\item L'individu de groupe [AB] ne possède aucune agglutinine : c'est le \cle{receveur universel}.
\item Le facteur Rhésus doit aussi être respecté : un receveur Rh$^-$ ne doit pas recevoir de sang Rh$^+$.
\end{itemize}

\begin{savaistu}[Les transfusions au Cameroun]
Dans les hôpitaux camerounais, chaque poche de sang donnée est systématiquement groupée (ABO et Rhésus) avant toute transfusion. Le don de sang volontaire et gratuit est encouragé, car les besoins sont importants : accidents de la route, accouchements difficiles, paludisme grave chez l'enfant. Le groupe le plus fréquent dans la population camerounaise est le groupe [O] Rh$^+$.
\end{savaistu}

\courssub{Tests de paternité : les limites du système ABO}

\begin{attention}[Le système ABO peut exclure, jamais prouver]
Le groupage sanguin permet d'\textbf{exclure} une paternité quand les groupes sont incompatibles : par exemple, un enfant [AB] ne peut pas être le fils d'un homme [O] (l'allèle $A$ ou $B$ de l'enfant devrait venir du père). En revanche, si les groupes sont compatibles, on ne peut rien conclure : des milliers d'hommes du même groupe pourraient être le père. Le système ABO \textbf{ne prouve jamais une paternité à lui seul} ; seule l'analyse de l'ADN le permet aujourd'hui.
\end{attention}

\courssub{Mariage et incompatibilité Rhésus}

Un problème se pose lorsqu'une \textbf{femme Rhésus négatif} ($d/d$) attend un enfant \textbf{Rhésus positif} (le père étant Rh$^+$). Lors d'un premier accouchement, des hématies Rh$^+$ du fœtus peuvent passer dans le sang de la mère, qui fabrique alors des anticorps anti-D. Lors d'une grossesse suivante avec un autre fœtus Rh$^+$, ces anticorps traversent le placenta et détruisent les hématies du fœtus : c'est la \textbf{maladie hémolytique du nouveau-né}. C'est pourquoi, avant le mariage ou en début de grossesse, le médecin demande le groupage sanguin complet du couple : si la femme est Rh$^-$, une surveillance médicale (et une injection préventive d'immunoglobulines anti-D) protège les enfants suivants.

\courssub{Méthodologie : rédiger une réponse scientifique au BEPC}

Au BEPC, beaucoup d'élèves perdent des points non pas parce qu'ils ne savent pas, mais parce qu'ils rédigent mal leur raisonnement. Voici la démarche complète à suivre devant tout problème de génétique.

\begin{methode}[Résoudre un problème de génétique en cinq étapes]
\begin{enumerate}
\item \textbf{Dégager le problème} : reformule la question posée (exemple : « le groupe sanguin de l'enfant est-il compatible avec celui du père présumé ? »).
\item \textbf{Formuler une hypothèse} : propose une réponse provisoire à vérifier (exemple : « l'homme pourrait être le père »).
\item \textbf{Construire la démarche} : identifie les génotypes possibles des parents à partir des phénotypes, puis construis l'échiquier de croisement.
\item \textbf{Exploiter les résultats} : lis dans l'échiquier les génotypes et phénotypes possibles des enfants, avec leurs proportions.
\item \textbf{Conclure en justifiant} : compare les résultats aux données de l'énoncé et réponds clairement à la question posée.
\end{enumerate}
\end{methode}

\begin{methode}[Rédiger une conclusion scientifique en trois temps]
Une bonne conclusion se rédige toujours en trois phrases :
\begin{enumerate}
\item \textbf{Le résultat} : « On observe que... » (rappelle les faits, les données, sans les commenter).
\item \textbf{L'interprétation} : « Cela signifie que... » (explique le résultat à l'aide des savoirs de la leçon : dominance, récessivité, transmission des allèles).
\item \textbf{La généralisation} : « Donc, de façon générale,... » (énonce la règle ou la loi qui dépasse le cas particulier étudié).
\end{enumerate}
\textbf{Exemple rédigé :} « On observe que les deux parents de groupe [O] n'ont que des enfants de génotype $O/O$. Cela signifie qu'aucun des deux parents ne possède les allèles $A$ ou $B$ à transmettre. Donc, de façon générale, deux parents de groupe [O] ne peuvent avoir que des enfants de groupe [O]. »
\end{methode}

\begin{attention}[Pièges fréquents au BEPC]
\begin{itemize}
\item Écrire « le parent donne son groupe sanguin à l'enfant » : \textbf{faux}. Le parent transmet un \emph{allèle} par gamète, pas son phénotype.
\item Oublier qu'un individu [A] peut être $A/A$ \emph{ou} $A/O$ : pense toujours aux deux génotypes possibles avant de construire l'échiquier.
\item Affirmer qu'un test de groupe sanguin « prouve » la paternité : il ne peut que l'\emph{exclure}.
\item Conclure sans justifier : toute conclusion doit s'appuyer sur l'échiquier ou sur les résultats du groupage.
\end{itemize}
\end{attention}

\courssub{Auto-évaluation sur les objectifs de la leçon}

Coche mentalement chaque objectif que tu te sens capable de réaliser seul, sans le cours. Si un point te résiste, relis la section correspondante avant l'évaluation.

\begin{center}
\begin{tabularx}{\linewidth}{|X|c|c|}
\hline
\rowcolor{popblueL} \textbf{Objectif de la leçon} & \textbf{Acquis} & \textbf{À revoir} \\
\hline
Décrire la formule chromosomique de l'espèce humaine (46 chromosomes, 23 paires, XX ou XY) & $\square$ & $\square$ \\
\hline
Réaliser et exploiter une maquette de caryotype humain (TP N°3) & $\square$ & $\square$ \\
\hline
Définir gène, allèle, génotype, phénotype, dominance, récessivité, codominance & $\square$ & $\square$ \\
\hline
Citer les génotypes possibles de chaque groupe du système ABO & $\square$ & $\square$ \\
\hline
Déterminer expérimentalement un groupe sanguin ABO et Rhésus (TP N°4) & $\square$ & $\square$ \\
\hline
Construire et exploiter un échiquier de croisement & $\square$ & $\square$ \\
\hline
Expliquer les règles du don de sang et de la transfusion & $\square$ & $\square$ \\
\hline
Expliquer l'incompatibilité Rhésus mère-enfant et ses conséquences & $\square$ & $\square$ \\
\hline
Rédiger une conclusion scientifique en trois temps (résultat, interprétation, généralisation) & $\square$ & $\square$ \\
\hline
\end{tabularx}
\end{center}

\begin{exercice}[Problème de synthèse final]
Dans une maternité de Yaoundé, on soupçonne un échange de nouveau-nés entre deux familles. La famille Mballa : le père est de groupe [O], la mère de groupe [A] ($A/O$) ; leur enfant présumé est de groupe [AB]. La famille Fouda : le père est de groupe [B] ($B/O$), la mère de groupe [A] ($A/A$) ; leur enfant présumé est de groupe [O].
\begin{enumerate}
\item Construis les échiquiers de croisement des deux couples.
\item Montre qu'un échange d'enfants a bien eu lieu, en justifiant par les génotypes.
\item Rédige ta conclusion en trois temps (résultat, interprétation, généralisation).
\end{enumerate}
\end{exercice}

\begin{corrige}[Problème de synthèse final]
\begin{enumerate}
\item \textbf{Couple Mballa} : père $O/O$ $\times$ mère $A/O$. Gamètes du père : $O$ ; gamètes de la mère : $A$ ou $O$. Enfants possibles : $A/O$ [A] ou $O/O$ [O]. \textbf{Couple Fouda} : père $B/O$ $\times$ mère $A/A$. Gamètes du père : $B$ ou $O$ ; gamètes de la mère : $A$. Enfants possibles : $A/B$ [AB] ou $A/O$ [A].
\item L'enfant [AB] attribué aux Mballa possède les allèles $A$ et $B$ : or ni le père $O/O$ ni la mère $A/O$ ne possèdent l'allèle $B$. Il ne peut donc pas être leur enfant ; en revanche, il est compatible avec le couple Fouda ($A/B$ figure dans leur échiquier). L'enfant [O] attribué aux Fouda est $O/O$ : la mère $A/A$ ne possède pas d'allèle $O$, il ne peut donc pas être leur enfant ; il est compatible avec le couple Mballa ($O/O$ figure dans leur échiquier). L'échange est donc démontré.
\item \textbf{Conclusion rédigée :} « On observe que l'enfant [AB] est incompatible avec le couple Mballa mais compatible avec le couple Fouda, et inversement pour l'enfant [O]. Cela signifie que les allèles transmis par chaque couple ne peuvent produire que les génotypes prévus par les échiquiers de croisement. Donc, de façon générale, le groupage sanguin permet d'exclure une filiation lorsque les génotypes des parents ne peuvent pas produire le génotype de l'enfant. »
\end{enumerate}
\end{corrige}

\begin{aretenir}[L'essentiel de la leçon]
L'information génétique est portée par les gènes, situés sur les 46 chromosomes de l'espèce humaine. Chaque gène existe en plusieurs versions appelées allèles, dont la combinaison (génotype) détermine le caractère observable (phénotype). L'exemple du groupe sanguin ABO (allèles $A$ et $B$ codominants, $O$ récessif) et du facteur Rhésus illustre cette expression de l'information génétique et explique une partie de la diversité humaine, avec des applications concrètes : transfusions sûres, prévention de la maladie hémolytique du nouveau-né et exclusion de paternité.
\end{aretenir}

\newpage

\coursec{Activité d'intégration}

\courssub{Objectifs de l'activité}
Cette activité d'intégration mobilise l'ensemble des savoirs et savoir-faire de la leçon « Expression de l'information génétique » pour résoudre une situation complexe, authentique et contextualisée. À l'issue de cette activité, tu dois être capable de :
\begin{itemize}
    \item \textbf{Identifier} le groupe sanguin (système ABO et facteur Rhésus) à partir de l'observation de plaques de groupage (réactions d'agglutination).
    \item \textbf{Déterminer} les compatibilités transfusionnelles en justifiant par la relation antigènes (sur les globules rouges) / anticorps (dans le plasma).
    \item \textbf{Analyser} un caryotype humain : déterminer le sexe, le nombre de chromosomes et identifier une éventuelle anomalie du nombre de chromosomes.
    \item \textbf{Vérifier} une filiation biologique par la construction d'un échiquier de croisement à partir des groupes sanguins des parents et de l'enfant.
    \item \textbf{Rédiger} un rapport médical structuré, argumenté et utilisant un vocabulaire scientifique précis.
\end{itemize}

\courssub{Situation : Urgence au Centre Hospitalier d'Essos (Yaoundé)}
\emph{Contexte.} Un grave accident de la circulation vient de se produire sur l'axe routier Yaoundé--Douala, au niveau du pont sur la Sanaga. Trois victimes sont évacuées en urgence vers le Centre Hospitalier d'Essos. Elles présentent une anémie hémorragique sévère nécessitant une transfusion sanguine immédiate. La banque de sang de l'hôpital dispose de poches de sang prélevées sur cinq volontaires (D1 à D5) ce matin même. Le biologiste de l'hôpital te transmet les résultats des tests de groupage sanguin (plaques de groupage) et d'autres analyses complémentaires.

\textbf{Document 1 : Résultats des plaques de groupage (hémagglutination).}\par
Le test de groupage utilise trois sérums de référence : anti-A, anti-B et anti-D (Rhésus). L'agglutination (formation d'amas visibles de globules rouges) est notée « + », son absence « – ».

\begin{center}
\begin{tabularx}{\linewidth}{|c|c|c|c|X|}\hline
\rowcolor{popblueL}\textbf{Individu} & \textbf{Anti-A} & \textbf{Anti-B} & \textbf{Anti-D} & \textbf{Phénotype à déterminer} \\ \hline
Victime 1 (V1) & + & – & + & \ldots \\ \hline
Victime 2 (V2) & – & + & – & \ldots \\ \hline
Victime 3 (V3) & – & – & + & \ldots \\ \hline
Donneur 1 (D1) & + & – & + & \ldots \\ \hline
Donneur 2 (D2) & – & + & + & \ldots \\ \hline
Donneur 3 (D3) & + & + & – & \ldots \\ \hline
Donneur 4 (D4) & – & – & – & \ldots \\ \hline
Donneur 5 (D5) & + & – & – & \ldots \\ \hline
\end{tabularx}
\end{center}
\legende{Résultats des tests d'hémagglutination pour les 3 victimes et les 5 donneurs potentiels.}

\textbf{Document 2 : Caryotype de la Victime 1 (V1).}\par
La victime 1, inconsciente, porte une carte d'identité indiquant un sexe masculin. Le laboratoire de l'hôpital a réalisé son caryotype à partir d'une prise de sang. Le cliché montre \textbf{47 chromosomes}. On observe la présence de \textbf{trois chromosomes homologues au niveau de la paire 21} (petits chromosomes). Les chromosomes sexuels (gonosomes) sont X et Y.

\textbf{Document 3 : Vérification de filiation pour la Victime 2 (V2).}\par
La victime 2, un adolescent conscient mais choqué, affirme être le fils biologique de M. \textsc{Mvondo} (père) et de Mme \textsc{Mvondo} (mère), présents à son chevet. Le père est de groupe \textbf{B Rhésus +} (phénotype B+), la mère de groupe \textbf{A Rhésus +} (phénotype A+). La victime 2 (V2) a été typée (voir Document 1). Le père émet un doute sur sa paternité biologique car il trouve que l'enfant ne lui ressemble pas. L'équipe médicale demande un avis génétique rapide basé sur les groupes sanguins.

\courssub{Consignes}
\emph{Traite les questions suivantes en justifiant rigoureusement chaque réponse.}
\begin{enumerate}
    \item \textbf{Groupage sanguin :} À partir du Document 1, détermine le phénotype complet (système ABO et Rhésus) de chaque victime (V1, V2, V3) et de chaque donneur (D1 à D5). Rappelle la règle de correspondance entre agglutination et présence d'antigène.
    \item \textbf{Compatibilité transfusionnelle :} Pour chaque victime, dresse la liste des donneurs compatibles parmi D1 à D5. Justifie chaque compatibilité ou incompatibilité en précisant les antigènes du donneur et les anticorps du receveur (règle du receveur).
    \item \textbf{Analyse du caryotype (Document 2) :}
    \begin{enumerate}
        \item Écris la formule chromosomique (formule caryotypique) complète de la Victime 1.
        \item Précise le sexe chromosomique et le nombre total de chromosomes.
        \item Nomme l'anomalie chromosomique identifiée. Quelle est la maladie associée ? Cette anomalie est-elle héréditaire au sens d'un caractère transmis par les gènes des parents ? Justifie.
    \end{enumerate}
    \item \textbf{Étude de la filiation (Document 3) :}
    \begin{enumerate}
        \item Déduis les génotypes probables (allèles ABO) des parents M. et Mme \textsc{Mvondo} en tenant compte de leur phénotype et de celui de l'enfant (V2).
        \item Construis l'échiquier de croisement correspondant à l'union de ces deux parents pour le système ABO.
        \item Conclus sur la compatibilité du groupe sanguin de V2 avec la paternité de M. \textsc{Mvondo}. La filiation est-elle exclue ou possible ? Précise la probabilité d'avoir un enfant de groupe B pour ce couple.
    \end{enumerate}
    \item \textbf{Synthèse :} Rédige le \textbf{rapport médical d'urgence} adressé au médecin chef de service. Ce rapport doit comporter : l'identification des victimes, les groupes sanguins déterminés, le tableau des transfusions réalisables (Victime $\leftarrow$ Donneur), le diagnostic chromosomique pour V1, l'avis sur la filiation pour V2, et les signatures.
\end{enumerate}

\courssub{Ressources à mobiliser}
\begin{itemize}
    \item \textbf{Système ABO :} Trois allèles $I^A$, $I^B$, $i$ (ou $I^O$). $I^A$ et $I^B$ sont codominants entre eux et dominants sur $i$. Antigènes A et B à la surface des globules rouges (hématies). Anticorps anti-A et anti-B dans le plasma (règle de Landsteiner : on possède les anticorps dirigés contre les antigènes que l'on \emph{ne possède pas}).
    \item \textbf{Système Rhésus (facteur D) :} Allèle $D$ (Rh+) dominant sur $d$ (Rh–). Antigène D à la surface des hématies. Il n'existe pas d'anticorps anti-D naturels : ils n'apparaissent qu'après un contact avec du sang Rh+ (immunisation).
    \item \textbf{Compatibilité ABO (règle du receveur) :}
    \begin{center}
    \begin{tabular}{|c|c|c|}\hline
    \rowcolor{popblueL}\textbf{Groupe du receveur} & \textbf{Anticorps du receveur} & \textbf{Donneurs compatibles (globules rouges)} \\ \hline
    A & anti-B & A, O \\ \hline
    B & anti-A & B, O \\ \hline
    AB & aucun & A, B, AB, O (receveur universel) \\ \hline
    O & anti-A et anti-B & O seulement (O = donneur universel) \\ \hline
    \end{tabular}
    \end{center}
    \item \textbf{Compatibilité Rhésus :} Receveur Rh+ $\leftarrow$ donneur Rh+ ou Rh–. Receveur Rh– $\leftarrow$ donneur Rh– uniquement.
    \item \textbf{Caryotype humain normal :} 46 chromosomes = 22 paires de chromosomes non sexuels (autosomes) + 1 paire de chromosomes sexuels (XX chez la femme, XY chez l'homme). Formule : 46,XX ou 46,XY.
    \item \textbf{Anomalies du nombre de chromosomes :} trisomie (47 chromosomes, 3 exemplaires d'une paire), monosomie (45 chromosomes, 1 seul exemplaire). Exemples : trisomie 21, monosomie X (45,X).
    \item \textbf{Croisement ABO :} utilisation de l'échiquier de croisement (gamètes sur les bords, zygotes dans les cases).
\end{itemize}

\courssub{Démarche et corrigé commenté}
\begin{exoresolu}[Corrigé détaillé de l'activité d'intégration : Urgence à l'hôpital]
\textbf{Question 1 : Détermination des phénotypes (Document 1).}\par
\emph{Principe :} l'agglutination (+) avec un sérum anti-X signale la présence de l'antigène X à la surface des hématies. L'absence d'agglutination (–) signale l'absence de cet antigène. Ainsi : agglutination avec anti-A seul $\rightarrow$ groupe A ; avec anti-B seul $\rightarrow$ groupe B ; avec anti-A et anti-B $\rightarrow$ groupe AB ; avec aucun des deux $\rightarrow$ groupe O. L'agglutination avec anti-D indique le Rhésus positif (Rh+).

\begin{center}
\begin{tabularx}{\linewidth}{|c|c|c|c|X|}\hline
\rowcolor{popblueL}\textbf{Individu} & \textbf{Anti-A} & \textbf{Anti-B} & \textbf{Anti-D} & \textbf{Phénotype (ABO / Rhésus)} \\ \hline
Victime 1 (V1) & + & – & + & \textbf{A Rhésus + (A+)} \\ \hline
Victime 2 (V2) & – & + & – & \textbf{B Rhésus – (B–)} \\ \hline
Victime 3 (V3) & – & – & + & \textbf{O Rhésus + (O+)} \\ \hline
Donneur 1 (D1) & + & – & + & \textbf{A+} \\ \hline
Donneur 2 (D2) & – & + & + & \textbf{B+} \\ \hline
Donneur 3 (D3) & + & + & – & \textbf{AB–} \\ \hline
Donneur 4 (D4) & – & – & – & \textbf{O–} \\ \hline
Donneur 5 (D5) & + & – & – & \textbf{A–} \\ \hline
\end{tabularx}
\end{center}
\legende{Tableau récapitulatif des phénotypes déterminés.}

\textbf{Question 2 : Compatibilité transfusionnelle.}\par
\emph{Rappel :} on transfuse des globules rouges. La compatibilité dépend des antigènes du \textbf{donneur} (sur l'hématie) et des anticorps du \textbf{receveur} (dans le plasma). Il ne doit y avoir \textbf{aucune réaction antigène-anticorps}, sinon les globules rouges transfusés sont détruits (agglutination puis destruction des hématies : accident transfusionnel).

\begin{itemize}
    \item \textbf{V1 (A+) :} anticorps du receveur = anti-B. Elle reçoit du sang A ou O. Étant Rh+, elle reçoit du sang Rh+ ou Rh–.
    \begin{itemize}
        \item D1 (A+) : \textbf{compatible} (antigène A seul, pas d'antigène B ; Rh+ accepté).
        \item D2 (B+) : \textbf{incompatible} (l'antigène B est reconnu par les anticorps anti-B du receveur $\rightarrow$ agglutination).
        \item D3 (AB–) : \textbf{incompatible} (antigène B présent).
        \item D4 (O–) : \textbf{compatible} (ni antigène A ni antigène B ; le sang Rh– convient à un receveur Rh+).
        \item D5 (A–) : \textbf{compatible} (antigène A seul ; Rh– accepté).
    \end{itemize}
    \cle{Donneurs pour V1 : D1, D4, D5.}

    \item \textbf{V2 (B–) :} anticorps du receveur = anti-A. Il reçoit du sang B ou O. Étant Rh–, il ne reçoit que du sang Rh–.
    \begin{itemize}
        \item D1 (A+) : \textbf{incompatible} (antigène A reconnu par les anti-A ; de plus Rh+).
        \item D2 (B+) : \textbf{incompatible} (sang Rh+ transfusé à un receveur Rh– : risque de fabrication d'anticorps anti-D par le receveur).
        \item D3 (AB–) : \textbf{incompatible} (antigène A présent).
        \item D4 (O–) : \textbf{compatible} (ni antigène A ni antigène B ; Rh–).
        \item D5 (A–) : \textbf{incompatible} (antigène A présent).
    \end{itemize}
    \cle{Donneur pour V2 : D4 uniquement.}

    \item \textbf{V3 (O+) :} anticorps du receveur = anti-A et anti-B. Elle reçoit du sang O uniquement. Étant Rh+, elle reçoit du sang Rh+ ou Rh–.
    \begin{itemize}
        \item D1 (A+) : \textbf{incompatible} (antigène A).
        \item D2 (B+) : \textbf{incompatible} (antigène B).
        \item D3 (AB–) : \textbf{incompatible} (antigènes A et B).
        \item D4 (O–) : \textbf{compatible} (ni antigène A ni antigène B ; Rh– accepté).
        \item D5 (A–) : \textbf{incompatible} (antigène A).
    \end{itemize}
    \cle{Donneur pour V3 : D4 uniquement.}
\end{itemize}

\textbf{Question 3 : Analyse du caryotype de V1 (Document 2).}
\begin{enumerate}
    \item \textbf{Formule chromosomique :} \textbf{47,XY,+21} (47 chromosomes, gonosomes XY, un chromosome 21 en surnombre).
    \item \textbf{Sexe chromosomique :} masculin (présence des gonosomes X et Y). \textbf{Nombre total :} 47 chromosomes au lieu de 46.
    \item \textbf{Anomalie :} \textbf{trisomie 21} : il existe trois chromosomes 21 au lieu de deux.
    \item \textbf{Maladie associée :} la \textbf{trisomie 21} (ou syndrome de Down). Elle se traduit par un retard mental plus ou moins marqué, des traits du visage caractéristiques et parfois des malformations du c\oe ur.
    \item \textbf{Hérédité :} cette anomalie n'est \textbf{pas héréditaire au sens d'un caractère transmis par les gènes des parents} : les parents de V1 possèdent un caryotype normal (46 chromosomes). L'anomalie est apparue \emph{accidentellement} lors de la formation des gamètes de l'un des parents (le plus souvent de la mère) : un gamète a reçu deux chromosomes 21 au lieu d'un. Le risque de cette anomalie augmente avec l'âge de la mère.
\end{enumerate}

\textbf{Question 4 : Étude de la filiation pour V2 (Document 3).}
\begin{enumerate}
    \item \textbf{Génotypes parentaux :}
    \begin{itemize}
        \item Père (M. Mvondo), phénotype B+ : génotype ABO $I^B/i$ (hétérozygote) \emph{ou} $I^B/I^B$ (homozygote) ; génotype Rhésus $D/d$ ou $D/D$.
        \item Mère (Mme Mvondo), phénotype A+ : génotype ABO $I^A/i$ \emph{ou} $I^A/I^A$ ; génotype Rhésus $D/d$ ou $D/D$.
        \item Enfant V2, phénotype B– : génotype ABO \textbf{$I^B/i$} et génotype Rhésus \textbf{$d/d$} (le phénotype Rh– est récessif, il impose deux allèles $d$).
    \end{itemize}
    \emph{Déduction cruciale :} V2 est Rh– ($d/d$) : il a reçu un allèle $d$ de chaque parent. Les deux parents, bien que Rh+, sont donc \textbf{obligatoirement hétérozygotes $D/d$}. Pour le système ABO, V2 ($I^B/i$) a reçu $I^B$ de son père (groupe B) et l'allèle $i$ de sa mère : la mère, bien que de groupe A, est donc \textbf{obligatoirement hétérozygote $I^A/i$}.

    \item \textbf{Échiquier de croisement (système ABO) :}\par
    \emph{Parents :} mère $I^A/i$ $\times$ père $I^B/i$ (cas le plus informatif ; si le père était $I^B/I^B$, la conclusion sur la possibilité de la filiation serait la même).
    \begin{center}
    \begin{tabular}{|c|c|c|}\hline
    \rowcolor{popblueL}\textbf{Gamètes} & $I^B$ (père) & $i$ (père) \\ \hline
    $I^A$ (mère) & $I^A/I^B$ : groupe \textbf{AB} & $I^A/i$ : groupe \textbf{A} \\ \hline
    $i$ (mère) & $I^B/i$ : groupe \textbf{B} & $i/i$ : groupe \textbf{O} \\ \hline
    \end{tabular}
    \end{center}
    \legende{Échiquier de croisement ABO : mère $I^A/i$ $\times$ père $I^B/i$. Les quatre groupes A, B, AB et O sont possibles dans la descendance.}

    \item \textbf{Conclusion sur la filiation :}\par
    Les phénotypes possibles pour les enfants de ce couple sont : \textbf{A, B, AB et O}. La victime 2 est de groupe \textbf{B} : ce phénotype figure bien dans la descendance théorique (génotype $I^B/i$). De même, son phénotype Rh– est possible puisque deux parents Rh+ hétérozygotes $D/d$ peuvent avoir un enfant $d/d$ (probabilité 1/4).
    \cle{La filiation n'est PAS exclue par les groupes sanguins : elle est génétiquement possible.}
    \emph{Probabilité :} pour ce couple ($I^A/i \times I^B/i$), la probabilité d'avoir un enfant de groupe B est de 1/4, soit 25\,\%. (Si le père était $I^B/I^B$, elle serait de 1/2, soit 50\,\%.)
    \begin{attention}[Limite du test de groupe sanguin]
    Les groupes sanguins permettent d'\textbf{exclure} une paternité (quand le groupe de l'enfant est impossible avec les groupes des parents), mais jamais de la \textbf{prouver} avec certitude : beaucoup d'hommes de groupe B pourraient être le père de V2. Seule une analyse de l'ADN permettrait une conclusion définitive.
    \end{attention}
\end{enumerate}

\textbf{Question 5 : Rapport médical d'urgence (Synthèse).}

\begin{center}
\begin{tabularx}{\linewidth}{|X|}\hline
\textbf{RÉPUBLIQUE DU CAMEROUN} \\
\textbf{Paix -- Travail -- Patrie} \\ \hline
\textbf{CENTRE HOSPITALIER D'ESSOS --- SERVICE DES URGENCES} \\
\textbf{LABORATOIRE DE BIOLOGIE MÉDICALE} \\ \hline
\textbf{RAPPORT D'URGENCE TRANSFUSIONNELLE ET GÉNÉTIQUE} \\
\textbf{Référence :} ACC/2026/045/Yaoundé \hfill \textbf{Date :} 15 octobre 2026 \\ \hline
\textbf{Objet :} prise en charge transfusionnelle de 3 victimes d'un accident sur l'axe Yaoundé--Douala ; avis sur caryotype et filiation. \\ \hline
\end{tabularx}
\end{center}

\vspace{0.3cm}
\noindent \textbf{1. Identification des patients :}
\begin{itemize}
    \item Victime 1 (V1) : sexe masculin, \textbf{groupe A Rhésus + (A+)}. Caryotype : 47,XY,+21 (trisomie 21).
    \item Victime 2 (V2) : sexe masculin, \textbf{groupe B Rhésus – (B–)}.
    \item Victime 3 (V3) : \textbf{groupe O Rhésus + (O+)}.
\end{itemize}

\noindent \textbf{2. Bilan de compatibilité transfusionnelle (globules rouges) :}
\begin{center}
\begin{tabularx}{\linewidth}{|c|c|X|X|}\hline
\rowcolor{popblueL}\textbf{Receveur} & \textbf{Groupe} & \textbf{Donneurs compatibles} & \textbf{Commentaire} \\ \hline
V1 & A+ & \textbf{D1 (A+), D4 (O–), D5 (A–)} & 3 poches disponibles. Priorité à D1 (même groupe). \\ \hline
V2 & B– & \textbf{D4 (O–) uniquement} & \textbf{Stock critique :} une seule poche compatible. Rechercher en urgence des donneurs B– ou O–. \\ \hline
V3 & O+ & \textbf{D4 (O–) uniquement} & \textbf{Conflit :} D4 est le seul donneur compatible pour V2 ET pour V3. \\ \hline
\end{tabularx}
\end{center}
\legende{Tableau de décision transfusionnelle.}

\noindent \textbf{3. Avis sur le caryotype (Victime 1) :}
Le caryotype 47,XY,+21 confirme une \textbf{trisomie 21}. Cette anomalie du nombre de chromosomes est apparue accidentellement lors de la formation des gamètes : elle n'est pas transmise par les gènes des parents. Un bilan cardiaque et une prise en charge adaptée sont recommandés ; les parents doivent être informés.

\noindent \textbf{4. Avis sur la filiation (Victime 2) :}
Analyse des groupes sanguins ABO/Rhésus du couple parental (père B+, mère A+) et de l'enfant (V2 : B–).
\begin{itemize}
    \item L'enfant est Rh– ($d/d$) : compatible avec deux parents Rh+ hétérozygotes ($D/d \times D/d$), probabilité 1/4.
    \item L'enfant est de groupe B ($I^B/i$) : compatible avec une mère A ($I^A/i$) et un père B ($I^B/i$ ou $I^B/I^B$), probabilité 1/4 à 1/2.
    \item \textbf{Conclusion :} \textbf{aucune exclusion de paternité} par les systèmes ABO et Rhésus. La filiation biologique est \textbf{possible}. Pour une certitude, seule une analyse de l'ADN pourrait trancher définitivement.
\end{itemize}

\noindent \textbf{5. Recommandations immédiates :}
\begin{enumerate}
    \item Transfuser V1 avec D1 (A+) en priorité.
    \item Transfuser V2 (B–) avec l'unique poche D4 (O–) \textbf{immédiatement} (pronostic vital engagé).
    \item Pour V3 (O+) : contacter le Centre National de Transfusion Sanguine (CNTS) de Yaoundé pour l'acheminement en urgence de poches O+ ou O–, ainsi que de poches B– ou O– pour V2.
    \item Signaler la trisomie 21 (V1) au service de pédiatrie pour suivi.
\end{enumerate}

\vspace{0.8cm}
\noindent \textbf{Fait à Yaoundé, le 15/10/2026} \\
\noindent \textbf{Le Biologiste de garde} \hfill \textbf{Le Médecin chef de service} \\
\noindent \emph{(Signature et cachet)} \hfill \emph{(Signature et cachet)}
\end{exoresolu}

\courssub{Bilan de l'activité}
Cette activité d'intégration a permis de mettre en évidence la \textbf{transversalité des concepts génétiques} abordés dans l'unité :
\begin{itemize}
    \item \textbf{Du chromosome au phénotype :} l'analyse du caryotype (niveau chromosomique) a révélé une anomalie du nombre de chromosomes (trisomie 21) chez V1, indépendante de son groupe sanguin (niveau des gènes et des allèles) : l'information génétique s'exprime à plusieurs échelles.
    \item \textbf{Génotype et phénotype :} la détermination des groupes sanguins a nécessité de passer du phénotype observé (agglutination) au génotype déduit (allèles $I^A$, $I^B$, $i$, $D$, $d$) pour résoudre le problème de filiation.
    \item \textbf{Application médicale concrète :} la gestion des poches de sang (ressource rare) a imposé un raisonnement rigoureux sur la compatibilité antigène/anticorps (règle du receveur) et la gestion des priorités (conflit autour du donneur universel D4).
    \item \textbf{Argumentation scientifique :} la rédaction du rapport médical a exigé une structuration logique (données $\rightarrow$ analyse $\rightarrow$ conclusion $\rightarrow$ décision), un vocabulaire précis (caryotype, allèle, phénotype, agglutination, trisomie) et une attitude responsable (information sur la trisomie, réserve sur la preuve de paternité).
\end{itemize}
\cle{Compétence clé validée :} « Mobiliser les connaissances sur l'organisation et l'expression de l'information génétique pour expliquer la diversité humaine et résoudre des problèmes de santé (transfusion sanguine, diagnostic d'anomalie chromosomique, filiation). »

\newpage

\coursec{Méthodes et exercices résolus}

\courssub{Méthodes et exercices résolus}

\begin{methode}[Réaliser une maquette de caryotype humain (TP N°3)]
    \textbf{Objectif :} Passer d'une photographie de chromosomes en métaphase (étalement chromosomique) à un caryotype ordonné pour identifier le sexe et d'éventuelles anomalies numériques ou structurales.
    \begin{enumerate}
        \item \textbf{Préparer le matériel :} Photographie d'un étalement chromosomique (ou jeu de chromosomes découpés), feuille de papier format A3, ciseaux, colle, règle, crayon à papier.
        \item \textbf{Découper les chromosomes :} Découper soigneusement \emph{chaque} chromosome visible sur la photographie en laissant une marge fine autour du chromatine. Compter le nombre total de chromosomes découpés (normalement \textbf{46}).
        \item \textbf{Classer par paires homologues :} Rassembler les chromosomes deux par deux selon trois critères morphologiques : 
        \begin{itemize}
            \item la \textbf{longueur} (du plus long au plus court) ;
            \item la \textbf{position du centromère} (métacentrique $\approx$ médian, sous-métacentrique, acrocentrique $\approx$ terminal) ;
            \item le \textbf{bandage} (alternance de bandes claires et sombres après coloration G, R ou Q) — \emph{critère de confirmation}.
        \end{itemize}
        \item \textbf{Ordonner les paires :} Disposer les 22 paires d'\textbf{autosomes} par ordre décroissant de taille (paires 1 à 22) en 7 groupes (A à G) selon la classification internationale (Denver/Paris). Placer la paire de \textbf{gonosomes} (chromosomes sexuels) à la fin (paire 23).
        \item \textbf{Coller et légender :} Coller les paires alignées par leur centromère sur une ligne imaginaire horizontale (bras courts \textbf{p} vers le haut, bras longs \textbf{q} vers le bas). Indiquer le numéro de chaque paire. Écrire la formule caryotypique : \textbf{46,XX} (femme) ou \textbf{46,XY} (homme).
        \item \textbf{Analyser :} Vérifier le nombre total (46 = normal ; 47 = trisomie ; 45 = monosomie) et la structure (délétion, translocation, inversion visibles si bande anormale).
    \end{enumerate}
    \textbf{Réflexe :} Toujours aligner les centromères sur une même ligne horizontale pour faciliter la comparaison des bras courts (p) et longs (q) et détecter toute anomalie structurale.
\end{methode}

\begin{exoresolu}[Analyse d'un caryotype humain — Détection d'une trisomie 21]
Sur une photographie d'étalement chromosomique obtenu à partir de lymphocytes fœtaux (amniocentèse), un élève compte 47 chromosomes après découpage. Il classe les chromosomes et obtient l'arrangement suivant pour les trois dernières paires :
\begin{itemize}
    \item Paire 21 : \textbf{trois} chromosomes acrocentriques de taille très petite, au bandage identique.
    \item Paire 22 : deux chromosomes acrocentriques de taille très petite.
    \item Gonosomes : deux chromosomes X.
\end{itemize}
\textbf{Questions :}
\begin{enumerate}
    \item Établir la formule caryotypique complète de cet individu.
    \item Préciser le sexe chromosomique.
    \item Nommer l'anomalie chromosomique identifiée et en donner la conséquence phénotypique majeure.
    \item Expliquer l'origine cellulaire de cette anomalie (mécanisme au cours de la méiose).
\end{enumerate}
\tcblower
\textbf{Solution détaillée :}
\begin{enumerate}
    \item \textbf{Formule caryotypique :} Le comptage total donne 47 chromosomes. Les autosomes sont au nombre de 44 (22 paires) + 1 chromosome en excès sur la paire 21. Les gonosomes sont XX. La formule s'écrit : \textbf{47,XX,+21}. (Notation : nombre total, gonosomes, anomalie avec signe + pour gain sur la paire concernée).
    \item \textbf{Sexe chromosomique :} La présence de deux chromosomes X (XX) indique un sexe chromosomique \textbf{féminin}.
    \item \textbf{Anomalie et conséquence :} Il s'agit d'une \textbf{trisomie libre de la paire 21} (ou trisomie 21 standard). C'est l'anomalie à l'origine du \textbf{syndrome de Down}. Conséquences phénotypiques majeures : retard mental modéré, traits faciaux caractéristiques (fente palpébrale oblique, épicanthus, langue proéminente), hypotonie musculaire, malformations cardiaques fréquentes (communication inter-ventriculaire), risque accru de leucémie et de maladie d'Alzheimer précoce.
    \item \textbf{Origine cellulaire (Non-disjonction) :} Cette anomalie résulte d'une \textbf{non-disjonction} survenue au cours de la \textbf{méiose} lors de la gamétogenèse (le plus souvent en méiose I chez la mère, plus rarement en méiose II ou chez le père). 
    \begin{itemize}
        \item \textbf{Méiose I normale :} Les deux chromosomes 21 homologues se séparent $\rightarrow$ chaque gamète reçoit un chromosome 21 (formé de deux chromatides).
        \item \textbf{Non-disjonction en méiose I :} Les deux chromosomes 21 homologues ne se séparent pas $\rightarrow$ un gamète reçoit \textbf{deux} chromosomes 21 (disomique 21), l'autre \textbf{zéro} (nullique 21).
        \item \textbf{Fécondation :} Si le gamète disomique (n+1 = 24 chromosomes) est fécondé par un gamète normal (n = 23 chromosomes), le zygote possède 47 chromosomes dont trois chromosomes 21.
    \end{itemize}
    \textbf{Conclusion :} Le caryotype 47,XX,+21 confirme une trisomie 21 d'origine méiotique (non-disjonction), responsable du syndrome de Down chez une fille.
\end{enumerate}
\end{exoresolu}

\begin{methode}[Déterminer le groupe sanguin (système ABO et Rhésus) par test d'agglutination (TP N°4)]
    \textbf{Objectif :} Identifier le phénotype sanguin d'un individu à l'aide de sérums de référence (Anti-A, Anti-B, Anti-AB, Anti-D) en observant la réaction d'agglutination (grumeaux) ou de non-agglutination.
    \begin{enumerate}
        \item \textbf{Préparer le matériel :} Sang total (prélèvement au doigt ou tube EDTA), 3 à 4 lames de verre ou plaque de test, sérums de test : \textbf{Anti-A}, \textbf{Anti-B}, \textbf{Anti-AB} (contrôle), \textbf{Anti-D} (Anti-Rhésus), baguettes de mélange en plastique, chronomètre.
        \item \textbf{Réaliser le test (technique sur lame) :}
        \begin{itemize}
            \item Déposer une goutte de chaque sérum sur des zones distinctes de la lame (étiqueter : A, B, AB, D).
            \item Ajouter une goutte de sang \emph{à côté} de chaque goutte de sérum (pas dedans pour éviter la dilution avant mélange).
            \item Mélanger soigneusement sang + sérum sur une surface d'environ 2 cm² à l'aide d'une baguette \textbf{différente pour chaque sérum} (éviter la contamination croisée).
        \end{itemize}
        \item \textbf{Observer et chronométrer :} Incliner doucement la lame (mouvements de rotation lents) pendant \textbf{2 à 3 minutes}. Observer à l'œil nu (ou loupe $\times$10) la formation de \textbf{grumeaux (agglutinats)} rouges au centre d'un fond clair (surnageant).
        \item \textbf{Noter les résultats :} 
        \begin{center}
            \begin{tabular}{|c|c|c|}
                \hline
                \textbf{Sérum testé} & \textbf{Agglutination (+)} & \textbf{Pas d'agglutination (-)} \\ \hline
                Anti-A & Ag A présent & Ag A absent \\ \hline
                Anti-B & Ag B présent & Ag B absent \\ \hline
                Anti-AB & Ag A \textbf{ou} B présent & Ag A \textbf{et} B absents \\ \hline
                Anti-D (Rh) & Ag D présent (Rh+) & Ag D absent (Rh-) \\ \hline
            \end{tabular}
        \end{center}
        \item \textbf{Déduire le phénotype ABO et Rhésus :}
        \begin{center}
            \begin{tabular}{|c|c|c|c|}
                \hline
                \textbf{Résultats (A / B / AB)} & \textbf{Phénotype ABO} & \textbf{Génotypes possibles} \\ \hline
                + / - / + & A (ou A1) & $I^A I^A$ ou $I^A i$ \\ \hline
                - / + / + & B & $I^B I^B$ ou $I^B i$ \\ \hline
                + / + / + & AB & $I^A I^B$ \\ \hline
                - / - / - & O & $ii$ \\ \hline
            \end{tabular}
        \end{center}
        Phénotype Rhésus : Agglutination Anti-D (+) $\rightarrow$ \textbf{Rh+} ($DD$ ou $Dd$) ; Pas d'agglutination (-) $\rightarrow$ \textbf{Rh-} ($dd$).
        \item \textbf{Contrôles obligatoires :} Témoin auto (hématies du patient + son propre sérum = -), hématies témoins A1, B, O connues + sérums tests (validation des réactifs).
    \end{enumerate}
    \textbf{Réflexe :} L'agglutination = réaction Ag/Ac visible. Anti-AB est un contrôle : s'il est (-) alors que A ou B est (+), le test est invalide (réactif périmé ou erreur de manipulation).
\end{methode}

\begin{exoresolu}[Interprétation d'une fiche de groupage sanguin]
Lors d'une campagne de don de sang au Centre National de Transfusion Sanguine (CNTS) de Yaoundé, les résultats suivants sont obtenus pour un donneur :
\begin{center}
    \begin{tabular}{|c|c|}
        \hline
        \textbf{Sérum de test} & \textbf{Résultat (Agglutination)} \\ \hline
        Anti-A & \textbf{Négatif} (-) \\ \hline
        Anti-B & \textbf{Positif} (+) \\ \hline
        Anti-AB & \textbf{Positif} (+) \\ \hline
        Anti-D & \textbf{Négatif} (-) \\ \hline
    \end{tabular}
\end{center}
\textbf{Questions :}
\begin{enumerate}
    \item Déterminer le phénotype ABO et Rhésus complet de ce donneur.
    \item Préciser son ou ses génotype(s) possible(s) pour le système ABO et pour le système Rhésus.
    \item Ce donneur peut-il donner son sang à un receveur de phénotype \textbf{A+} ? Justifier en précisant la règle de compatibilité globulaire (don de globules rouges concentrés).
    \item Ce donneur peut-il recevoir du sang d'un donneur \textbf{O-} ? Justifier.
\end{enumerate}
\tcblower
\textbf{Solution détaillée :}
\begin{enumerate}
    \item \textbf{Phénotype :} 
    \begin{itemize}
        \item Anti-A (-) : Ag A absent. Anti-B (+) : Ag B présent. Anti-AB (+) : cohérent (Ag B présent). $\rightarrow$ \textbf{Groupe B}.
        \item Anti-D (-) : Ag D absent. $\rightarrow$ \textbf{Rhésus négatif (Rh-)}.
    \end{itemize}
    Phénotype complet : \textbf{B-} (ou B Rhésus négatif).
    \item \textbf{Génotypes :}
    \begin{itemize}
        \item Système ABO (allèles $I^A, I^B, i$ ; $I^A$ et $I^B$ codominants, dominants sur $i$) : Phénotype B $\rightarrow$ Génotypes possibles : \textbf{$I^B I^B$} (homozygote) \textbf{ou} \textbf{$I^B i$} (hétérozygote).
        \item Système Rhésus (allèles $D$ dominant sur $d$) : Phénotype Rh- ($dd$ seulement) $\rightarrow$ Génotype unique : \textbf{$dd$} (homozygote récessif).
    \end{itemize}
    \item \textbf{Don à un receveur A+ :} \textbf{NON.}
    \textbf{Justification :} Règle de compatibilité pour les \textbf{globules rouges (GR)} : \emph{Le donneur ne doit pas posséder d'antigènes que le receveur n'a pas (risque d'agglutination par les anticorps du receveur)}.
    \begin{itemize}
        \item Receveur A+ : Ag A et Ag D présents. Ac anti-B dans le plasma.
        \item Donneur B- : Ag B présent, Ag D absent.
        \item Problème : Le donneur possède l'antigène \textbf{B} que le receveur \textbf{n'a pas}. Les anti-B du receveur agglutineront les globules rouges du donneur $\rightarrow$ \textbf{Accident hémolytique transfusionnel grave}.
    \end{itemize}
    \item \textbf{Recevoir d'un donneur O- :} \textbf{OUI.}
    \textbf{Justification :} Règle de compatibilité pour le \textbf{receveur} : \emph{Il ne doit pas recevoir d'antigènes qu'il ne possède pas}.
    \begin{itemize}
        \item Donneur O- : Aucun antigène A, B, D sur les GR (donneur universel pour GR).
        \item Receveur B- : Possède Ag B. Pas d'Ag A, pas d'Ag D.
        \item Les GR du donneur O- n'apportent aucun antigène étranger (ni A, ni B, ni D). Pas de risque d'agglutination par les anticorps du receveur (anti-A, anti-D). Compatible.
    \end{itemize}
    \textbf{Conclusion :} Ce donneur B- peut donner aux groupes B et AB (Rh- ou Rh+ pour le Rh, mais idéalement Rh- pour une femme en âge de procréer) et recevoir des groupes B- et O-.
\end{enumerate}
\end{exoresolu}

\begin{methode}[Résoudre un problème de génétique formelle : Système ABO (Allèles multiples et Codominance)]
    \textbf{Objectif :} Prédire les phénotypes et génotypes possibles de la descendance à partir des phénotypes des parents pour le système ABO (3 allèles : $I^A, I^B, i$).
    \begin{enumerate}
        \item \textbf{Identifier les allèles et relations de dominance :} $I^A$ et $I^B$ sont \textbf{codominants} (s'expriment ensemble chez $I^A I^B$ $\rightarrow$ phénotype AB). $I^A$ et $I^B$ sont \textbf{dominants} sur $i$ (récessif). $ii$ $\rightarrow$ phénotype O.
        \item \textbf{Déduire les génotypes parentaux possibles :} 
        \begin{itemize}
            \item Phénotype A : $I^A I^A$ ou $I^A i$
            \item Phénotype B : $I^B I^B$ ou $I^B i$
            \item Phénotype AB : $I^A I^B$ (unique)
            \item Phénotype O : $ii$ (unique)
        \end{itemize}
        \item \textbf{Construire le(s) croisement(s) (Carré de Punnett) :} Pour chaque combinaison possible de génotypes parentaux, croiser les gamètes.
        \begin{itemize}
            \item Gamètes possibles : $I^A$, $I^B$, $i$.
            \item Remplir le tableau : lignes = gamètes mère, colonnes = gamètes père.
        \end{itemize}
        \item \textbf{Lister les génotypes et phénotypes de la descendance :} Compter les cases (total 4, 8 ou 16 selon homozygotie/hétérozygotie).
        \item \textbf{Calculer les fréquences (probabilités) :} $\frac{\text{Nombre de cases pour un phénotype}}{\text{Nombre total de cases}} \times 100\%$.
        \item \textbf{Rédiger la conclusion :} « Ce croisement peut donner des enfants de groupes [liste] aux fréquences respectives de [\%]. Le groupe [X] est impossible. »
    \end{enumerate}
    \textbf{Réflexe :} Toujours écrire les allèles en \emph{italique} avec exposant ($I^A, I^B, i$). Ne pas confondre $I$ (majuscule i) et $l$ (L minuscule). Vérifier la cohérence : somme des pourcentages = 100\%.
\end{methode}

\begin{exoresolu}[Croisement ABO : Parents de groupes A et B — Enfant possible de groupe O ?]
Un couple consulte pour un conseil génétique prénatal. Le père est de groupe sanguin \textbf{A}, la mère de groupe \textbf{B}. Leur premier enfant est de groupe \textbf{O}.
\textbf{Questions :}
\begin{enumerate}
    \item Déterminer avec certitude les génotypes du père et de la mère.
    \item Réaliser le croisement entre ces deux génotypes (Carré de Punnett).
    \item Calculer les probabilités (en \%) pour chaque groupe sanguin possible chez leurs enfants futurs.
    \item Un test de paternité exclut-il le père si un deuxième enfant est de groupe \textbf{AB} ? Justifier.
\end{enumerate}
\tcblower
\textbf{Solution détaillée :}
\begin{enumerate}
    \item \textbf{Génotypes parentaux déduits :}
    \begin{itemize}
        \item Enfant groupe O = phénotype $ii$. Il a reçu un allèle $i$ de \textbf{chaque} parent.
        \item Père groupe A : doit porter $i$ $\rightarrow$ Génotype \textbf{$I^A i$} (hétérozygote). ($I^A I^A$ impossible car ne donnerait que $I^A$).
        \item Mère groupe B : doit porter $i$ $\rightarrow$ Génotype \textbf{$I^B i$} (hétérozygote).
    \end{itemize}
    \item \textbf{Croisement : $I^A i \times I^B i$}
    \begin{itemize}
        \item Gamètes Père : $I^A$ ; $i$
        \item Gamètes Mère : $I^B$ ; $i$
    \end{itemize}
    \textbf{Carré de Punnett :}
    \begin{center}
        \begin{tabular}{|c|c|c|}
            \hline
            \textbf{Gamètes Mère $\downarrow$ / Père $\rightarrow$} & \textbf{$I^A$} & \textbf{$i$} \\ \hline
            \textbf{$I^B$} & $I^A I^B$ (\textbf{AB}) & $I^B i$ (\textbf{B}) \\ \hline
            \textbf{$i$} & $I^A i$ (\textbf{A}) & $ii$ (\textbf{O}) \\ \hline
        \end{tabular}
    \end{center}
    \item \textbf{Probabilités (4 cas équiprobables = 25\% chacun) :}
    \begin{itemize}
        \item Groupe A ($I^A i$) : \textbf{25\%}
        \item Groupe B ($I^B i$) : \textbf{25\%}
        \item Groupe AB ($I^A I^B$) : \textbf{25\%}
        \item Groupe O ($ii$) : \textbf{25\%}
    \end{itemize}
    \item \textbf{Test de paternité et enfant AB :} \textbf{NON, le test n'exclut pas le père.}
    \textbf{Justification :} Le croisement $I^A i \times I^B i$ prévoit explicitement la naissance d'enfants de génotype $I^A I^B$ (phénotype AB) avec une probabilité de 25\%. La naissance d'un enfant AB est \textbf{totalement compatible} avec la paternité de cet homme. Seuls les groupes impossibles (ex: un parent O et un parent AB ne peuvent pas faire d'enfant O) permettraient une exclusion. Ici, les 4 groupes sont possibles.
\end{enumerate}
\end{exoresolu}

\begin{methode}[Résoudre un problème de génétique formelle : Facteur Rhésus (Allèles Dominant/Récessif)]
    \textbf{Objectif :} Prédire le phénotype Rhésus de la descendance (Rh+ ou Rh-) à partir des phénotypes des parents (Allèles $D$ dominant, $d$ récessif).
    \begin{enumerate}
        \item \textbf{Rappel des génotypes :} Rh+ $\rightarrow$ $DD$ ou $Dd$ ; Rh- $\rightarrow$ $dd$ (uniquement).
        \item \textbf{Analyser les cas de figure parents :}
        \begin{itemize}
            \item \textbf{Rh- $\times$ Rh- :} $dd \times dd$ $\rightarrow$ 100\% Rh- ($dd$). \emph{Seul croisement donnant 100\% Rh-}.
            \item \textbf{Rh+ $\times$ Rh- :} 
                \begin{itemize}
                    \item Si Rh+ = $DD$ : $DD \times dd$ $\rightarrow$ 100\% $Dd$ (Rh+).
                    \item Si Rh+ = $Dd$ : $Dd \times dd$ $\rightarrow$ 50\% $Dd$ (Rh+) ; 50\% $dd$ (Rh-).
                \end{itemize}
            \item \textbf{Rh+ $\times$ Rh+ :} 
                \begin{itemize}
                    \item $DD \times DD$ : 100\% Rh+.
                    \item $DD \times Dd$ : 100\% Rh+.
                    \item $Dd \times Dd$ : 75\% Rh+ ($DD, Dd$) ; 25\% Rh- ($dd$). \emph{Seul cas donnant des Rh-}.
                \end{itemize}
        \end{itemize}
        \item \textbf{Utiliser l'information sur les enfants (rétro-croisement) :} Si un couple Rh+ a un enfant Rh- $\rightarrow$ Les \textbf{deux} parents sont obligatoirement hétérozygotes ($Dd$).
        \item \textbf{Calculer les probabilités demandées.}
    \end{enumerate}
    \textbf{Réflexe :} Ne jamais écrire « Génotype Rh+ » sans préciser $DD$ ou $Dd$. La mention « Rh+ » est un \textbf{phénotype}.
\end{methode}

\begin{exoresolu}[Transmission du Rhésus : Couple Rh+ ayant un enfant Rh-]
Au Cameroun, un couple dont les deux partenaires sont de phénotype \textbf{Rhésus positif (Rh+)} a déjà un enfant de phénotype \textbf{Rhésus négatif (Rh-)}. La mère est enceinte du deuxième enfant.
\textbf{Questions :}
\begin{enumerate}
    \item Déterminer avec certitude les génotypes des deux parents pour le facteur Rhésus.
    \item Réaliser le croisement correspondant.
    \item Calculer la probabilité (en \%) que le deuxième enfant soit Rhésus négatif.
    \item Calculer la probabilité que le deuxième enfant soit Rhésus positif \textbf{homozygote}.
    \item Si le père était Rh+ homozygote ($DD$) et la mère Rh- ($dd$), quelle aurait été la probabilité d'avoir un enfant Rh- ? Justifier.
\end{enumerate}
\tcblower
\textbf{Solution détaillée :}
\begin{enumerate}
    \item \textbf{Génotypes parentaux :} Un enfant Rh- a pour génotype \textbf{$dd$}. Il a reçu un allèle $d$ de chaque parent. Comme les deux parents sont phénotypiquement Rh+, ils portent obligatoirement l'allèle $D$. Ils sont donc tous les deux \textbf{hétérozygotes : $Dd$}.
    \item \textbf{Croisement : $Dd \times Dd$}
    \begin{itemize}
        \item Gamètes Père : $D$ ; $d$
        \item Gamètes Mère : $D$ ; $d$
    \end{itemize}
    \textbf{Carré de Punnett :}
    \begin{center}
        \begin{tabular}{|c|c|c|}
            \hline
            \textbf{Gamètes Mère $\downarrow$ / Père $\rightarrow$} & \textbf{$D$} & \textbf{$d$} \\ \hline
            \textbf{$D$} & $DD$ (Rh+) & $Dd$ (Rh+) \\ \hline
            \textbf{$d$} & $Dd$ (Rh+) & $dd$ (\textbf{Rh-}) \\ \hline
        \end{tabular}
    \end{center}
    \item \textbf{Probabilité Rh- :} 1 cas sur 4 ($dd$) $\rightarrow$ \textbf{25\%} (ou 1/4). \emph{Indépendant du premier enfant (événements indépendants).}
    \item \textbf{Probabilité Rh+ homozygote ($DD$) :} 1 cas sur 4 $\rightarrow$ \textbf{25\%}.
    \item \textbf{Cas hypothétique $DD \times dd$ :} 
    \begin{itemize}
        \item Gamètes Père : $D$ seulement. Gamètes Mère : $d$ seulement.
        \item Descendance : 100\% $Dd$ (Rh+ hétérozygote).
        \item Probabilité enfant Rh- ($dd$) = \textbf{0\%}.
    \end{itemize}
    \textbf{Justification :} Le père $DD$ ne possède pas l'allèle $d$ à transmettre. L'enfant reçoit obligatoirement $D$ du père et $d$ de la mère $\rightarrow$ $Dd$ (Rh+).
\end{enumerate}
\end{exoresolu}

\begin{methode}[Analyser un arbre généalogique (Pedigree) pour les groupes sanguins]
    \textbf{Objectif :} Déduire les génotypes des membres d'une famille à partir de leurs phénotypes (ABO et/ou Rhésus) et des relations de parenté, pour identifier le mode de transmission ou valider une paternité/maternité.
    \begin{enumerate}
        \item \textbf{Lire les symboles :} $\square$ = Homme ; $\bigcirc$ = Femme ; \textbf{Symboles pleins (noirs)} = Phénotype étudié présent (ex: Groupe A ou Rh-) ; \textbf{Symboles vides} = Phénotype absent (ex: Non-A ou Rh+). Trait horizontal = union (mariage) ; Trait vertical = descendance ; Fratrie = ordre de naissance (gauche = aîné).
        \item \textbf{Repérer les individus « informatifs » :} 
        \begin{itemize}
            \item Individus \textbf{homozygotes récessifs connus} (Phénotype O = $ii$ ; Phénotype Rh- = $dd$) : leurs génotypes sont \textbf{certains}.
            \item Individus ayant un parent homozygote récessif : ils ont \textbf{obligatoirement} reçu l'allèle récessif de ce parent.
        \end{itemize}
        \item \textbf{Propulser les allèles (Déduction descendante/ascendante) :} 
        \begin{itemize}
            \item Descendante : Parents $\rightarrow$ Enfants (vérifier compatibilité).
            \item Ascendante : Enfant $\rightarrow$ Parents (déduire allèle obligatoire).
        \end{itemize}
        \item \textbf{Tester les hypothèses pour les hétérozygotes :} Si un individu phénotype A (génotype $I^A I^A$ ou $I^A i$) a un enfant O ($ii$), il est \textbf{forcément} $I^A i$.
        \item \textbf{Conclure sur la compatibilité :} Si un enfant possède un allèle que \textbf{ni le père ni la mère} ne peuvent lui transmettre $\rightarrow$ \textbf{Non-paternité / Non-maternité} (ou mutation rare, non au programme).
    \end{enumerate}
    \textbf{Réflexe :} Toujours annoter les génotypes déduits \emph{sous} les symboles de l'arbre (ex: $\square_{I^A i}$). Ne pas deviner : écrire « $I^A-$ » si l'allèle secondaire est inconnu (peut être $I^A$ ou $i$).
\end{methode}

\begin{exoresolu}[Arbre généalogique : Transmission du groupe A et du Rhésus]
On étudie une famille camerounaise sur trois générations pour les systèmes ABO et Rhésus.
\begin{itemize}
    \item \textbf{Génération I (Grands-parents) :} 
        \begin{itemize}
            \item I-1 (Homme) : Groupe \textbf{A+} 
            \item I-2 (Femme) : Groupe \textbf{O-}
        \end{itemize}
    \item \textbf{Génération II (Parents) :} 3 enfants :
        \begin{itemize}
            \item II-1 (Homme) : Groupe \textbf{A+} $\rightarrow$ Épouse II-4 (Femme) : Groupe \textbf{AB-}
            \item II-2 (Femme) : Groupe \textbf{O+}
            \item II-3 (Homme) : Groupe \textbf{B+}
        \end{itemize}
    \item \textbf{Génération III (Enfants de II-1 et II-4) :} 2 enfants :
        \begin{itemize}
            \item III-1 (Fille) : Groupe \textbf{A-}
            \item III-2 (Garçon) : Groupe \textbf{B+}
        \end{itemize}
\end{itemize}
\textbf{Questions :}
\begin{enumerate}
    \item Déterminer les génotypes \textbf{complets (ABO + Rhésus)} de I-1, I-2, II-1, II-4, III-1, III-2.
    \item Expliquer comment le phénotype B+ de II-3 est possible avec ces grands-parents.
    \item Calculer la probabilité qu'un futur enfant de II-1 et II-4 soit de groupe \textbf{O+}.
\end{enumerate}
\tcblower
\textbf{Solution détaillée :}
\begin{enumerate}
    \item \textbf{Déduction des génotypes (Notation : Génotype ABO // Génotype Rhésus) :}
    \begin{itemize}
        \item \textbf{I-2 (O-)} : Phénotype O = $ii$ ; Phénotype Rh- = $dd$. $\rightarrow$ \textbf{$ii // dd$} (Certitude absolue).
        \item \textbf{I-1 (A+)} : Doit transmettre $i$ et $d$ à l'enfant II-2 (O+).
            \begin{itemize}
                \item Pour ABO : Enfant O ($ii$) $\rightarrow$ I-1 donne $i$ $\rightarrow$ I-1 = \textbf{$I^A i$}.
                \item Pour Rhésus : Enfant O+ ($D-$) mais I-2 = $dd$ $\rightarrow$ Enfant a reçu $d$ de I-2. Pour être Rh+, il a reçu $D$ de I-1 $\rightarrow$ I-1 = \textbf{$Dd$}.
            \end{itemize}
            $\rightarrow$ \textbf{I-1 : $I^A i // Dd$}.
        \item \textbf{II-1 (A+, fils de I-1 x I-2)} : 
            \begin{itemize}
                \item ABO : Reçu $i$ de I-2. Phénotype A $\rightarrow$ Reçu $I^A$ de I-1. $\rightarrow$ \textbf{$I^A i$}.
                \item Rhésus : Reçu $d$ de I-2. Phénotype + $\rightarrow$ Reçu $D$ de I-1. $\rightarrow$ \textbf{$Dd$}.
            \end{itemize}
            $\rightarrow$ \textbf{II-1 : $I^A i // Dd$}.
        \item \textbf{II-4 (AB-, épouse de II-1)} : 
            \begin{itemize}
                \item ABO : AB $\rightarrow$ \textbf{$I^A I^B$}.
                \item Rhésus : - $\rightarrow$ \textbf{$dd$}.
            \end{itemize}
            $\rightarrow$ \textbf{II-4 : $I^A I^B // dd$}.
        \item \textbf{III-1 (A-, fille de II-1 x II-4)} :
            \begin{itemize}
                \item ABO : Phénotype A. Parents : $I^A i \times I^A I^B$. Gamètes M : $I^A, i$ ; Gamètes F : $I^A, I^B$. Enfant A possible : $I^A I^A$ ou $I^A i$. \textbf{Indéterminé} ($I^A-$).
                \item Rhésus : Phénotype -. Parents : $Dd \times dd$. Enfant - $\rightarrow$ $dd$. $\rightarrow$ \textbf{$dd$}.
            \end{itemize}
            $\rightarrow$ \textbf{III-1 : $I^A- // dd$}.
        \item \textbf{III-2 (B+, garçon de II-1 x II-4)} :
            \begin{itemize}
                \item ABO : Phénotype B. Parents : $I^A i \times I^A I^B$. Pour être B, il faut $I^B$ (venant de II-4) et $i$ (venant de II-1). $\rightarrow$ \textbf{$I^B i$} (Certitude).
                \item Rhésus : Phénotype +. Parents : $Dd \times dd$. Enfant + $\rightarrow$ $Dd$. $\rightarrow$ \textbf{$Dd$}.
            \end{itemize}
            $\rightarrow$ \textbf{III-2 : $I^B i // Dd$}.
    \end{itemize}
    \item \textbf{Origine du B+ de II-3 (Frère de II-1) :} 
    Parents : I-1 ($I^A i // Dd$) $\times$ I-2 ($ii // dd$).
    \begin{itemize}
        \item ABO : Gamètes I-1 : $I^A, i$ ; I-2 : $i$. Enfants possibles : $I^A i$ (A) ou $ii$ (O). \textbf{Impossible d'avoir B ($I^B$)} car aucun parent ne porte $I^B$.
    \end{itemize}
    \textbf{Conclusion :} Il y a une \textbf{incohérence génétique majeure} (non-paternité de I-1 pour II-3, ou erreur de groupage, ou adoption). Le groupe B nécessite l'allèle $I^B$ absent chez les deux grands-parents.
    \item \textbf{Probabilité O+ pour futur enfant de II-1 ($I^A i // Dd$) $\times$ II-4 ($I^A I^B // dd$) :}
    \begin{itemize}
        \item ABO : Croisement $I^A i \times I^A I^B$. Gamètes M : $I^A, i$ ; F : $I^A, I^B$.
            \begin{center}
            \begin{tabular}{|c|c|c|}\hline 
            \textbf{Gamètes Mère $\downarrow$ / Père $\rightarrow$} & \textbf{$I^A$} & \textbf{$I^B$} \\ \hline 
            \textbf{$I^A$} & $I^A I^A$ (A) & $I^A I^B$ (AB) \\ \hline 
            \textbf{$i$} & $I^A i$ (A) & $I^B i$ (B) \\ \hline \end{tabular}
            \end{center}
            $\rightarrow$ Phénotypes possibles : A (50\%), AB (25\%), B (25\%). \textbf{Groupe O (ii) impossible (0\%)}. 
        \item Rhésus : Croisement $Dd \times dd$. $\rightarrow$ 50\% $Dd$ (+), 50\% $dd$ (-).
    \end{itemize}
    \textbf{Probabilité O+ = Probabilité(O) $\times$ Probabilité(+) = 0\% $\times$ 50\% = \textbf{0\%}.}
    \textbf{Conclusion :} Ce couple ne peut \textbf{jamais} avoir d'enfant de groupe O (l'allèle $i$ ne vient que du père, la mère n'en a pas).
\end{enumerate}
\end{exoresolu}

\begin{methode}[Rédiger une conclusion scientifique structurée (CER : Claim - Evidence - Reasoning)]
    \textbf{Objectif :} Répondre à une question problème (ex: « Ce père est-il le père biologique ? », « Cette transfusion est-elle compatible ? ») par un paragraphe rigoureux, noté au BEPC.
    \begin{enumerate}
        \item \textbf{Revendication (Claim / Réponse directe) :} Une phrase affirmative ou négative répondant \textbf{exactement} à la question. \emph{Ex: « L'homme X ne peut pas être le père biologique de l'enfant Y. »}
        \item \textbf{Preuves (Evidence / Données) :} Citer \textbf{uniquement les faits observés ou donnés} dans l'énoncé (résultats de groupage, caryotype, arbre généalogique). \emph{Ex: « L'enfant est de groupe O ($ii$). L'homme X est de groupe AB ($I^A I^B$). »} Pas d'interprétation ici.
        \item \textbf{Raisonnement (Reasoning / Mécanisme) :} Expliquer \textbf{le principe biologique} qui relie les preuves à la revendication. Citer la loi de la ségrégation, la codominance, la transmission des allèles. \emph{Ex: « L'allèle $i$ de l'enfant vient obligatoirement de son père biologique. Or l'homme X, de génotype $I^A I^B$, ne possède pas l'allèle $i$ et ne peut pas le transmettre. »}
        \item \textbf{Conclusion finale :} Reformuler la réponse en intégrant la certitude. \emph{Ex: « Par conséquent, la paternité de l'homme X est exclue avec certitude. »}
    \end{enumerate}
    \textbf{Réflexe :} Utiliser les connecteurs logiques : \textbf{Or, Par conséquent, En effet, Donc, Car}. Éviter : « Je pense que », « Il me semble ». Écrire à l'impersonnel ou « Nous concluons que ».
\end{methode}

\begin{exoresolu}[Rédaction d'une conclusion : Test de paternité par groupes sanguins]
\textbf{Situation :} Un litige oppose un homme (M. K.) à une femme (Mme M.) concernant la paternité de l'enfant E. Les groupages sanguins donnent :
\begin{itemize}
    \item Mère (Mme M.) : Groupe \textbf{A+} (Génotype déterminé par antécédents : $I^A i // Dd$)
    \item Enfant (E) : Groupe \textbf{O-} ($ii // dd$)
    \item Père allégué (M. K.) : Groupe \textbf{AB+} (Génotype : $I^A I^B // Dd$)
\end{itemize}
\textbf{Consigne :} Rédiger la conclusion officielle du biologiste pour le tribunal, en suivant la méthode CER.
\tcblower
\textbf{Solution détaillée (Modèle de rédaction attendue) :}

\textbf{Conclusion du test de paternité}

\textbf{Revendication :} L'homme M. K. \textbf{ne peut pas être le père biologique} de l'enfant E.

\textbf{Preuves (Données expérimentales) :}
\begin{itemize}
    \item L'enfant E présente le phénotype \textbf{O-}, correspondant au génotype \textbf{$ii // dd$} (homozygote récessif pour les deux systèmes).
    \item La mère Mme M. présente le phénotype \textbf{A+} avec le génotype \textbf{$I^A i // Dd$}.
    \item L'homme M. K. présente le phénotype \textbf{AB+} avec le génotype \textbf{$I^A I^B // Dd$}.
\end{itemize}

\textbf{Raisonnement (Interprétation génétique) :}
\begin{itemize}
    \item \textbf{Système ABO :} L'enfant E étant de génotype $ii$, il a obligatoirement reçu un allèle $i$ de sa mère et un allèle $i$ de son père biologique. La mère, de génotype $I^A i$, peut transmettre l'allèle $i$ (probabilité 1/2). L'homme M. K., de génotype $I^A I^B$, ne possède \textbf{pas l'allèle $i$} ; ses gamètes portent exclusivement les allèles $I^A$ ou $I^B$. Il est \textbf{génétiquement impossible} qu'il transmette un allèle $i$.
    \item \textbf{Système Rhésus :} L'enfant est $dd$ (Rh-). Il a reçu un allèle $d$ de chaque parent. La mère ($Dd$) peut le transmettre. L'homme M. K. ($Dd$) peut aussi le transmettre. Ce système \textbf{ne permet pas d'exclure} la paternité.
\end{itemize}

\textbf{Conclusion finale :} L'incompatibilité formelle sur le \textbf{système ABO} (absence de l'allèle $i$ chez l'homme M. K. alors qu'il est obligatoire chez le père biologique d'un enfant $ii$) permet d'\textbf{exclure avec certitude absolue} la paternité de M. K. à l'égard de l'enfant E. Le système Rhésus, bien que compatible, ne contredit pas cette exclusion.
\end{exoresolu}

\begin{methode}[Calculer des probabilités de transmission combinées (ABO + Rhésus) — Indépendance des systèmes]
    \textbf{Objectif :} Déterminer la probabilité d'un phénotype combiné (ex: A+) chez la descendance, sachant que les gènes ABO (chr 9) et Rhésus (chr 1) sont sur \textbf{chromosomes différents} (loi de l'indépendance / 2ème loi de Mendel).
    \begin{enumerate}
        \item \textbf{Séparer les deux systèmes :} Traiter le croisement ABO \textbf{séparément} du croisement Rhésus.
        \item \textbf{Calculer les probabilités pour chaque système :} $P(\text{Phénotype ABO})$ et $P(\text{Phénotype Rhésus})$.
        \item \textbf{Appliquer la règle du produit (indépendance) :} 
        \[ P(\text{Phénotype Combiné}) = P(\text{Phénotype ABO}) \times P(\text{Phénotype Rhésus}) \]
        \item \textbf{Exprimer en pourcentage ou fraction :} Ex: $1/4 \times 1/2 = 1/8 = 12,5\%$.
        \item \textbf{Vérifier :} La somme des probabilités de tous les phénotypes combinés possibles doit faire 1 (ou 100\%).
    \end{enumerate}
    \textbf{Réflexe :} Ne jamais faire un seul carré de Punnett géant 16x16 (gamètes à 4 allèles) sauf si liaison génique (pas le cas ici). La méthode séparée est plus rapide, moins source d'erreur et attendue au BEPC.
\end{methode}

\begin{exoresolu}[Probabilité combinée ABO + Rhésus : Croisement double hétérozygote]
Un couple a les génotypes suivants :
\begin{itemize}
    \item Père : $I^A i // Dd$ (Phénotype A+)
    \item Mère : $I^B i // Dd$ (Phénotype B+)
\end{itemize}
\textbf{Questions :}
\begin{enumerate}
    \item Calculer la probabilité d'avoir un enfant de groupe \textbf{O-}.
    \item Calculer la probabilité d'avoir un enfant de groupe \textbf{AB+}.
    \item Lister tous les phénotypes combinés possibles avec leurs probabilités (en \%). Vérifier la somme.
\end{enumerate}
\tcblower
\textbf{Solution détaillée :}
\begin{enumerate}
    \item \textbf{Probabilité O- :}
    \begin{itemize}
        \item \textbf{ABO :} Croisement $I^A i \times I^B i$. Carré de Punnett $\rightarrow$ $P(O) = P(ii) = \frac{1}{4}$.
        \item \textbf{Rhésus :} Croisement $Dd \times Dd$. Carré de Punnett $\rightarrow$ $P(-) = P(dd) = \frac{1}{4}$.
        \item \textbf{Combiné :} $P(O-) = \frac{1}{4} \times \frac{1}{4} = \frac{1}{16} = \textbf{6,25\%}$.
    \end{itemize}
    \item \textbf{Probabilité AB+ :}
    \begin{itemize}
        \item \textbf{ABO :} $P(AB) = P(I^A I^B) = \frac{1}{4}$.
        \item \textbf{Rhésus :} $P(+) = P(DD \text{ ou } Dd) = \frac{3}{4}$.
        \item \textbf{Combiné :} $P(AB+) = \frac{1}{4} \times \frac{3}{4} = \frac{3}{16} = \textbf{18,75\%}$.
    \end{itemize}
    \item \textbf{Tableau récapitulatif (Produit des probabilités marginales) :}
    \begin{center}
    \begin{tabular}{|c|c|c|c|c|}
        \hline
        \textbf{Phénotype ABO} & \textbf{P(ABO)} & \textbf{Phénotype Rhésus} & \textbf{P(Rh)} & \textbf{P(Combiné)} \\ \hline
        A ($I^A i$) & 1/4 & + (DD, Dd) & 3/4 & 3/16 (18,75\%) \\ \hline
        A & 1/4 & - (dd) & 1/4 & 1/16 (6,25\%) \\ \hline
        B ($I^B i$) & 1/4 & + & 3/4 & 3/16 (18,75\%) \\ \hline
        B & 1/4 & - & 1/4 & 1/16 (6,25\%) \\ \hline
        AB ($I^A I^B$) & 1/4 & + & 3/4 & 3/16 (18,75\%) \\ \hline
        AB & 1/4 & - & 1/4 & 1/16 (6,25\%) \\ \hline
        O ($ii$) & 1/4 & + & 3/4 & 3/16 (18,75\%) \\ \hline
        O & 1/4 & - & 1/4 & 1/16 (6,25\%) \\ \hline
    \end{tabular}
    \end{center}
    \textbf{Vérification :} 
    $4 \times \frac{3}{16} + 4 \times \frac{1}{16} = \frac{12}{16} + \frac{4}{16} = \frac{16}{16} = 1$ (soit 100\%). \textbf{Correct.}
\end{enumerate}
\end{exoresolu}

\begin{attention}[Les 5 erreurs qui coûtent des points au BEPC]
    \begin{enumerate}
        \item \textbf{Confondre Phénotype et Génotype :} Écrire « Le père est $I^A i$ » alors que l'énoncé dit « Le père est de groupe A ». \textbf{Correction :} « Le père est de groupe A, son génotype est $I^A I^A$ \textbf{ou} $I^A i$ ». Ne jamais affirmer un génotype hétérozygote sans preuve (enfant récessif ou parent récessif connu).
        \item \textbf{Oublier l'allèle récessif $i$ (Groupe O) dans les gamètes :} Pour un parent $I^A i$, écrire seulement le gamète $I^A$. \textbf{Conséquence :} Carré de Punnett faux, enfant O impossible alors qu'il est possible. \textbf{Réflexe :} Lister \textbf{tous} les gamètes possibles (séparation des allèles).
        \item \textbf{Inverser la règle de compatibilité transfusionnelle :} Dire « Le donneur O- reçoit de tout le monde » ou « Le receveur AB+ donne à tout le monde ». \textbf{Règle :} \textbf{Donneur} = pas d'Ag étranger pour le receveur. \textbf{Receveur} = pas d'Ac contre les Ag du donneur. \emph{Mnémo :} « Le donneur universel (O-) \textbf{donne} ; Le receveur universel (AB+) \textbf{reçoit} ».
        \item \textbf{Ne pas aligner les centromères sur le caryotype (TP3) :} Coller les chromosomes n'importe comment. \textbf{Conséquence :} Impossible de comparer la taille des bras p et q, impossible de voir une délétion ou inversion. \textbf{Réflexe :} Ligne de centromères horizontale, bras courts (p) vers le haut.
        \item \textbf{Rédiger une conclusion sans justification (Assertion gratuite) :} Écrire « L'enfant est bien le sien car les groupes sanguins correspondent ». \textbf{Exigence BEPC :} \textbf{Preuve} (génotypes) + \textbf{Mécanisme} (transmission allèle) + \textbf{Conclusion}. Sans « Car / Or / Par conséquent », la réponse est incomplète (perte de 50\% des points de l'item « Rédiger »).
    \end{enumerate}
\end{attention}

\newpage

\coursec{Exercices}

\courssub{Je vérifie mes connaissances}

\begin{exercice}[QCM : les bons réflexes]
Pour chaque question, une seule réponse est exacte. Recopie la lettre de la bonne réponse.
\begin{enumerate}
\item Le nombre de chromosomes d'une cellule humaine ordinaire (cellule somatique) est :
\begin{itemize}
\item[a)] 23 ; \quad b) 44 ; \quad c) 46 ; \quad d) 92.
\end{itemize}
\item Un gène est :
\begin{itemize}
\item[a)] un chromosome entier ;
\item[b)] un fragment de chromosome qui porte une information héréditaire (un caractère) ;
\item[c)] une protéine du sang ;
\item[d)] une cellule reproductrice.
\end{itemize}
\item Les allèles du système ABO sont :
\begin{itemize}
\item[a)] $A$, $B$ et $O$ ; \quad b) $A$ et $B$ seulement ; \quad c) $Rh+$ et $Rh-$ ; \quad d) $X$ et $Y$.
\end{itemize}
\item Une personne de génotype $A//A$ est dite :
\begin{itemize}
\item[a)] hétérozygote ; \quad b) homozygote ; \quad c) receveur universel ; \quad d) donneur universel.
\end{itemize}
\end{enumerate}
\end{exercice}

\begin{exercice}[Vrai ou faux ?]
Recopie chaque affirmation en écrivant « Vrai » ou « Faux », puis justifie ta réponse en une ou deux phrases.
\begin{enumerate}
\item Les chromosomes sexuels d'une femme sont $X$ et $Y$.
\item L'allèle $O$ est dominant par rapport aux allèles $A$ et $B$.
\item Un individu de groupe sanguin O possède nécessairement le génotype $O//O$.
\item Une personne de phénotype $[A]$ peut avoir deux génotypes différents.
\item Le caryotype est la représentation ordonnée des chromosomes d'une cellule.
\end{enumerate}
\end{exercice}

\begin{exercice}[Définitions à compléter]
Recopie et complète les phrases suivantes avec les mots de la liste : \emph{gène, allèles, génotype, phénotype, chromosome, caryotype} (chaque mot peut n'être utilisé qu'une seule fois).
\begin{enumerate}
\item Le \ldots\ldots\ldots{} est un bâtonnet formé d'ADN, visible au microscope dans le noyau d'une cellule en division.
\item Un \ldots\ldots\ldots{} est un fragment de chromosome qui porte une information héréditaire.
\item Les différentes versions d'un même gène s'appellent des \ldots\ldots\ldots{}.
\item L'ensemble des gènes d'un individu constitue son \ldots\ldots\ldots{} ; l'ensemble des caractères visibles qui en résultent constitue son \ldots\ldots\ldots{}.
\item La représentation ordonnée et classée par paires des chromosomes d'une cellule s'appelle le \ldots\ldots\ldots{}.
\end{enumerate}
\end{exercice}

\begin{exercice}[Génotypes et phénotypes]
On rappelle que dans le système ABO, les allèles $A$ et $B$ sont dominants et l'allèle $O$ est récessif.
\begin{enumerate}
\item Donne tous les génotypes possibles d'un individu de groupe A.
\item Donne le génotype d'un individu de groupe O.
\item Donne le génotype d'un individu de groupe AB.
\item Une élève de 3ème affirme : « Un individu de groupe B est forcément homozygote. » A-t-elle raison ? Justifie.
\end{enumerate}
\end{exercice}

\courssub{Je m'entraîne}

\begin{exercice}[Échiquier de croisement]
Un homme de groupe sanguin AB épouse une femme de groupe O.
\begin{enumerate}
\item Écris les génotypes du père et de la mère.
\item Indique les types de gamètes (spermatozoïdes et ovules) produits par chaque parent.
\item Construis l'échiquier de croisement.
\item Déduis-en les groupes sanguins possibles des enfants et leur probabilité (en pourcentage).
\item Ce couple peut-il avoir un enfant de groupe O ? Justifie.
\end{enumerate}
\end{exercice}

\begin{exercice}[Interprétation d'un test de groupage]
Au laboratoire d'un centre de santé de Bafoussam, on dépose une goutte de sang de quatre élèves sur deux sérums de test : le sérum anti-A et le sérum anti-B. Le signe « + » signifie qu'il y a agglutination (formation de grumeaux), le signe « $-$ » qu'il n'y en a pas.
\begin{center}
\begin{tabular}{|l|c|c|}
\hline
\rowcolor{popblueL} \textbf{Élève} & \textbf{Sérum anti-A} & \textbf{Sérum anti-B} \\
\hline
Manga & + & $-$ \\
\hline
Aïcha & $-$ & + \\
\hline
Tchoupo & + & + \\
\hline
Ngo Bassa & $-$ & $-$ \\
\hline
\end{tabular}
\end{center}
\begin{enumerate}
\item Rappelle ce que contient le sérum anti-A et ce que signifie une agglutination.
\item Détermine le groupe sanguin de chacun des quatre élèves. Justifie chaque réponse.
\item Lequel de ces élèves est donneur universel ? Justifie.
\end{enumerate}
\end{exercice}

\begin{exercice}[Le facteur Rhésus]
On complète le test précédent par un sérum anti-Rhésus (anti-D). Résultats : Manga : agglutination ; Aïcha : pas d'agglutination ; Tchoupo : agglutination ; Ngo Bassa : pas d'agglutination.
\begin{enumerate}
\item Que signifie l'agglutination avec le sérum anti-Rhésus ?
\item Donne le groupe sanguin complet (ABO et Rhésus) de chaque élève, en utilisant les résultats de l'exercice précédent.
\item Parmi ces quatre élèves, qui peut recevoir du sang de groupe A$+$ ? Justifie.
\end{enumerate}
\end{exercice}

\begin{exercice}[Qui sont les donneurs compatibles ?]
Madame Etoundi, de groupe sanguin B$-$, doit recevoir une transfusion sanguine à l'hôpital de district.
\begin{enumerate}
\item Rappelle la règle de compatibilité pour le système ABO (quels groupes peut recevoir un individu de groupe B ?).
\item Rappelle la règle de compatibilité pour le facteur Rhésus.
\item Les volontaires au don sont : un donneur O$+$, un donneur O$-$, un donneur B$+$, un donneur B$-$, un donneur AB$-$. Établis la liste des donneurs compatibles avec Madame Etoundi et justifie chaque réponse.
\end{enumerate}
\end{exercice}

\courssub{Je me prépare au BEPC}

\begin{exercice}[Type BEPC : le caryotype humain]
Au cours d'un TP, un élève de 3ème photographie au microscope les chromosomes d'une cellule humaine en division, puis les découpe et les range par paires, du plus grand au plus petit.
\begin{enumerate}
\item Comment appelle-t-on la représentation ordonnée ainsi obtenue ?
\item L'élève compte 46 chromosomes. Combien de paires cela représente-t-il ? Combien de paires de chromosomes ordinaires (autosomes) et combien de paires de chromosomes sexuels (gonosomes) y a-t-il ?
\item Sur sa planche, la dernière paire est formée de deux chromosomes identiques en forme de X. Quel est le sexe de l'individu ? Justifie.
\item Qu'aurait observé l'élève si l'individu avait été de sexe masculin ?
\item Sur une autre planche, un élève observe 47 chromosomes, avec trois chromosomes au lieu de deux à la paire n°21. Que peut-on conclure ?
\item Décris en étapes numérotées la démarche à suivre pour confectionner une maquette de caryotype humain à partir de la photographie des chromosomes (TP N°3).
\end{enumerate}
\end{exercice}

\begin{exercice}[Type BEPC : problème de filiation]
À la maternité d'un hôpital de Douala, un doute s'élève sur la filiation d'un nouveau-né. Madame Mbappe est de groupe sanguin O ; son enfant est de groupe B. Deux hommes prétendent être le père : monsieur Ndi, de groupe O, et monsieur Fotso, de groupe B (de génotype $B//O$).
\begin{enumerate}
\item Donne le génotype de la mère. Justifie.
\item L'enfant est de groupe B. Quel(s) génotype(s) peut-il posséder \emph{a priori} ? Sachant que sa mère est de groupe O, quel est obligatoirement son génotype ? Explique.
\item D'où vient l'allèle $B$ de l'enfant ?
\item Monsieur Ndi peut-il être le père de cet enfant ? Justifie ta réponse à l'aide d'un échiquier de croisement.
\item Monsieur Fotso peut-il être le père ? Justifie à l'aide d'un échiquier de croisement et donne la probabilité pour ce couple d'avoir un enfant de groupe B.
\item Quels sont les génotypes possibles du père biologique de cet enfant ?
\end{enumerate}
\end{exercice}

\begin{exercice}[Type BEPC : les expériences historiques de Landsteiner]
En 1900, le médecin autrichien Karl Landsteiner mélange le sérum (liquide du sang contenant les anticorps) de plusieurs personnes avec les globules rouges d'autres personnes, et observe parfois la formation de grumeaux (agglutination). Le tableau ci-dessous résume ses observations : « + » = agglutination, « $-$ » = pas d'agglutination.
\begin{center}
\begin{tabular}{|l|c|c|c|c|}
\hline
\rowcolor{popblueL} \textbf{Sérum de} & \textbf{Globules de Paul} & \textbf{Globules de Marie} & \textbf{Globules de Jean} & \textbf{Globules d'Anne} \\
\hline
Paul & $-$ & + & + & $-$ \\
\hline
Marie & $-$ & $-$ & + & $-$ \\
\hline
Jean & $-$ & + & $-$ & $-$ \\
\hline
Anne & $-$ & + & + & $-$ \\
\hline
\end{tabular}
\end{center}
\begin{enumerate}
\item Que se passe-t-il lorsqu'il y a agglutination ? Pourquoi ce phénomène est-il dangereux lors d'une transfusion sanguine ?
\item En comparant les lignes et les colonnes du tableau, regroupe les personnes ayant le même comportement. Combien de groupes sanguins différents Landsteiner a-t-il ainsi mis en évidence dans cet échantillon ?
\item Sachant que Paul et Anne sont de groupe O, déduis le groupe de Marie et celui de Jean, en justifiant à partir des agglutinations observées.
\item Quel groupe sanguin, absent de ce tableau, complète le système ABO ? Quels résultats d'agglutination observerait-on avec les globules et le sérum d'une personne de ce groupe ?
\item Explique pourquoi un individu de groupe O est appelé « donneur universel » et un individu de groupe AB « receveur universel ».
\item Rédige une conclusion de quatre à six lignes expliquant comment les travaux de Landsteiner ont rendu les transfusions sanguines plus sûres.
\end{enumerate}
\end{exercice}

\courssub{Je vérifie mes réponses (corrigés)}



\newpage

\coursec{Corrigés des exercices}

\begin{corrige}[Exercice 1 — QCM : les bons réflexes]
\begin{enumerate}
\item \textbf{c)} 46. Toute cellule somatique humaine contient 46 chromosomes, soit 23 paires (22 paires d'autosomes + 1 paire de gonosomes).
\item \textbf{b)} Un gène est un fragment de chromosome (d'ADN) qui porte une information héréditaire, c'est-à-dire un caractère.
\item \textbf{a)} Le système ABO comporte trois allèles : $A$, $B$ et $O$.
\item \textbf{b)} Le génotype $A//A$ comporte deux allèles identiques : l'individu est homozygote.
\end{enumerate}
\textbf{Point de vigilance :} ne pas confondre 23 (nombre de paires, ou nombre de chromosomes d'un gamète) avec 46 (nombre de chromosomes d'une cellule somatique).
\end{corrige}

\begin{corrige}[Exercice 2 — Vrai ou faux ?]
\begin{enumerate}
\item \textbf{Faux.} Les chromosomes sexuels d'une femme sont $X$ et $X$ (paire $XX$) ; la paire $XY$ caractérise l'homme.
\item \textbf{Faux.} L'allèle $O$ est \emph{récessif} : il ne s'exprime que lorsqu'il est présent en deux exemplaires ($O//O$). Ce sont les allèles $A$ et $B$ qui sont dominants par rapport à $O$.
\item \textbf{Vrai.} $O$ étant récessif, le phénotype $[O]$ n'est possible qu'avec le génotype $O//O$.
\item \textbf{Vrai.} Le phénotype $[A]$ correspond à deux génotypes possibles : $A//A$ (homozygote) ou $A//O$ (hétérozygote), car $A$ domine $O$.
\item \textbf{Vrai.} C'est la définition du caryotype : les chromosomes d'une cellule sont photographiés, découpés et rangés par paires, du plus grand au plus petit.
\end{enumerate}
\end{corrige}

\begin{corrige}[Exercice 3 — Définitions à compléter]
\begin{enumerate}
\item Le \textbf{chromosome} est un bâtonnet formé d'ADN, visible au microscope dans le noyau d'une cellule en division.
\item Un \textbf{gène} est un fragment de chromosome qui porte une information héréditaire.
\item Les différentes versions d'un même gène s'appellent des \textbf{allèles}.
\item L'ensemble des gènes d'un individu constitue son \textbf{génotype} ; l'ensemble des caractères visibles qui en résultent constitue son \textbf{phénotype}.
\item La représentation ordonnée et classée par paires des chromosomes d'une cellule s'appelle le \textbf{caryotype}.
\end{enumerate}
\end{corrige}

\begin{corrige}[Exercice 4 — Génotypes et phénotypes]
\begin{enumerate}
\item Un individu de groupe A peut avoir deux génotypes : $A//A$ (homozygote) ou $A//O$ (hétérozygote), car l'allèle $A$ est dominant sur $O$.
\item Un individu de groupe O a un seul génotype possible : $O//O$ (l'allèle $O$ est récessif, il faut deux exemplaires pour qu'il s'exprime).
\item Un individu de groupe AB a un seul génotype possible : $A//B$ ($A$ et $B$ sont tous les deux dominants et s'expriment ensemble : codominance).
\item Elle a \textbf{tort} : un individu de groupe B peut être homozygote $B//B$ \emph{ou} hétérozygote $B//O$, car l'allèle $B$ est dominant sur l'allèle $O$. Le phénotype $[B]$ ne permet donc pas de conclure sur l'homozygotie.
\end{enumerate}
\end{corrige}

\begin{corrige}[Exercice 5 — Échiquier de croisement]
\begin{enumerate}
\item Père de groupe AB : génotype $A//B$ ; mère de groupe O : génotype $O//O$.
\item Gamètes du père : $A$ ou $B$ (deux types) ; gamètes de la mère : $O$ uniquement (un seul type).
\item Échiquier de croisement :
\begin{center}
\begin{tabular}{|c|c|c|}
\hline
\rowcolor{popblueL} \textbf{Gamètes} & $A$ (père) & $B$ (père) \\
\hline
$O$ (mère) & $A//O$ & $B//O$ \\
\hline
\end{tabular}
\end{center}
\item Groupes sanguins possibles des enfants : groupe \textbf{A} (génotype $A//O$) avec une probabilité de 50 \%, et groupe \textbf{B} (génotype $B//O$) avec une probabilité de 50 \%.
\item \textbf{Non}, ce couple ne peut pas avoir d'enfant de groupe O. Un enfant de groupe O devrait recevoir un allèle $O$ de chaque parent ; or le père $A//B$ ne possède aucun allèle $O$ à transmettre (il ne transmet que $A$ ou $B$). Le génotype $O//O$ est donc impossible chez les enfants de ce couple.
\end{enumerate}
\textbf{Point de vigilance :} au BEPC, l'échiquier doit présenter en en-tête les gamètes des deux parents, et chaque case doit contenir le génotype complet de l'enfant (deux allèles).
\end{corrige}

\begin{corrige}[Exercice 6 — Interprétation d'un test de groupage]
\begin{enumerate}
\item Le sérum anti-A contient des \cle{anticorps} anti-A. L'agglutination (formation de grumeaux) signifie que les globules rouges testés portent l'\cle{antigène} A : les anticorps du sérum se fixent sur ces antigènes et collent les globules rouges entre eux.
\item Détermination des groupes :
\begin{itemize}
\item \textbf{Manga} : agglutination avec l'anti-A seulement $\Rightarrow$ ses globules portent l'antigène A $\Rightarrow$ groupe \textbf{A}.
\item \textbf{Aïcha} : agglutination avec l'anti-B seulement $\Rightarrow$ antigène B $\Rightarrow$ groupe \textbf{B}.
\item \textbf{Tchoupo} : agglutination avec les deux sérums $\Rightarrow$ antigènes A \emph{et} B $\Rightarrow$ groupe \textbf{AB}.
\item \textbf{Ngo Bassa} : aucune agglutination $\Rightarrow$ aucun antigène A ni B $\Rightarrow$ groupe \textbf{O}.
\end{itemize}
\item \textbf{Ngo Bassa} (groupe O) est le donneur universel : ses globules rouges ne portent aucun antigène A ni B, ils ne peuvent donc être agglutinés par le sérum d'aucun receveur.
\end{enumerate}
\end{corrige}

\begin{corrige}[Exercice 7 — Le facteur Rhésus]
\begin{enumerate}
\item L'agglutination avec le sérum anti-Rhésus (anti-D) signifie que les globules rouges portent l'antigène (facteur) Rhésus : l'individu est Rhésus positif ($Rh+$). En l'absence d'agglutination, il est Rhésus négatif ($Rh-$).
\item Groupes sanguins complets :
\begin{center}
\begin{tabular}{|l|c|c|c|}
\hline
\rowcolor{popblueL} \textbf{Élève} & \textbf{Groupe ABO} & \textbf{Rhésus} & \textbf{Groupe complet} \\
\hline
Manga & A & $+$ & A$+$ \\
\hline
Aïcha & B & $-$ & B$-$ \\
\hline
Tchoupo & AB & $+$ & AB$+$ \\
\hline
Ngo Bassa & O & $-$ & O$-$ \\
\hline
\end{tabular}
\end{center}
\item Peuvent recevoir du sang A$+$ : \textbf{Manga} (A$+$ : même groupe, même Rhésus) et \textbf{Tchoupo} (AB$+$ : receveur universel, il n'a ni anticorps anti-A ni anti-B, et il est $Rh+$). Aïcha (B$-$) et Ngo Bassa (O$-$) ne le peuvent pas : leurs anticorps anti-A agglutineraient les globules A du donneur (incompatibilité ABO), et leur Rhésus négatif interdit en outre le sang $Rh+$.
\end{enumerate}
\end{corrige}

\begin{corrige}[Exercice 8 — Qui sont les donneurs compatibles ?]
\begin{enumerate}
\item Un individu de groupe B possède des anticorps anti-A dans son sérum : il peut recevoir du sang de groupe \textbf{B} ou de groupe \textbf{O} (globules dépourvus d'antigène A).
\item Un individu Rhésus négatif ne doit recevoir que du sang \textbf{Rhésus négatif} : les globules $Rh+$ portent l'antigène Rhésus, qui déclencherait une réaction de rejet chez un receveur $Rh-$.
\item Madame Etoundi (B$-$) peut donc recevoir le sang de :
\begin{itemize}
\item donneur \textbf{O$-$} : compatible (donneur universel, aucun antigène A, B ni Rhésus) ;
\item donneur \textbf{B$-$} : compatible (même groupe ABO, même Rhésus).
\end{itemize}
Donneurs refusés : \textbf{O$+$} et \textbf{B$+$} (Rhésus positif, dangereux pour une receveuse $Rh-$) et \textbf{AB$-$} (l'antigène A de ses globules serait agglutiné par les anticorps anti-A de Madame Etoundi).
\end{enumerate}
\textbf{Point de vigilance :} la compatibilité doit être vérifiée pour \emph{les deux} systèmes (ABO \emph{et} Rhésus) ; une réponse qui n'en cite qu'un est incomplète.
\end{corrige}

\begin{corrige}[Exercice 9 — Type BEPC : le caryotype humain]
\begin{enumerate}
\item Cette représentation ordonnée des chromosomes s'appelle le \cle{caryotype}.
\item 46 chromosomes représentent \textbf{23 paires} : \textbf{22 paires d'autosomes} (chromosomes ordinaires, n°~1 à n°~22) et \textbf{1 paire de gonosomes} (chromosomes sexuels).
\item Deux chromosomes $X$ identiques (paire $XX$) : l'individu est de sexe \textbf{féminin}, car la paire de gonosomes d'une femme est $XX$.
\item Pour un individu de sexe masculin, la dernière paire serait formée de deux chromosomes \emph{différents} : un grand chromosome $X$ et un petit chromosome $Y$ (paire $XY$).
\item 47 chromosomes avec trois chromosomes à la paire n°~21 : il s'agit d'une \cle{trisomie 21}, c'est-à-dire une anomalie du nombre de chromosomes (un chromosome 21 surnuméraire).
\item Démarche de confection de la maquette (TP n°~3) :
\begin{enumerate}
\item Photographier (ou observer au microscope) les chromosomes d'une cellule humaine en division, bien étalés.
\item Découper chaque chromosome sur la photographie.
\item Associer les chromosomes deux par deux selon leur taille et leur forme : ce sont les paires de chromosomes homologues.
\item Ranger les paires par ordre de taille décroissante, de la paire n°~1 à la paire n°~22 (autosomes).
\item Placer en dernier la paire de chromosomes sexuels ($XX$ ou $XY$).
\item Coller les chromosomes sur la planche, numéroter les paires et légender ; compter les chromosomes et identifier le sexe de l'individu.
\end{enumerate}
\end{enumerate}
\textbf{Barème indicatif (sur 10) :} question 1 : 1 pt ; question 2 : 2 pts ; question 3 : 1,5 pt ; question 4 : 1,5 pt ; question 5 : 1 pt ; question 6 : 3 pts (0,5 pt par étape correctement ordonnée).\\
\textbf{Points de vigilance :} citer le terme exact « caryotype » ; justifier le sexe par la composition de la paire de gonosomes ($XX$ ou $XY$) et non par une simple affirmation ; dans la démarche, respecter l'ordre logique des étapes (découpage $\rightarrow$ appariement $\rightarrow$ classement $\rightarrow$ collage).
\end{corrige}

\begin{corrige}[Exercice 10 — Type BEPC : problème de filiation]
\begin{enumerate}
\item La mère est de groupe O ; l'allèle $O$ étant récessif, son génotype est obligatoirement $O//O$.
\item \emph{A priori}, un enfant de groupe B peut être $B//B$ ou $B//O$. Mais sa mère $O//O$ ne transmet qu'un allèle $O$ : l'enfant a donc obligatoirement reçu $O$ de sa mère, et son génotype est nécessairement $B//O$.
\item L'allèle $B$ de l'enfant vient donc forcément du père biologique.
\item Monsieur Ndi, de groupe O, a le génotype $O//O$ : il ne produit que des gamètes $O$.
\begin{center}
\begin{tabular}{|c|c|}
\hline
\rowcolor{popblueL} \textbf{Gamètes} & $O$ (M.~Ndi) \\
\hline
$O$ (mère) & $O//O$ [O] \\
\hline
\end{tabular}
\end{center}
Tous les enfants de ce couple seraient $O//O$, de groupe O. Monsieur Ndi \textbf{ne peut pas} être le père d'un enfant de groupe B.
\item Monsieur Fotso ($B//O$) produit deux types de gamètes : $B$ et $O$.
\begin{center}
\begin{tabular}{|c|c|c|}
\hline
\rowcolor{popblueL} \textbf{Gamètes} & $B$ (M.~Fotso) & $O$ (M.~Fotso) \\
\hline
$O$ (mère) & $B//O$ [B] & $O//O$ [O] \\
\hline
\end{tabular}
\end{center}
Ce couple peut avoir un enfant de groupe B ($B//O$) avec une probabilité de \textbf{50 \%} : monsieur Fotso \textbf{peut} être le père.
\item Le père biologique doit posséder au moins un allèle $B$ à transmettre : ses génotypes possibles sont $B//B$, $B//O$ ou $A//B$ (groupes B ou AB).
\end{enumerate}
\textbf{Barème indicatif (sur 10) :} question 1 : 1 pt ; question 2 : 2 pts ; question 3 : 1 pt ; question 4 : 2 pts (échiquier + conclusion) ; question 5 : 2,5 pts (échiquier + probabilité) ; question 6 : 1,5 pt.\\
\textbf{Points de vigilance :} justifier chaque génotype par la dominance/récessivité des allèles ; l'échiquier est exigé pour conclure, une simple affirmation ne suffit pas ; attention, le test de groupe sanguin permet d'\emph{exclure} un père (M.~Ndi) mais pas de prouver avec certitude la paternité de M.~Fotso.
\end{corrige}

\begin{corrige}[Exercice 11 — Type BEPC : les expériences historiques de Landsteiner]
\begin{enumerate}
\item Lors de l'agglutination, les anticorps du sérum se fixent sur les antigènes étrangers portés par les globules rouges et les collent en grumeaux. Lors d'une transfusion, ces grumeaux peuvent boucher les vaisseaux sanguins et provoquer la destruction des globules rouges, ce qui peut entraîner la mort du receveur.
\item Paul et Anne ont exactement le même comportement (mêmes résultats en ligne et en colonne) : ils appartiennent au même groupe. Marie et Jean présentent chacun un comportement différent de tous les autres. Landsteiner met donc en évidence \textbf{trois groupes} sanguins différents dans cet échantillon.
\item Les globules de Paul et d'Anne (groupe O) ne sont agglutinés par aucun sérum : ils ne portent aucun antigène ; leur sérum agglutine les globules de Marie et de Jean : il contient des anticorps anti-A et anti-B.
\begin{itemize}
\item Les globules de \textbf{Marie} sont agglutinés par les sérums de Paul (anti-A + anti-B) et de Jean, mais pas par son propre sérum : Marie porte l'antigène A (et des anticorps anti-B) ; elle est de groupe \textbf{A}.
\item Les globules de \textbf{Jean} sont agglutinés par les sérums de Paul et de Marie : Jean porte l'antigène B (et des anticorps anti-A) ; il est de groupe \textbf{B}.
\end{itemize}
\item Le groupe manquant est le groupe \textbf{AB}. Ses globules, portant les antigènes A \emph{et} B, seraient agglutinés par tous les sérums du tableau ; son sérum, ne contenant ni anticorps anti-A ni anticorps anti-B, n'agglutinerait les globules d'aucune des quatre personnes.
\item Un individu de groupe O ne porte aucun antigène A ni B sur ses globules : son sang ne peut pas être agglutiné par le sérum d'un receveur, d'où le nom de « \cle{donneur universel} ». Un individu de groupe AB ne possède aucun anticorps anti-A ni anti-B dans son sérum : il peut recevoir les globules de tous les groupes, d'où le nom de « \cle{receveur universel} ».
\item \textbf{Conclusion rédigée :} Avant les travaux de Landsteiner, les transfusions sanguines échouaient souvent et tuaient parfois le receveur, car on ignorait l'existence des groupes sanguins. En découvrant les groupes A, B et O (le groupe AB sera identifié ensuite) et en expliquant le phénomène d'agglutination par la rencontre entre antigènes et anticorps, Landsteiner a montré qu'il faut obligatoirement vérifier la compatibilité entre le sang du donneur et celui du receveur. Le groupage sanguin, réalisé avec les sérums anti-A et anti-B, est ainsi devenu un contrôle systématique avant toute transfusion, ce qui a rendu celles-ci sûres et a sauvé d'innombrables vies.
\end{enumerate}
\textbf{Barème indicatif (sur 10) :} question 1 : 2 pts ; question 2 : 1,5 pt ; question 3 : 2 pts ; question 4 : 1,5 pt ; question 5 : 1,5 pt ; question 6 : 1,5 pt.\\
\textbf{Points de vigilance :} bien distinguer antigène (porté par les globules rouges) et anticorps (présent dans le sérum) ; la déduction des groupes doit s'appuyer explicitement sur les agglutinations observées dans le tableau ; la conclusion rédigée doit être structurée (problème initial $\rightarrow$ découverte $\rightarrow$ conséquence pratique) et respecter la longueur demandée.
\end{corrige}

\newpage

\coursec{Fiche bilan}

\begin{popfigure}
\begin{tikzpicture}[mindmap, grow cyclic, every node/.style={concept}, text width=2.6cm, align=center,
  root concept/.append style={concept color=popblue, font=\bfseries, text width=3.2cm},
  level 1/.append style={level distance=4.6cm, sibling angle=60, font=\small},
  level 2/.append style={level distance=2.9cm, sibling angle=45, font=\scriptsize, text width=2.2cm}]
\node[root concept]{Expression de l'information génétique}
  child[concept color=poporange]{ node {Chromosomes humains}
    child { node {23 paires = 46 chromosomes} }
    child { node {22 paires d'autosomes + 1 paire de gonosomes (XX ou XY)} }
  }
  child[concept color=popgreen]{ node {Caryotype (TP 3)}
    child { node {Classement par taille et forme} }
    child { node {Garçon : XY ; Fille : XX} }
  }
  child[concept color=poppurple]{ node {Les gènes}
    child { node {Portés par les chromosomes} }
    child { node {Un gène = information d'un caractère} }
    child { node {Allèles : versions d'un gène} }
  }
  child[concept color=poppink]{ node {Système ABO}
    child { node {3 allèles : $A$, $B$, $O$} }
    child { node {4 groupes : A, B, AB, O} }
  }
  child[concept color=popgold]{ node {Facteur Rhésus}
    child { node {Rh$+$ : présence de l'antigène D} }
    child { node {Rh$-$ : absence de l'antigène D} }
  }
  child[concept color=popblue!70]{ node {TP 4 : groupage}
    child { node {Sérums tests anti-A, anti-B, anti-D} }
    child { node {Agglutination = présence de l'antigène} }
  };
\end{tikzpicture}
\legende{Carte mentale de la leçon : l'information génétique est portée par les chromosomes ; les gènes (et leurs allèles) expliquent la diversité humaine, illustrée par les groupes sanguins ABO et Rhésus.}
\end{popfigure}

\begin{aretenir}[L'essentiel de la leçon : définitions clés]
\begin{itemize}
  \item \cle{Chromosome} : structure du noyau de la cellule qui porte l'information génétique ; il est visible au microscope lors de la division cellulaire.
  \item \cle{Caryotype} : photographie ordonnée de l'ensemble des chromosomes d'une cellule, classés par paires, par taille et par forme.
  \item \cle{Gène} : fragment de chromosome qui porte l'information déterminant un caractère (ex. : groupe sanguin).
  \item \cle{Allèle} : version possible d'un gène. Exemple : le gène du groupe sanguin existe en trois allèles $A$, $B$ et $O$.
  \item \cle{Antigène} : substance portée par les globules rouges, reconnue par les anticorps (antigènes A, B, D).
  \item \cle{Agglutination} : formation d'un dépôt (grumeaux) quand un anticorps rencontre l'antigène correspondant : c'est le signe d'une réaction positive.
\end{itemize}
\end{aretenir}

\begin{aretenir}[Les chromosomes de l'espèce humaine]
\begin{itemize}
  \item La cellule humaine possède \cle{46 chromosomes}, répartis en \cle{23 paires} : 22 paires d'\cle{autosomes} (communes aux deux sexes) et 1 paire de \cle{gonosomes} (chromosomes sexuels).
  \item Fille : $44 + XX$ ; garçon : $44 + XY$. C'est le père (spermatozoïde porteur de $X$ ou de $Y$) qui détermine le sexe de l'enfant.
  \item Les chromosomes d'une même paire sont dits \cle{homologues} : même taille, même forme, mêmes gènes.
\end{itemize}
\end{aretenir}

\begin{aretenir}[Groupes sanguins : à connaître par c\oe ur]
\begin{center}
\begin{tabularx}{\linewidth}{|l|X|X|X|}
\hline
\rowcolor{popblueL} \textbf{Groupe} & \textbf{Antigènes sur les globules rouges} & \textbf{Agglutination avec anti-A} & \textbf{Agglutination avec anti-B} \\
\hline
A & antigène A & oui & non \\
\hline
B & antigène B & non & oui \\
\hline
AB & antigènes A et B & oui & oui \\
\hline
O & aucun (ni A ni B) & non & non \\
\hline
\end{tabularx}
\end{center}
\begin{itemize}
  \item Facteur Rhésus : agglutination avec le sérum \cle{anti-D} $\Rightarrow$ Rh$+$ ; pas d'agglutination $\Rightarrow$ Rh$-$.
  \item Notation complète d'un groupe : lettre + signe Rhésus, ex. A$+$, O$-$.
  \item Les allèles $A$ et $B$ dominent l'allèle $O$ ; $A$ et $B$ sont codominants entre eux (groupe AB).
\end{itemize}
\end{aretenir}

\begin{methode}[Méthodes express]
\begin{itemize}
  \item \cle{Construire un caryotype} : découper les chromosomes $\rightarrow$ former les paires homologues $\rightarrow$ classer par taille décroissante $\rightarrow$ identifier les gonosomes (XX ou XY) $\rightarrow$ conclure sur le sexe.
  \item \cle{Déterminer un groupe sanguin} : déposer une goutte de sang dans trois cases $\rightarrow$ ajouter anti-A, anti-B, anti-D $\rightarrow$ observer : agglutination = antigène présent $\rightarrow$ conclure (lettre + Rhésus).
  \item \cle{Rédiger une conclusion} : rappeler l'observation, l'interpréter avec les savoirs, puis conclure en répondant au problème posé.
\end{itemize}
\end{methode}

\begin{aretenir}[Checklist avant le BEPC]
\begin{itemize}
  \item[$\square$] Je sais donner la formule chromosomique de l'espèce humaine : 46 chromosomes = 23 paires.
  \item[$\square$] Je distingue autosomes (22 paires) et gonosomes (1 paire : XX ou XY).
  \item[$\square$] Je sais expliquer pourquoi c'est le père qui détermine le sexe de l'enfant.
  \item[$\square$] Je sais définir : chromosome, caryotype, gène, allèle, antigène, agglutination.
  \item[$\square$] Je sais décrire les étapes de la conception d'une maquette de caryotype humain.
  \item[$\square$] Je connais les trois allèles du gène ABO ($A$, $B$, $O$) et les quatre groupes (A, B, AB, O).
  \item[$\square$] Je sais interpréter un test de groupage avec les sérums anti-A, anti-B et anti-D.
  \item[$\square$] Je sais compléter et exploiter le tableau des réactions d'agglutination.
  \item[$\square$] Je sais écrire un groupe sanguin complet avec le facteur Rhésus (ex. : B$+$, AB$-$).
  \item[$\square$] Je sais rédiger une démarche scientifique : observation, hypothèse, expérience, interprétation, conclusion.
\end{itemize}
\end{aretenir}

\findecours

\end{document}
